\makeatletter
\@ifundefined{ifresearchtrackincluded}{%
  \newif\ifresearchtrackincluded
  \researchtrackincludedfalse
}{}
\makeatother

%! Author = Omar Iskandarani
%! Companion to Swirl-String-Theory Canon-v0.8.21 (included from 	exttt{SST\_CANON-v0.8.21.tex} before the manual bibliography).

\section{Research-track sectors migrated from the main canon}

\noindent
This companion document collects those sectors that remain useful for continuity and future development but are too long, too provisional, or too phenomenology-specific for the core body of \textit{Swirl-String-Theory\_Canon-v0.8.21}. Their migration out of the main canon is editorial rather than negative: each subsection remains part of the broader SST research program, but none is required for the minimal logical closure of the core canon.

\medskip
\noindent
\textbf{Editorial note (v0.8.21):} the maximal swirl tension sector ($F_{\rm swirl}^{\max}$, Rydberg bridge identities, Gibbons-class maximum-tension comparison, and the threefold interpretation table) is developed in the core document under \verb|been_processed/v0.8.21/SST_CANON-v0.8.21.tex|, subsection \texttt{Maximal Swirl Tension} (label \texttt{sec:Fmax}); it is not duplicated here.

\medskip








\subsection{Thermodynamic Radiation Interface}
A recurring legacy idea is that a local swirl-energy density may define an effective temperature field and thereby an emergent radiation sector. Let
\begin{align}
    T_{\mathrm{eff}} = T_{\mathrm{eff}}\big(\rhoE,\, \omega,\, \SwirlClock\big)
\end{align}
denote an unspecified constitutive relation. A minimal legacy interface then asks whether Planck-type spectra can emerge from a medium whose local excitation scale is set by $T_{\mathrm{eff}}$.

\textbf{[ORTHODOX]}: blackbody spectra are governed by Planck's law once a valid temperature variable is defined.

\textbf{[SPECULATIVE]}: SST does not yet derive the constitutive relation above, so this sector remains a research program rather than canonically closed physics.

\subsection{Minimal Coupling Interface to QED}
Early research notes repeatedly explored whether an effective vector potential could be built from circulation variables. The most conservative placeholder takes the form
\begin{align}
    \nabla \rightarrow \nabla - i\frac{m}{\hbar}A_{\mathrm{eff}},
    \qquad
    A_{\mathrm{eff}} = \chi \, v,
\end{align}
where $\chi$ is a scalar weight and $v$ is the local carrier velocity field.

\textbf{[ORTHODOX]}: minimal coupling is the standard gateway from free to gauge-coupled wave dynamics.

\textbf{[SPECULATIVE]}: the above substitution is only an analogy template unless a full gauge-covariant derivation is supplied.

\subsection{Ribbon-Invariant Chirality Refinement}
The older canon layers also suggested using the Calugareanu--White--Fuller relation
\begin{align}
    Lk = Tw + Wr
\end{align}
as a refinement of the knot-based matter/antimatter dictionary.

\textbf{[ORTHODOX]}: this is a standard geometric identity for framed curves and ribbons.

\textbf{[DERIVED BOOKKEEPING]}: once a framed curve is specified, the identity separates writhe and twist contributions.

\textbf{[SPECULATIVE]}: the dynamical stability of writhe-dominated or twist-dominated realizations, and the mapping from $\operatorname{sign}(Tw)$ or related chirality measures to particle/antiparticle sectors, remains a model hypothesis.

\subsection{Particle-Candidate Mapping Layer}
\textbf{[SPECULATIVE] Working dictionary:} a minimal knot/composition mapping retained from earlier SST work is summarized in Table~\ref{tab:particle-candidates}. This table is organizational rather than definitive; it records the current topological assignments used by the mass and taxonomy program.

\begin{table}[h]
\centering
\begin{tabularx}{\linewidth}{p{0.28\linewidth}p{0.22\linewidth}X}
\hline
Knot class / composition & Candidate & Status / role \\
\hline
$3_1$ trefoil & electron / positron sector & baseline chiral lepton calibration \\
$T(2,5)$ or equivalent higher torus class & muon & lepton hierarchy candidate \\
$T(2,7)$ or equivalent higher torus class & tau & lepton hierarchy candidate \\
twist/clasp knots, quasi-unknot loops & quark-like defects & research-track fractional boundary winding candidates \\
$5_2+5_2+6_1$ and nearby composites & baryon candidates & legacy mass-functional dictionary, still under audit \\
triple-gear locked analogue & proton-like charged baryon analogue & mechanical model for nonzero exterior collective rotation \\
Borromean-type analogue & neutron-like neutral baryon analogue & mechanical/topological model for zero exterior charge with internal triple linking \\
\hline
\end{tabularx}
\caption{Minimal knot-class and analogue summary (v0.8.21 research-track), aligned with main canon semantic corrections.}
\label{tab:particle-candidates}
\end{table}

\textbf{[CRITICAL NOTE]}: Table~\ref{tab:particle-candidates} is a working dictionary only. A full canonical assignment requires an explicit invariant set and a mass-functional comparison program.


\subsection{Quark-Like Twist Knots as Quasi-Unknot Defects}
\label{subsec:rt_quark_twist_quasi_unknot}

\textbf{[RESEARCH TRACK / CANON CANDIDATE]} Twist knots and related clasp knots may be interpreted as quasi-unknot carrier loops with a localized twist/clasp defect. Geometrically such a knot can be stretched so that it resembles an unknot with a central aperture, while topologically remaining nontrivial until a crossing change or reconnection occurs.

A useful diagnostic decomposition is
\begin{align}
    K_{\mathrm{twist}}
    \simeq
    0_1 + \Delta_{\mathrm{clasp}},
\end{align}
where $\Delta_{\mathrm{clasp}}$ denotes a localized topological defect region. Candidate audit quantities are
\begin{align}
    \mathcal{U}_{\mathrm{quasi}},\qquad
    A_{\mathrm{aperture}},\qquad
    f_{\mathrm{clasp}},\qquad
    f_{\kappa,\mathrm{localized}},
\end{align}
standing for quasi-unknot score, central aperture, clasp-localization fraction, and curvature-concentration fraction.

\textbf{[CHARGE PLACEHOLDER]} Fractional quark-like charge should not be assigned to the trefoil. If retained, it is represented as fractional exterior boundary winding,
\begin{align}
    \frac{Q}{e}=q_{\partial},
    \qquad
    q_u=+\frac{2}{3},
    \qquad
    q_d=-\frac{1}{3}.
\end{align}
The topological knot supplies the defect carrier; $q_{\partial}$ supplies the charge-sector bookkeeping.

\subsection{Twist-Ladder Closure from a Folded Clasp}
\label{subsec:rt_twist_ladder_closure}

\textbf{[ORTHODOX / RESEARCH TRACK].}
The rope experiment described in the project discussion is naturally modeled as
a twist box plus a single clasp closure. If \(m\) denotes the number of signed
half-twists retained in the folded region at closure, the resulting twist-knot
ladder is
\begin{align}
    K_m
    =
    \text{clasp}+\text{\(m\)-half-twist box},
    \qquad
    m\in\mathbb{Z}.
    \label{eq:rt_twist_ladder_Km}
\end{align}
For the positive convention used here,
\begin{align}
    K_1=3_1,\qquad
    K_2=4_1,\qquad
    K_3=5_2,\qquad
    K_4=6_1,\qquad
    K_5=7_2,\qquad
    K_6=8_1 .
    \label{eq:rt_twist_ladder_first_terms}
\end{align}
The standard invariant checks are
\begin{align}
    c_{\min}(K_m)=m+2,
    \qquad
    \det(K_m)=2m+1
    \quad (m\ge 1).
    \label{eq:rt_twist_ladder_invariants}
\end{align}
Changing the twist sign gives the mirror sector,
\begin{align}
    K_{-m}=\overline{K_m}.
    \label{eq:rt_twist_ladder_mirror}
\end{align}

\textbf{[CANON-COMPATIBLE INTERPRETATION].}
End-rotation and plectonemic folding redistribute twist and writhe according
to the framed-ribbon bookkeeping \(Lk=Tw+Wr\), but they do not change the
topological class of an already closed ideal carrier. The knot class is fixed
only when the clasp/boundary closure is performed. This module is therefore a
research-track generator for twist/clasp candidates and for the \(m=1\) trefoil
closure picture; it is not yet a dynamical theorem of Euler--SST.

\subsection{Triple-Gear and Borromean Baryon Analogues}
\label{subsec:rt_triple_gear_borromean_baryons}

\textbf{[MECHANICAL ANALOGUE / NOT A PROOF]} The triple-gear print and Borromean-type links are retained as baryon analogues, not as final particle assignments.

The proton-like analogue is a three-component locked system with nonzero exterior collective rotation,
\begin{align}
    \mathcal{B}_{p}^{\mathrm{gear}}
    =
    (C_1,C_2,C_3;\ q_{\partial}=+1,\ \Theta_{\mathrm{ext}}\neq 0),
\end{align}
where the charge proxy is the boundary winding $q_{\partial}$, not raw pairwise helicity.

The neutron-like analogue is a Borromean or Borromean-adjacent system with no net exterior boundary winding but nontrivial three-component topology,
\begin{align}
    \mathcal{B}_{n}^{\mathrm{Bor}}
    =
    (C_1,C_2,C_3;\ q_{\partial}=0,\ \mu_{123}\neq 0).
\end{align}
This preserves the intuition that the neutron is neutral externally while still carrying internal topological structure.

\textbf{[FALSIFIERS]} The proton analogue fails if the locked triple system has no stable exterior collective phase. The neutron analogue fails if the Borromean sector necessarily induces nonzero exterior boundary winding.

\subsection{Trefoil--Gear Kinematic Non-Equivalence}
\label{subsec:rt_trefoil_gear_kinematic_nonequivalence}

\textbf{[RESEARCH TRACK]} The triple gear and segmented trefoil are not kinematically identical even if they use related printed parts. The gear model measures rotation about a straight axle; the trefoil model measures phase transport along a curved framed tube.

For the gear analogue,
\begin{align}
    \omega_{\mathrm{gear}}
    =
    \frac{1}{5}\,\omega_{\mathrm{core}}
\end{align}
if five full core rotations produce one full gear rotation.

For the segmented trefoil Blender model, the shell/blade winding budget is
\begin{align}
    N_{\mathrm{shell,total}}=10,
    \qquad
    N_{\mathrm{shell,segment}}=\frac{10}{3}.
\end{align}
If the same core-shaft mechanism gives five core rotations per segment, then
\begin{align}
    \frac{N_{\mathrm{shell,segment}}}{N_{\mathrm{core,segment}}}
    =
    \frac{10/3}{5}
    =
    \frac{2}{3}.
\end{align}
This $2/3$ is therefore a conditional shell/core kinematic ratio, not a trefoil-electron charge assignment.

\subsection{Higgs-as-Trefoil High-Temperature Hypothesis}
\label{subsec:rt_higgs_trefoil_parked}

\textbf{[PARKED / SPECULATIVE]} The idea that a Higgs-like state is a thermally activated trefoil phase is not canon-ready. A single trefoil is chiral and lepton-like in the current taxonomy, whereas a Higgs-like excitation must behave effectively as a scalar sector.

A viable future version would need at least one of the following mechanisms:
\begin{enumerate}
    \item a trefoil--anti-trefoil paired condensate with scalar effective symmetry,
    \item a high-temperature unframed phase where chirality averages out,
    \item an independent scalar amplitude mode coupled to trefoil attachment without identifying the Higgs with the trefoil itself.
\end{enumerate}
Until such a mechanism is derived, the hypothesis remains outside the canon.



\subsection{Hydrodynamic Exchange and the Pauli Barrier}
\textbf{[ORTHODOX]}: for two thin filamentary loops $C_1$ and $C_2$ in an incompressible medium, the Biot--Savart interaction energy is
\begin{align}
    E_{\mathrm{int}}
    =
    \frac{\rhoF\Gamma_1\Gamma_2}{4\pi}
    \oint_{C_1}\oint_{C_2}
    \frac{d\ell_1\cdot d\ell_2}{\norm{r_1-r_2}}.
\end{align}

\textbf{[DERIVED]}: if two identical electron-candidate loops are forced into the same spatial state, regularization at a short-distance cutoff $a_{\rm cut}$ yields the overlap penalty
\begin{align}
    V_{\mathrm{Pauli}}(d)
    \approx
    \frac{\rhoF\Gamma_0^2}{4\pi}
    \iint
    \frac{d\ell_1\cdot d\ell_2}{\sqrt{\norm{r_1(\theta_1)-r_2(\theta_2)+d}^2+a_{\rm cut}^2}},
\end{align}
with maximal-overlap scale
\begin{align}
    V_{\mathrm{Pauli}}^{\max}
    \approx
    \frac{\rhoF\Gamma_0^2}{4\pi}\left(\frac{L}{a_{\rm cut}}\right)\mathcal{O}(1).
\end{align}

Using the canonical values
\begin{align}
    \rhoF &= 7.0\times 10^{-7}\ \mathrm{kg\,m^{-3}}, \\
    \rc &= 1.40897017\times 10^{-15}\ \mathrm{m}, \\
    \Gamma_0 &= 2\pi \rc\,\vnorm \approx 9.6836\times 10^{-9}\ \mathrm{m^2\,s^{-1}},
\end{align}
and \(L\sim 2\pi a_0\) at the hydrogenic scale, the legacy benchmark corresponds to the cutoff calibration \(a_{\rm cut}=\rc\). This does \emph{not} identify the resolved tube radius \(a_{\rm core}\) with \(\rc\); it only fixes the ultraviolet cutoff used in this benchmark. One obtains
\begin{align}
    V_{\mathrm{Pauli}}^{\max} \sim 1.23\times 10^{-18}\ \mathrm{J} \sim 7.6\ \mathrm{eV}.
\end{align}

\textbf{[CRITICAL NOTE]}: this is retained as a cutoff-calibrated scale benchmark, not yet as a closed theorem. Its role is to show that the exclusion-energy estimate falls in an atomically relevant window rather than at obviously pathological scales. A physical resolved-core calculation must use an independently specified \(a_{\rm core}\) or profile model.

\subsection{Numerical Benchmark Layer}
\textbf{[DERIVED] Computational-program requirement:} earlier SST work contained explicit benchmark tables for particle and atomic masses. In the present canon, the benchmark layer is retained as a reproducibility requirement rather than as a stand-alone authority claim.

\begin{table}[h]
\centering
\begin{tabular}{lllll}
\hline
Particle & $m_{\mathrm{exp}}$ (MeV) & $m_{\mathrm{SST}}$ (MeV) & rel. error & mode class \\
\hline
e & 0.51099895 & 0.51099895 & 0 & predictive/closure \\
$\mu$ & 105.6583745 & \dots & \dots & predictive candidate \\
$\tau$ & 1776.86 & \dots & \dots & predictive candidate \\
p & 938.272088 & \dots & \dots & predictive or closure \\
n & 939.565420 & \dots & \dots & predictive or closure \\
\hline
\end{tabular}
\caption{Benchmark summary structure to be populated from canonical scripts and archived CSV output.}
\end{table}

\textbf{[CRITICAL NOTE]}: a future v0.8.x benchmark appendix should include the exact script or package path used for generation, archived CSV outputs, mode definitions (predictive vs. exact-closure), and uncertainty propagation on canonical constants and geometric factors.

\subsection{Helicity-by-base archive (diagnostic layer)}
\textbf{[DERIVED] Archive role:} the SST helicity-by-base outputs are retained as a measurable balance diagnostic for taxonomy support, not as an independent law.
Define the midpoint-distance diagnostic
\begin{align}
    \Delta_H(K)=\abs{\widehat{\mathcal{H}}_{b}(K)+\frac{1}{2}},
\end{align}
where $\widehat{\mathcal{H}}_{b}(K)$ is the archived normalized helicity-like base indicator for class $K$.

Near-balanced values with $\Delta_H(K)\ll 1$ indicate candidate dark or quasi-dark sectors only when symmetry and instability filters are also satisfied.

\textbf{[CRITICAL NOTE]}: no single helicity-like scalar is sufficient for sector assignment; classifier decisions remain conjunction rules over symmetry, helicity balance, and low-order stability content.

\textbf{[CANON WORKFLOW NOTE]}: if SST-labelled and parallel VAM-labelled helicity outputs are numerically identical, the SST-labelled archive is the only required source for canon workflows.

\subsection{Trefoil-closure archive (benchmark discipline)}
\textbf{[DERIVED] Benchmark role:} ideal-trefoil robustness sweeps, final-best-estimate tables, and theorem-style reports are retained as benchmark-discipline infrastructure.

\textbf{[CRITICAL NOTE]}: practical best-estimate closure and theorem-level closure are distinct.
A result may remain useful for calibration while still failing strict theorem-level tail or postcheck requirements.

\textbf{[REQUIRED METADATA]}: benchmark reports should include explicit pass/fail logic, postchecks, extrapolation summaries, and local-curve diagnostics so that closure claims are auditable.

\subsection{G7: Ring Constant versus Trefoil-Geometric Closure}
\label{subsec:g7_ring_constant_vs_trefoil_closure}

\textbf{[NUMERICAL OBSERVATION / OPEN GATE / NOT CANONICAL].}
Track-B calculations give the numerical vortex-ring constant
\begin{align}
    \alpha_{\mathrm{ring}}
    =
    1.61935\ldots
    =
    \varphi + 0.00132\ldots ,
\end{align}
with
\begin{align}
    \frac{\alpha_{\mathrm{ring}}-\varphi}{\varphi}
    \approx
    8.16\times 10^{-4}
    =
    0.0816\%.
\end{align}
This is retained as a genuine sub-per-mille numerical proximity between the
GP/NLSE ring-energy constant and the golden ratio.

However, the identification
\begin{align}
    \alpha_{\mathrm{ring}}
    \stackrel{?}{=}
    \mathcal{C}_{\mathrm{trefoil}}
    =
    \frac{R_{\mathrm{eff}}}{a_{\mathrm{eff}}}
\end{align}
is not derived. The quantity \(\alpha_{\mathrm{ring}}\) is a vortex-ring
energy constant, whereas \(\mathcal{C}_{\mathrm{trefoil}}\) is a geometric
centerline-thickness observable of an ideal trefoil.

The recursive closure relation
\begin{align}
    \mathcal{C}_{\mathrm{cl}}
    =
    1+\frac{1}{\mathcal{C}_{\mathrm{cl}}}
\end{align}
implies
\begin{align}
    \mathcal{C}_{\mathrm{cl}}=\varphi,
\end{align}
but this proves only the consequence of the ansatz. It does not prove that
the ideal trefoil realizes the ansatz.

We therefore split G7 into three independent subgates:
\begin{align}
    \mathrm{G7a}:&\quad \alpha_{\mathrm{ring}}\to\varphi,\\
    \mathrm{G7b}:&\quad
    \mathcal{C}_{\mathrm{trefoil}}
    =
    \frac{R_{\mathrm{eff}}}{a_{\mathrm{eff}}}
    \to\varphi,\\
    \mathrm{G7c}:&\quad \xi=r_c.
\end{align}
At present all three remain open. In particular, \(\xi=r_c\) is an imposed
core-matching condition in the NLSE/GP diagnostic layer, not a derived
identity.

\textbf{Canon rule.}
The canon may cite
\begin{align}
    \alpha_{\mathrm{ring}}\approx\varphi
\end{align}
as a numerical observation, but it must not cite this as evidence that
\begin{align}
    \mathcal{C}_{\mathrm{trefoil}}\approx\varphi
\end{align}
without an independent geometric convergence test on ideal-trefoil data.

\subsection{\texorpdfstring{$\phi_{3_1}$}{phi3_1} transform archive (deferred)}
\textbf{[SPECULATIVE] Status:} retained as research-track spectral evidence for later zeta/spectral programs.

\textbf{[SCOPE RULE]}: this archive is not part of the current dark-sector or galactic canonization closure path in v0.8.x.

\subsection{Gauge-Sector Roadmap}
\textbf{[SPECULATIVE] Structural roadmap:} a retained minimal statement of the gauge program is
\begin{align}
    g_{\mathrm{eff}} = su(3) \oplus su(2) \oplus u(1),
\end{align}
with effective gauge phases organized as holonomy data of multi-director transport.

A minimal phase-based representation is
\begin{align}
    A^a_{\mu}(x) = \partial_{\mu}\theta^a(x),
\end{align}
and clock coupling acts as a common phase reference,
\begin{align}
    \theta^a(x) \mapsto \theta^a(x)+q^a\chi(x)
    \quad\Rightarrow\quad
    A^a_{\mu}\mapsto A^a_{\mu}+q^a\partial_{\mu}\chi.
\end{align}

\textbf{[CRITICAL NOTE]}: the gauge layer is retained only as a roadmap statement. Running couplings, anomaly cancellation, and realistic mixing matrices remain open work.

\section{Dark-sector and galactic canonization execution package (v0.8.x)}
\label{sec:dark-sector-execution-package}

\subsection{Stage-0 freeze: topology toolkit and benchmark protocol}
\textbf{[DERIVED] Scope freeze:} the mandatory topology toolkit for v0.8.x canonization is fixed to:
\begin{enumerate}
    \item Gauss linking integral,
    \item helicity in ideal fluids,
    \item thin-tube helicity reduction,
    \item Hopf invariant,
    \item basic torus-knot invariants $(c,b,g)$,
    \item Calugareanu--White--Fuller relation $Lk=Tw+Wr$.
\end{enumerate}

\textbf{[DERIVED] Notation freeze:} one notation family is used in this package:
\[
\rhoF,\quad v_\theta(r),\quad p_{\mathrm{swirl}}(r),\quad
\mathcal{H}(K),\quad N_u^{(1)}(K),\quad \Gamma,\quad \kappa,\quad \SwirlClock.
\]
No parallel symbol aliases are introduced for the same quantity.

\textbf{[DERIVED] Mandatory benchmark protocol template:}
\begin{enumerate}
    \item \textbf{Inputs:} canonical constants, topology identifiers, profile parameters, and dataset/version tags.
    \item \textbf{Procedure:} exact script or package path, run command, and branch/commit hash.
    \item \textbf{Outputs:} archived CSV files with fixed column schema and an attached metadata note.
    \item \textbf{Uncertainty:} first-order propagation for calibrated constants and sensitivity sweep on profile parameters.
    \item \textbf{Rerun:} deterministic rerun steps and acceptance tolerances.
\end{enumerate}

\textbf{[CANON WORKFLOW NOTE]}: every promoted quantitative claim in this package is invalid without an explicit input-to-CSV trace.

\subsection{Stage-1 freeze: dark-sector Euler pressure anchor}
\textbf{[ORTHODOX] Euler radial balance (steady, axisymmetric, azimuthal-dominant):}
\begin{align}
    \frac{1}{\rhoF}\frac{dp_{\mathrm{swirl}}}{dr}
    =
    \frac{v_\theta(r)^2}{r}.
    \label{eq:exec_swirl_pressure}
\end{align}

\textbf{[DERIVED] Flat-tail limit (constant tail speed):}
\begin{align}
    v_\theta(r)\to v_0
    \quad\Longrightarrow\quad
    p_{\mathrm{swirl}}(r)
    =
    p_0 + \rhoF v_0^2\ln\!\left(\frac{r}{r_0}\right).
\end{align}

\textbf{[DERIVED] Assumption block:}
incompressible, inviscid, stationary, axisymmetric, azimuthal-dominant flow, with validity restricted to regions where radial drift and stratification corrections are subleading.

\textbf{[DERIVED] Dimensional check:}
\[
\left[\frac{1}{\rhoF}\frac{dp}{dr}\right]
=
\frac{\mathrm{kg\,m^{-2}\,s^{-2}}}{\mathrm{kg\,m^{-3}}}
=
\mathrm{m\,s^{-2}}
=
\left[\frac{v^2}{r}\right].
\]

\textbf{[SPECULATIVE] SST interpretation:}
Eq.~\eqref{eq:exec_swirl_pressure} is interpreted as a classifier/driver layer for dark-sector support. This interpretation is model-level and separate from the orthodox derivation.

\textbf{[CRITICAL NOTE]}: Applying Eq.~\eqref{eq:exec_swirl_pressure} unscreened with the microscopic $\rhoF=7.0\times10^{-7}\,\mathrm{kg\,m^{-3}}$ at galactic scales ($v_\theta\sim200\,\mathrm{km\,s^{-1}}$, $r\sim10\,\mathrm{kpc}$) gives $dp/dr\sim10^{-16}\,\mathrm{Pa\,m^{-1}}$ and $\Delta p\sim2.8\times10^{3}\,\mathrm{Pa}$ over $1\,\mathrm{kpc}$, far above the observed ISM. The swirl pressure is not the thermal ISM pressure, but no screening/scaling law for $\rhoF$ across scales is yet derived; this sector stays research-track and is falsifiable on exactly this point.

\textbf{[DERIVED] Attached observables/falsifiers:}
\begin{enumerate}
    \item Rotation-curve consistency after pressure reconstruction from archived $v_\theta(r)$ data.
    \item Null falsifier: systematic mismatch in the reconstructed pressure gradient sign or scale across validated radial windows.
\end{enumerate}

\subsection{Stage-2 freeze: measurable dark taxonomy}
\textbf{[SPECULATIVE] Classifier variables:}
symmetry class (including amphichirality status), chirality sign/state, helicity-balance indicator $\Delta_H(K)$, and low-order instability count $N_u^{(1)}(K)$.

\textbf{[SPECULATIVE] Decision rules (symmetry-first, then stability-filtered):}
\begin{align}
    \text{Dark}
    &:\ \text{amphichiral candidate}
    \ \wedge\
    \Delta_H(K)\ll 1
    \ \wedge\
    N_u^{(1)}(K)=0,\\
    \text{Quasi-dark}
    &:\ \text{near-balanced symmetry/helicity}
    \ \wedge\
    0 < N_u^{(1)}(K)\le 2,\\
    \text{Visible}
    &:\ \text{chiral-dominant and/or helicity-unbalanced sectors}.
\end{align}

\textbf{[CANON WORKFLOW NOTE] Ambiguity tie-break:}
if classes overlap, choose the class with the strongest conjunction score in the order
\[
\text{symmetry gate} \rightarrow \text{helicity balance gate} \rightarrow \text{stability gate}.
\]

\textbf{[SPECULATIVE] Baseline templates and falsifiers:}
\begin{itemize}
    \item \textbf{Visible baseline:} trefoil-like chiral template; falsifier = persistent amphichiral-plus-balanced signature under validated data updates.
    \item \textbf{Quasi-dark baseline:} Borromean-type weak-coupling template; falsifier = stable zero-instability assignment across independent scans.
    \item \textbf{Dark baseline:} figure-eight/amphichiral template; falsifier = robust nonzero chirality or high-instability assignment.
\end{itemize}

\textbf{[CANON WORKFLOW NOTE]}:
amphichiral $\Rightarrow$ dark remains a classifier proposal tested against observables, not a closed theorem.

\subsection{Stage-3 freeze: galactic law branch gate and closure route}
\textbf{[CANON WORKFLOW NOTE] Branch-gate decision:} Branch A (legacy profile refine) is selected for current v0.8.x closure because it preserves tested limits while allowing benchmark traceability.

Working profile:
\begin{align}
    v_\theta(r)
    =
    \frac{C_{\mathrm{core}}}{\sqrt{1+(r_c/r)^2}}
    +
    C_{\mathrm{tail}}\left(1-e^{-r/r_c}\right).
\end{align}

\textbf{[DERIVED] Required checks:}
small-$r$ behavior, flat-tail behavior, parameter semantics, and pressure reconstruction via Eq.~\eqref{eq:exec_swirl_pressure}.

\textbf{[DERIVED] Failure modes:}
\begin{enumerate}
    \item non-physical parameter combinations (wrong sign or unstable growth),
    \item failure to satisfy tail-limit behavior in benchmark windows,
    \item inability to reconstruct a consistent pressure profile from the same archived inputs.
\end{enumerate}

\subsection{Stage-4 freeze: probe falsification map}
\textbf{[CANON WORKFLOW NOTE]} probe classes are fixed as astrophysical, laboratory pressure/swirl analog, HV-coil/thrust-inspired, optical/torsional coupling, and null dark-sector tests.

\begin{table}[h]
\centering
\begin{tabular}{p{2.7cm}p{3.1cm}p{2.5cm}p{2.2cm}p{2.2cm}}
\hline
Probe class & Equation under test & Null-result meaning & Primary confounder & Expected sign/trend \\
\hline
Astrophysical rotation curves & Eq.~\eqref{eq:exec_swirl_pressure} + profile closure & weakens pressure-law applicability window & baryonic-model bias & reconstructed $dp/dr$ consistent with $v_\theta^2/r$ \\
Lab swirl analog & Euler radial balance scaling & weakens scale-transfer claim & viscosity/edge forcing & monotone pressure rise with imposed swirl \\
HV-coil/thrust-inspired & circulation-pressure bridge claims & kills quantized-bridge claim if no scaling & EM/thermal artifacts & response tracks circulation proxy \\
Optical/torsional coupling & swirl-clock bridge signatures & weakens coupling interpretation & detector drift/noise & phase/timing shift with controlled swirl proxy \\
Null dark tests & taxonomy conjunction rules & kills classifier branch if conjunction fails & mislabelled topology class & class score decreases under contradictory diagnostics \\
\hline
\end{tabular}
\caption{Frozen equation-to-probe map for falsification-first workflow in the v0.8.x dark-sector program.}
\label{tab:exec-probe-map}
\end{table}

\textbf{[CANON WORKFLOW NOTE]}: probe outcomes can support or reject already-stated equations, but they do not replace derivation-level closure.

\subsection{Promotion snapshot after execution}
\textbf{[DERIVED]} post-freeze status in this package:
\begin{itemize}
    \item Item 6 (topology toolkit): CANON-ready in scope and notation.
    \item Item 5 (benchmark protocol): CANON-CANDIDATE pending repository-wide protocol reuse.
    \item Item 1 (Euler pressure anchor): CANON-CANDIDATE with explicit assumptions, limits, and falsifier path.
    \item Item 2 (dark taxonomy): BRIDGE with measurable classifier gates and falsifiers.
    \item Item 3 (galactic swirl law): BRIDGE with branch gate closed and benchmark route fixed.
    \item Item 4 (probe map): CANON-BRIDGE with equation/null-result mapping frozen.
\end{itemize}

%===========================================================

\section{Candidate CANON-v0.8.x Modules from Swirl-Flux, Envelope, and Reconnection Analysis}
    \label{app:sst_v08_candidate_modules}

% --- SST macro prelude ---
    \providecommand{\swirlarrow}{%
        \mathchoice{\mkern-2mu\scriptstyle\boldsymbol{\circlearrowleft}}%
        {\mkern-2mu\scriptstyle\boldsymbol{\circlearrowleft}}%
        {\mkern-2mu\scriptscriptstyle\boldsymbol{\circlearrowleft}}%
        {\mkern-2mu\scriptscriptstyle\boldsymbol{\circlearrowleft}}%
    }
    \providecommand{\vswirl}{\mathbf{v}_{\!\swirlarrow}}
\providecommand{\vchar}{v_{\!\boldsymbol{\circlearrowleft}}^{\ast}}
\providecommand{\uswirl}{\mathbf{u}_{\!\boldsymbol{\circlearrowleft}}}
\providecommand{\omegaswirl}{\boldsymbol{\omega}_{\!\boldsymbol{\circlearrowleft}}}
    \providecommand{\rhoF}{\rho_{\!f}}
    \providecommand{\rhoE}{\rho_{\!E}}
    \providecommand{\rhoM}{\rho_{\!m}}
    \providecommand{\rhocore}{\rho_{\mathrm{core}}}
    \providecommand{\rc}{r_c}
    \providecommand{\Fswirlmax}{F_{\mathrm{swirl}}^{\max}}
    
    \subsection{Status Classification}
        
        The following modules are proposed as CANON-v0.8.x candidates.
        
        \begin{enumerate}
            \item The Schrödinger kinetic coefficient is retained as an effective
            envelope stiffness,
            \[
                \mathcal S_e=\frac{\hbar^2}{2M_e}.
            \]
            The replacement \(-\vswirl^2\nabla^2/2\) is rejected on dimensional grounds.
            
            \item The Meissner effect is interpreted as macroscopic swirl-flux phase
            locking in a coherent paired condensate.
            
            \item R/T wave--string interaction is represented by conservation of
            radiative pseudomomentum plus vorticity impulse.
            
            \item Reconnection is excluded in the ideal sector but allowed in
            measurement, decay, and dissipative sectors, where helicity is redistributed
            between topology, writhe, twist, radiation, and dissipation.
        \end{enumerate}
    
    \subsection{Envelope Stiffness and the Rejection of the \texorpdfstring{\(\vswirl^2\nabla^2\)}{vswirl2 Laplacian} Kinetic Term}
        
        For an effective R/T envelope amplitude \(\Psi(\mathbf r,t)\), the kinetic
        coefficient must have dimension \(\mathrm{J\,m^2}\). The correct envelope
        stiffness is therefore
        \[
            \mathcal S_e
            \equiv
            \frac{\hbar^2}{2M_e},
            \qquad
            [\mathcal S_e]=\mathrm{J\,m^2}.
        \]
        The effective low-energy envelope equation is
        \[
            i\hbar\frac{\partial\Psi}{\partial t}
            =
            \left[
                -\mathcal S_e\nabla^2
            +
            V_{\mathrm{SST}}(\mathbf r,t)
            \right]\Psi .
        \]
        The operator \(-\vswirl^2\nabla^2/2\) is not a Schrödinger kinetic energy
        operator because
        \[
            \left[\vswirl^2\nabla^2\right]
            =
            \frac{\mathrm{m^2}}{\mathrm{s^2}}\frac{1}{\mathrm{m^2}}
            =
            \mathrm{s^{-2}},
        \]
        which is not an energy.
        
        For tunneling through a one-dimensional effective SST barrier,
        \[
            T
            \sim
            \exp\left(
                    -2\int_{x_1}^{x_2}\kappa_{\mathrm{SST}}(x)\,dx
            \right),
        \]
        where
        \[
            \kappa_{\mathrm{SST}}(x)
            =
            \sqrt{
                \frac{V_{\mathrm{SST}}(x)-E}{\mathcal S_e}
            }
            =
            \sqrt{
                \frac{2M_e\left[V_{\mathrm{SST}}(x)-E\right]}{\hbar^2}
            }.
        \]
        This is dimensionally consistent because
        \[
            \left[
                \frac{V-E}{\mathcal S_e}
            \right]
            =
            \mathrm{m^{-2}},
            \qquad
            [\kappa_{\mathrm{SST}}]=\mathrm{m^{-1}}.
        \]
    
    \subsection{Numerical Bridge Anomaly}
        
        A numerically close but dimensionally non-identical bridge is
        \[
            \mathcal B_e
            =
            \frac{\Fswirlmax \rc^3}
            {5\lambda_c\vchar}.
        \]
        Using
        \[
            \lambda_c=\frac{h}{M_e c},
        \]
        one obtains
        \[
            \frac{\hbar^2}{2M_e}
            =
            6.10426431\times10^{-39}\,\mathrm{J\,m^2},
        \]
        whereas
        \[
            \mathcal B_e
            =
            6.12394931\times10^{-39}\,\mathrm{J\,s}.
        \]
        The numerical ratio is
        \[
            \frac{\mathcal B_e}{\mathcal S_e}\bigg|_{\mathrm{SI\ numbers}}
            =
            1.00322479.
        \]
        This is classified as a Research Track numerical bridge, not a canonical
        identity, because
        \[
            [\mathcal B_e]=\mathrm{J\,s}
            \neq
            \mathrm{J\,m^2}=[\mathcal S_e].
        \]
        
        The dimensionally meaningful relation is instead
        \[
            5\lambda_c\vchar
            =
            5\alpha\frac{h}{2M_e},
        \]
        with both sides having dimension \(\mathrm{m^2\,s^{-1}}\).
    
    \subsection{Swirl-Flux Locking and the Meissner Effect}
        \textbf{[CRITICAL NOTE]}: The values $2M_e$ and $q_{\mathrm{pair}}=-2e$ are imported from the orthodox Cooper-pair description. Per Sec.~8.6 (Charge and Circulation), SST does not yet derive gauge charges from circulation alone, so this module is a $[\text{CALIBRATED}]$ bridge using empirical pairing inputs, not a first-principles flux-exclusion derivation.
        
        For a coherent paired electronic condensate of mass \(2M_e\) and charge
        \(q_{\mathrm{pair}}=-2e\), phase single-valuedness gives
        \[
            \oint_{\partial\Sigma}
            \left(
                2M_e\mathbf v_{\mathrm{pair}}
                +
                q_{\mathrm{pair}}\mathbf A
            \right)\cdot d\boldsymbol\ell
            =
            nh .
        \]
        Since \(q_{\mathrm{pair}}=-2e\),
        \[
            2M_e\Gamma_{\mathrm{pair}}
            -
            2e\Phi_B
            =
            nh .
        \]
        Thus
        \[
            \Phi_B
            =
            \frac{M_e}{e}\Gamma_{\mathrm{pair}}
            -
            n\frac{h}{2e}.
        \]
        In magnitude,
        \[
            |\Phi_B|
            =
            \left|
                n\Phi_0
            -
                \frac{M_e}{e}\Gamma_{\mathrm{pair}}
            \right|,
            \qquad
            \Phi_0=\frac{h}{2e}.
        \]
        For a non-rotating coherent bulk,
        \[
            \Gamma_{\mathrm{pair}}\to 0,
        \]
        and therefore
        \[
            |\Phi_B|=n\Phi_0.
        \]
        The Meissner effect is then interpreted in SST as exclusion of non-quantized
        magnetic flux from a phase-locked coherent R/T condensate.
    
    \subsection{R/T Pseudomomentum--Impulse Exchange}
        
        For a closed R/T interaction domain,
        \[
            \frac{d}{dt}
            \left(
                \mathbf P_R+\mathbf I_T
            \right)
            =
            0.
        \]
        Here \(\mathbf P_R\) denotes radiative R-phase pseudomomentum and
        \(\mathbf I_T\) denotes T-phase vorticity impulse,
        \[
            \mathbf I_T
            =
            \frac{\rhoF}{2}
            \int_{\Omega}
            \mathbf x\times\boldsymbol\omega\,d^3x .
        \]
        The dimension is
        \[
            [\mathbf I_T]
            =
            \mathrm{kg\,m\,s^{-1}},
        \]
        so \(\mathbf I_T\) is a true momentum-like quantity.
        
        For a thin circular swirl string of circulation \(\Gamma\) and radius \(R\),
        \[
            \mathbf I_T
            =
            \rhoF\Gamma\pi R^2\hat{\mathbf n}.
        \]
        The core circulation scale is
        \[
            \Gamma_0
            =
            2\pi \rc\vchar
            =
            9.68361920\times10^{-9}\,\mathrm{m^2\,s^{-1}}.
        \]
    
    \subsection{Reconnection and Helicity Redistribution}
        
        In the ideal sector, topology is conserved:
        \[
            \frac{D\mathcal K}{Dt}=0,
        \]
        under the assumptions
        \[
            \nabla\cdot\mathbf v=0,
            \qquad
            \nu=0,
            \qquad
            \mathrm{barotropic},
            \qquad
            \mathrm{no\ reconnection}.
        \]
        In the dissipative, measurement, or decay sector, reconnection is allowed only as a non-ideal event:
        \[
            \frac{D\mathcal K}{Dt}\neq0.
        \]
        This is not a free creation mechanism. A reconnection claim must state the trigger class, e.g. \(\nu_{\rm eff}>0\), a boundary injection, a phase-slip singularity, a resolved-core collapse, or an explicit external forcing term.
        The helicity budget is written
        \[
            \Delta H_{\mathrm{topology}}
            +
            \Delta H_{\mathrm{writhe}}
            +
            \Delta H_{\mathrm{twist}}
            +
            \Delta H_{\mathrm{radiative}}
            =
            -\Delta H_{\mathrm{diss}}.
        \]
        In the weak-dissipation limit,
        \[
            \Delta H_{\mathrm{topology}}
            +
            \Delta H_{\mathrm{writhe}}
            +
            \Delta H_{\mathrm{twist}}
            +
            \Delta H_{\mathrm{radiative}}
            \approx
            0.
        \]
        For filamentary structures,
        \[
            H_{\mathrm{centerline}}
            =
            \sum_{i\neq j}
            \Gamma_i\Gamma_j\mathrm{Lk}_{ij}
            +
            \sum_i
            \Gamma_i^2
            \left(
                \mathrm{Wr}_i+\mathrm{Tw}_i
            \right).
        \]
    
    \subsection{Spherical Pressure Envelopes}
        
        For inviscid incompressible SST flow,
        \[
            \nabla p_{\mathrm{SST}}
            =
            \rhoF(\mathbf v\cdot\nabla)\mathbf v .
        \]
        For an azimuthal swirl profile \(v_\theta(r)\),
        \[
            \frac{dp}{dr}
            =
            \rhoF\frac{v_\theta^2(r)}{r}.
        \]
        If
        \[
            v_\theta(r)=\frac{\Gamma}{2\pi r},
        \]
        then
        \[
            \frac{dp}{dr}
            =
            \rhoF\frac{\Gamma^2}{4\pi^2r^3},
        \]
        and hence
        \[
            p(r)
            =
            p_\infty
            -
            \frac{\rhoF\Gamma^2}{8\pi^2r^2}.
        \]
        Therefore
        \[
            \Delta p(r)
            =
            p_\infty-p(r)
            =
            \frac{\rhoF\Gamma^2}{8\pi^2r^2}
            =
            \frac{1}{2}\rhoF v_\theta^2(r).
        \]
        
        At the canonical swirl speed,
        \[
            \frac{1}{2}\rhoF\vchar{}^2
            =
            4.18774392\times10^5\,\mathrm{Pa}.
        \]
        At the core-density stress scale,
        \[
            \rhocore\vchar{}^2
            =
            4.65848920\times10^{30}\,\mathrm{Pa}.
        \]
        The canonical maximum swirl force is recovered as
        \[
            \Fswirlmax
            =
            \rhocore\vchar{}^2\pi\rc^2
            =
            29.053507\,\mathrm{N}.
        \]
    
    \subsection{Golden-Layer Scaling}
        
        The universal rule
        \[
            E_n=\varphi^nE_0
        \]
        is rejected as too strong. The weaker Research Track ansatz is
        \[
            E_{K,k}
            =
            E_K^{(0)}\varphi^{-2k},
            \qquad
            k\in\mathbb Z_{\ge0}.
        \]
        Here \(K\) labels a knot family and \(k\) labels a relaxation layer inside
        that family.
    
    \subsection{Non-Holonomic Gravity Claims}
        
        Podkletnov-type anti-gravity claims are not adopted as Canon. At most, the
        formal analogy
        \[
            \mathrm{non\mbox{-}holonomic\ twist}
            \longleftrightarrow
            \mathrm{foliation\ twist/vorticity}
        \]
        may be retained in a historical or speculative appendix.
        No experimental anti-gravity claim is inferred.
        
        \ifresearchtrackincluded\else
\begin{thebibliography}{99}
            
            \bibitem{LondonLondon1935}
            F. London and H. London,
            ``The Electromagnetic Equations of the Supraconductor,''
            \textit{Proceedings of the Royal Society A},
            vol. 149, no. 866, pp. 71--88, 1935.
            DOI: \texttt{10.1098/rspa.1935.0048}.
            
            \bibitem{Tinkham1996}
            M. Tinkham,
            \textit{Introduction to Superconductivity},
            2nd ed.,
            McGraw--Hill, New York, 1996.
            ISBN: \texttt{978-0486435039}.
            
            \bibitem{Onsager1949}
            L. Onsager,
            ``Statistical hydrodynamics,''
            \textit{Il Nuovo Cimento},
            vol. 6, Suppl. 2, pp. 279--287, 1949.
            DOI: \texttt{10.1007/BF02780991}.
            
            \bibitem{Feynman1955}
            R. P. Feynman,
            ``Application of Quantum Mechanics to Liquid Helium,''
            in \textit{Progress in Low Temperature Physics},
            vol. 1, pp. 17--53, Elsevier, 1955.
            DOI: \texttt{10.1016/S0079-6417(08)60077-3}.
            
            \bibitem{Vinen2004}
            W. F. Vinen,
            ``The Physics of Superfluid Helium,''
            CERN Yellow Reports / Lecture Notes, 2004.
            Permalink: \texttt{https://cds.cern.ch/record/808382/files/p363.pdf}.
            
            \bibitem{BuhlerMcIntyre2005}
            O. B{\"u}hler and M. E. McIntyre,
            ``Wave capture and wave--vortex duality,''
            \textit{Journal of Fluid Mechanics},
            vol. 534, pp. 67--95, 2005.
            DOI: \texttt{10.1017/S0022112005004374}.
            
            \bibitem{KlecknerIrvine2013}
            D. Kleckner and W. T. M. Irvine,
            ``Creation and dynamics of knotted vortices,''
            \textit{Nature Physics},
            vol. 9, pp. 253--258, 2013.
            DOI: \texttt{10.1038/nphys2560}.
            
            \bibitem{KlecknerKauffmanIrvine2016}
            D. Kleckner, L. H. Kauffman, and W. T. M. Irvine,
            ``How superfluid vortex knots untie,''
            \textit{Nature Physics},
            vol. 12, pp. 650--655, 2016.
            DOI: \texttt{10.1038/nphys3679}.
            
            \bibitem{KellerCheviakov2024}
            J. M. Keller and A. F. Cheviakov,
            ``Exact spherical vortex-type equilibrium flows in fluids and plasmas,''
            \textit{Fundamental Plasma Physics},
            vol. 11, article 100063, 2024.
            DOI: \texttt{10.1016/j.fpp.2024.100063}.
            
            \bibitem{Brickman1969}
            N. F. Brickman,
            ``Rotating Superconductors,''
            \textit{Physical Review},
            vol. 184, pp. 460--470, 1969.
            DOI: \texttt{10.1103/PhysRev.184.460}.
        
        \end{thebibliography}
\fi

        
        
        
% ============================================================
% Appendix: Mode-Locked Swirl-Coil Excitation
% legacy insertion block, retained in v0.8.17 research track:
% Canon v0.8.1 insertion block
% ============================================================

% Safe local macro guards; harmless if the Canon preamble already defines these.
\providecommand{\swirlarrow}{\mathchoice{\mkern-2mu\scriptstyle\boldsymbol{\circlearrowleft}}{\mkern-2mu\scriptstyle\boldsymbol{\circlearrowleft}}{\mkern-2mu\scriptscriptstyle\boldsymbol{\circlearrowleft}}{\mkern-2mu\scriptscriptstyle\boldsymbol{\circlearrowleft}}}
\providecommand{\vswirl}{\mathbf{v}_{\!\swirlarrow}}
\providecommand{\rhoF}{{\rho_{\!f}}}
\providecommand{\rhoE}{\rho_{\!E}}
\providecommand{\rhoM}{\rho_{\!m}}
\providecommand{\rc}{{r_c}}
\providecommand{\vchar}{v_{\!\boldsymbol{\circlearrowleft}}^{\ast}}
\providecommand{\vnorm}{\vchar}

\section{Mode-Locked Swirl-Coil Excitation and \texorpdfstring{$\vchar=f\Delta x$}{v=f dx} Metrology}
\label{sec:appendix_mode_locked_swirl_coil}

\textbf{Status.} \textbf{[RESEARCH TRACK / CANON-CANDIDATE]}.
This appendix formulates a falsifiable experimental protocol for testing whether externally driven multi-phase coil geometries can inject coherent swirl modes into a bounded measurement region. The construction is not assumed as a new force law. It is a mode-selection and metrology rule for identifying when the canonical characteristic swirl speed
\begin{equation}
\vnorm = 1.09384563\times 10^{6}\,\mathrm{m\,s^{-1}}
\end{equation}
is sampled by a laboratory-scale frequency--length product.

\subsection{Motivation}

The basic dimensional identity
\begin{equation}
\vnorm = f\,\Delta x
\label{eq:appendix_swirl_speed_frequency_length}
\end{equation}
has the units of speed and is therefore not, by itself, a dynamical claim. Its SST use is metrological: a driven system is said to probe a swirl-resonant channel only when a temporal frequency $f$ and a geometrically defined spatial period $\Delta x$ satisfy Eq.~\eqref{eq:appendix_swirl_speed_frequency_length} within experimental bandwidth.

For a circular or quasi-circular coil, the relevant spatial period is generally not an arbitrary external distance. It is the azimuthal wavelength selected by the coil radius and spatial mode number. Hence the physical coil radius enters the resonance condition directly.

\subsection{Ring-mode closure condition}

Let $R_{\mathrm{coil}}$ be the effective radius of the active winding, $f_0$ the electrical base frequency, $n\in\mathbb{N}$ the temporal harmonic, and $m\in\mathbb{Z}\setminus\{0\}$ the azimuthal spatial mode number around the coil. The spatial wavelength of mode $m$ along the coil circumference is
\begin{equation}
\Delta x_m
=
\frac{2\pi R_{\mathrm{coil}}}{|m|}.
\label{eq:appendix_spatial_period}
\end{equation}
The (n)-th temporal harmonic has frequency
\begin{equation}
f_n = n f_0.
\end{equation}
The dimensionless swirl-lock parameter is therefore
\begin{equation}
\boxed{
\eta_{n,m}
=
\frac{f_n \Delta x_m}{\vnorm}
=
\frac{n f_0}{\vnorm}
\frac{2\pi R_{\mathrm{coil}}}{|m|}
}.
\label{eq:appendix_eta_nm}
\end{equation}
A ring mode is swirl-locked when
\begin{equation}
\boxed{\eta_{n,m}=1.}
\label{eq:appendix_eta_lock}
\end{equation}
Equivalently, the resonant coil radius for a given $(n,m,f_0)$ is
\begin{equation}
\boxed{
R_{\mathrm{coil}}^{(n,m)}
=
\frac{|m|\,\vnorm}{2\pi n f_0}
}.
\label{eq:appendix_resonant_radius}
\end{equation}

\paragraph{Dimensional check.}
\[
[R_{\mathrm{coil}}]
=
\frac{\mathrm{m\,s^{-1}}}{\mathrm{s^{-1}}}
=
\mathrm{m}.
\]
Thus Eq.~\eqref{eq:appendix_resonant_radius} is dimensionally well-defined.

\subsection{Numerical calibration examples}

For the canonical value
\[
\vnorm = 1.09384563\times10^6\,\mathrm{m\,s^{-1}},
\]
and a base drive frequency
\[
f_0 = 150\,\mathrm{kHz},
\]
the first three $m=1$ radius locks are
\begin{align}
R_{\mathrm{coil}}^{(1,1)}
&=
\frac{1.09384563\times10^6}{2\pi(1)(150000)}
=
1.1606\,\mathrm{m},
\\
R_{\mathrm{coil}}^{(2,1)}
&=
0.5803\,\mathrm{m},
\\
R_{\mathrm{coil}}^{(3,1)}
&=
0.3869\,\mathrm{m}.
\end{align}
For a smaller experimental coil with
\[
R_{\mathrm{coil}}=0.150\,\mathrm{m},
\]
the $(n=3,\ m=1)$ lock frequency is
\begin{equation}
f_0
=
\frac{\vnorm}{2\pi n R_{\mathrm{coil}}}
=
\frac{1.09384563\times10^6}
{2\pi(3)(0.150)}
=
3.8687\times10^5\,\mathrm{Hz}.
\end{equation}
Hence a $0.150\,\mathrm{m}$ coil driven at $150\,\mathrm{kHz}$ does not satisfy the $(n=3,m=1)$ ring-lock condition. Its lock parameter is
\begin{equation}
\eta_{3,1}
=
\frac{3(150000)\,2\pi(0.150)}
{1.09384563\times10^6}
=
0.3877.
\end{equation}
This is sub-resonant for the $(3,1)$ ring mode, although it may still couple to different $(n,m)$ channels.

\subsection{Three-phase selection rule}

Consider a three-phase coil drive with phase index $q=0,1,2$. If the $q$-th phase is shifted by $2\pi q/3$, the discrete phase sum for temporal harmonic $n$ and spatial azimuthal mode $m$ is
\begin{equation}
T_{n,m}^{(3\phi)}
=
\frac{1}{3}
\left|
\sum_{q=0}^{2}
\exp!\left[
i q \frac{2\pi}{3}(m-n)
\right]
\right|.
\label{eq:appendix_three_phase_factor}
\end{equation}
The exact selection rule is
\begin{equation}
\boxed{
T_{n,m}^{(3\phi)}
=
\begin{cases}
1, & m\equiv n \pmod{3},\\
0, & m\not\equiv n \pmod{3}.
\end{cases}
}
\label{eq:appendix_three_phase_rule}
\end{equation}
Thus a three-phase winding does not merely ``add power''; it filters the allowed temporal--spatial mode pairs.

\subsection{Double-stack phase-shift filter}

Let the coil be built as two vertically separated stacks, with the second stack rotated azimuthally by
\begin{equation}
\delta = 30^\circ = \frac{\pi}{6}.
\end{equation}
Let $\psi$ denote the electrical inter-stack phase offset. The two-stack spatial filter is
\begin{equation}
L_m(\delta,\psi)
=
\left|
1+\exp[i(m\delta+\psi)]
\right|.
\label{eq:appendix_stack_filter}
\end{equation}
For equal stack polarity and zero additional electrical offset $(\psi=0)$, the destructive condition is
\begin{equation}
m\delta=(2q+1)\pi,
\qquad q\in\mathbb{Z}.
\end{equation}
With $\delta=\pi/6$, this gives
\begin{equation}
\boxed{
m = 6,18,30,\ldots
}
\end{equation}
Therefore a $30^\circ$ double-stack geometry exactly suppresses the sixth azimuthal spatial harmonic and its odd multiples in this idealized equal-amplitude model.

\subsection{Golden-duty PWM spectrum}

Let the duty cycle be the inverse golden-square fraction
\begin{equation}
d = \varphi^{-2}
=
\left(\frac{1+\sqrt{5}}{2}\right)^{-2}
=
0.381966011\ldots .
\end{equation}
For an ideal unipolar rectangular pulse train, the magnitude of the $n$-th Fourier coefficient, ignoring the DC normalization convention, is proportional to
\begin{equation}
P_n(d)
=
\left|
\frac{\sin(\pi n d)}{\pi n}
\right|.
\label{eq:appendix_pwm_spectrum}
\end{equation}
For $d=\varphi^{-2}$, the first seven relative amplitudes normalized to $P_1$ are approximately
\begin{center}
\begin{tabular}{c|c|c}
\hline
$n$ & $P_n(d)$ & $P_n(d)/P_1(d)$ \\
\hline
1 & 0.296675 & 1.000 \\
2 & 0.107508 & 0.362 \\
3 & 0.046948 & 0.158 \\
4 & 0.079273 & 0.267 \\
5 & 0.017794 & 0.060 \\
6 & 0.042102 & 0.142 \\
7 & 0.038864 & 0.131 \\
\hline
\end{tabular}
\end{center}
Thus golden-duty PWM strongly suppresses the fifth harmonic, but does not by itself null the sixth and seventh harmonics. In the proposed three-phase double-stack geometry, the suppression of the fifth, sixth, and seventh bands must be understood as a composite filter effect:
\begin{equation}
\boxed{
\mathcal{C}_{n,m}
=
P_n(d)\,
T_{n,m}^{(3\phi)}\,
S_m(\mathrm{P40},+11,-9)\,
L_m(\delta,\psi)\,
G_m(k_nR_{\mathrm{coil}})
}.
\label{eq:appendix_composite_filter}
\end{equation}
Here:
\begin{itemize}
\item $P_n(d)$ is the PWM harmonic amplitude;
\item $T_{n,m}^{(3\phi)}$ is the three-phase selection factor;
\item $S_m(\mathrm{P40},+11,-9)$ is the spatial winding factor of the SawShape path;
\item $L_m(\delta,\psi)$ is the double-stack filter;
\item $G_m(k_nR_{\mathrm{coil}})$ is the geometric near-field coupling factor.
\end{itemize}
The wave number of the $n$-th temporal harmonic is
\begin{equation}
k_n
=
\frac{2\pi n f_0}{\vnorm}.
\end{equation}

\subsection{SawShape winding factor for \texorpdfstring{$\mathrm{P40},+11,-9$}{P40,+11,-9}}

For a discrete winding on $N=40$ angular sites, define the node angle
\begin{equation}
\theta_j=\frac{2\pi p_j}{40},
\qquad p_j\in\mathbb{Z}/40\mathbb{Z}.
\end{equation}
A SawShape path with alternating skip sequence $(+11,-9)$ can be represented by
\begin{equation}
p_{j+1}
=
p_j + s_j
\pmod{40},
\qquad
s_j\in\{+11,-9\}.
\end{equation}
The normalized spatial Fourier factor is then
\begin{equation}
\boxed{
S_m(\mathrm{P40},+11,-9)
=
\left|
\frac{1}{N_s}
\sum_{j=0}^{N_s-1}
w_j
\exp(-i m \theta_j)
\right|
}.
\label{eq:appendix_sawshape_factor}
\end{equation}
Here $w_j$ encodes segment weight, orientation, current direction, layer index, and local segment length. This factor is experimentally accessible: it can be inferred by measuring the spatial FFT of the near-field magnetic map over a circular sampling path.

\subsection{Canonical interpretation}

The proposed coil does not directly generate a canonical force. Instead, it prepares a structured, rotating, mode-filtered boundary condition for the local field. In SST language, the coil acts as a laboratory-scale swirl-mode injector.

The relevant energy-density comparison is
\begin{equation}
\rhoE
=
\frac{1}{2}\rhoF \vnorm^2.
\end{equation}
Using
\[
\rhoF = 7.0\times10^{-7}\,\mathrm{kg\,m^{-3}},
\qquad
\vnorm = 1.09384563\times10^{6}\,\mathrm{m\,s^{-1}},
\]
one obtains
\begin{equation}
\rhoE
=
\frac{1}{2}
(7.0\times10^{-7})
(1.09384563\times10^{6})^2
=
4.187\times10^{5}\,\mathrm{J\,m^{-3}}.
\end{equation}
The corresponding mass-equivalent density is
\begin{equation}
\rhoM
=
\frac{\rhoE}{c^2}
=
\frac{4.187\times10^{5}}
{(2.99792458\times10^{8})^2}
=
4.659\times10^{-12}\,\mathrm{kg\,m^{-3}}.
\end{equation}
These values should be understood as canonical scale references, not as a claim that a laboratory coil reaches the full canonical swirl speed in ordinary electromagnetic matter.

\subsection{Experimental falsification protocol}

A coil experiment is SST-relevant only if it passes the following null hierarchy.

\begin{enumerate}
\item \textbf{Electrical null:} repeat with identical RMS current and power but randomized phase ordering.
\item \textbf{Geometric null:} replace SawShape $(\mathrm{P40},+11,-9)$ by a symmetry-preserving circular winding with the same resistance and inductance.
\item \textbf{Duty-cycle null:} compare $d=0.381966$ against $d=0.500$, $d=0.333$, and $d=0.250$ at fixed RMS current.
\item \textbf{Stack null:} compare $30^\circ$ double-stack shift against $0^\circ$, $15^\circ$, $45^\circ$, and $60^\circ$.
\item \textbf{Mode-lock scan:} sweep $f_0$ and test whether anomalous coherence peaks occur at
\[
f_0^{(n,m)}
=
\frac{|m|\,\vnorm}{2\pi n R_{\mathrm{coil}}}.
\]
\item \textbf{Environmental null:} repeat in still air, shielded enclosure, and, where safely possible, reduced pressure to remove ion wind and acoustic coupling.
\item \textbf{Force-decomposition null:} decompose any measured force into electromagnetic, thermal, acoustic, EHD, and mechanical vibration contributions before assigning a residual.
\end{enumerate}

Only a residual signal satisfying both
\begin{equation}
\eta_{n,m}\approx 1
\end{equation}
and
\begin{equation}
\mathcal{C}_{n,m}\neq 0
\end{equation}
is admissible as an SST research-track candidate.

\subsection{Vorticity-force decomposition layer}

Any measured force must first be written schematically as
\begin{equation}
\mathbf{F}_{\mathrm{meas}}
=
\mathbf{F}_{\mathrm{EM}}
+
\mathbf{F}_{\mathrm{EHD}}
+
\mathbf{F}_{\mathrm{thermal}}
+
\mathbf{F}_{\mathrm{acoustic}}
+
\mathbf{F}_{\mathrm{mechanical}}
+
\mathbf{F}_{\omega}
+
\mathbf{F}_{\mathrm{res}} .
\label{eq:appendix_force_budget}
\end{equation}
Here $\mathbf{F}_{\omega}$ denotes force reconstructed from ordinary vorticity transport and impulse balance. Only the remaining term $\mathbf{F}_{\mathrm{res}}$, after conventional reconstruction and uncertainty propagation, can be tested against the SST mode-locking rule. The canonical research criterion is therefore
\begin{equation}
\boxed{
\mathbf{F}_{\mathrm{res}}
=
\mathbf{F}_{\mathrm{res}}(\eta_{n,m},\mathcal{C}_{n,m})
\quad\text{with a peak near}\quad
\eta_{n,m}=1.
}
\end{equation}

\subsection{Recommended measured observables}

The minimum observable set is:
\begin{enumerate}
\item drive current $I_A(t),I_B(t),I_C(t)$;
\item voltage $V_A(t),V_B(t),V_C(t)$;
\item real input power $P_{\mathrm{in}}(t)$;
\item local magnetic field map $\mathbf{B}(\mathbf{x},t)$;
\item pickup-coil FFT around the active circumference;
\item accelerometer spectrum on the coil frame;
\item air temperature and pressure field if operated in air;
\item force sensor output with synchronized phase reference.
\end{enumerate}
A positive SST candidate requires a phase-stable spectral feature at a predicted $(n,m)$ lock that disappears under at least one of the symmetry-breaking null tests.

\subsection{Conclusion}

The 3-phase SawShape $(\mathrm{P40},+11,-9)$ double-stack coil with $30^\circ$ inter-stack rotation and $d=\varphi^{-2}$ PWM is canonically useful as a mode-filtered swirl-excitation apparatus. Its SST value is not that it assumes a new force, but that it makes the frequency--length criterion
\[
\vnorm=f\Delta x
\]
experimentally precise.

The essential design rule is
\[
\eta_{n,m}
=
\frac{n f_0}{\vnorm}
\frac{2\pi R_{\mathrm{coil}}}{|m|}
\approx 1,
\]
with observable coupling only when
\[
\mathcal{C}_{n,m}
=
P_n(d)\,
T_{n,m}^{(3\phi)}\,
S_m(\mathrm{P40},+11,-9)\,
L_m(\delta,\psi)\,
G_m(k_nR_{\mathrm{coil}})
\]
is nonzero. This converts the coil concept into a falsifiable SST metrology protocol.

% ============================================================
% Optional bibliography additions, if not already present
% ============================================================
%
% Add these entries to the existing thebibliography environment only if needed.
%
% \bibitem{Bracewell2000}
% R.~N. Bracewell,
% \newblock \emph{The Fourier Transform and Its Applications},
% \newblock 3rd ed., McGraw--Hill (2000).
%
% \bibitem{Balanis2016}
% C.~A. Balanis,
% \newblock \emph{Antenna Theory: Analysis and Design},
% \newblock 4th ed., Wiley (2016).
%
% \bibitem{KlecknerIrvine2013}
% D.~Kleckner and W.~T.~M. Irvine,
% \newblock Creation and dynamics of knotted vortices,
% \newblock \emph{Nature Physics} \textbf{9} (2013), 253--258,
% \newblock DOI: 10.1038/nphys2560.
%
% \bibitem{Gehlert2023}
% P.~Gehlert, I.~Andreu-Angulo, and H.~Babinsky,
% \newblock Vortex force decomposition---forces associated with individual elements of a vorticity field,
% \newblock \emph{Experiments in Fluids} \textbf{64} (2023), 112,
% \newblock DOI: 10.1007/s00348-023-03654-3.


% ==========================================================
% legacy insertion block, retained in v0.8.18 research track:
% SST CANON v0.8.18 -- Research Track Appendices
% Four topology/stability appendices extracted from knot/vorticity literature review.
% Integrated as a research-track section in SST_CANON-v0.8.18-research-track.tex.
% Assumes the v0.8.18 macro prelude is already loaded.
% ==========================================================

\section{Research Appendices: Topology, Stability, and Fluid Mechanics}

\subsection{Research Appendix: Non-Abelian Topological Charge for Protected Swirl Strings}
\label{app:nonabelian_topological_charge}

\paragraph{Canon status.}
\textbf{[SPECULATIVE] Research Track.}
This appendix does not promote a new particle assignment to canon. It records a canon-compatible extension of the SST topology sector: knot type and helicity may be insufficient to determine dynamical stability when reconnection and strand-crossing moves are allowed. A non-Abelian coloring invariant may supply an additional obstruction to decay.

\subsubsection{Motivation}
The current SST knot sector uses geometric and topological descriptors such as crossing number, ropelength, helicity, chirality, and hyperbolic-volume proxies. A minimal research-track screening form is
\begin{equation}
    \frac{E_{\rm eff}[K]}{E_0}
    =\alpha_C C(K)+\alpha_L L(K)+\alpha_H\widehat{\mathcal H}(K)+\cdots,
    \label{eq:sst_eff_energy_old}
\end{equation}
where $C(K)$ is a crossing-complexity proxy, $L(K)$ is a ropelength or length proxy, and $\widehat{\mathcal H}(K)$ denotes nondimensional helicity or helicity-derived topological data.  This is a scoring functional, not a derived Hamiltonian.

Recent work on non-Abelian topological defects shows that some linked structures cannot be removed by local reconnections and strand crossings because the strands carry non-commuting elements of an order-parameter fundamental group \cite{Annala2022ProtectedVortexKnots}. This suggests that SST should distinguish two layers:
\begin{enumerate}
    \item the geometric knot type $K$ of the swirl string centerline;
    \item an internal topological coloring $q$ valued in a discrete or continuous group $G$.
\end{enumerate}

A concrete benchmark choice is the quaternion group
\begin{equation}
    Q_8=\{\pm 1,\pm i,\pm j,\pm k\},
    \qquad ij=k,
    \qquad ji=-k,
\end{equation}
so that $ij\neq ji$.  The general symbol $G$ below is retained because SST does not yet select a unique internal order-parameter group.

\subsubsection{Group-colored swirl-string sector}
Let a colored swirl-string configuration be denoted
\begin{equation}
    \mathcal K_G=(K,q),
\end{equation}
where
\begin{equation}
    q: \pi_1\!\left(\mathbb R^3\setminus K\right)\longrightarrow G
    \label{eq:coloring_homomorphism}
\end{equation}
assigns a group element to each Wirtinger generator of the complement. For a diagram with arc generators $a_i$, the coloring must satisfy the usual crossing consistency relations. In SST language, $q(a_i)$ labels an internal director or phase-sector attached to a local swirl-string strand.

\textbf{[ORTHODOX] Constraint.}
For a non-Abelian group $G$, strand exchange is obstructed when the corresponding labels do not commute:
\begin{equation}
    [q(a_i),q(a_j)]
    = q(a_i)q(a_j)q(a_i)^{-1}q(a_j)^{-1}
    \neq e_G.
    \label{eq:noncommuting_obstruction}
\end{equation}

\textbf{[SPECULATIVE] SST interpretation.}
Equation~\eqref{eq:noncommuting_obstruction} is interpreted as a local reconnection barrier in the swirl medium. If two strands belong to non-commuting internal sectors, a direct crossing or reconnection requires an additional defect cord or phase bridge. In SST terms this is a candidate mechanism for topological protection beyond ordinary helicity conservation.

\subsubsection{Protected-sector invariant}
Define a protection marker
\begin{equation}
    \mathcal P_G(\mathcal K_G)=
    \begin{cases}
        1, & \text{if no allowed local reconnection/crossing sequence reduces }\mathcal K_G\text{ to the trivial sector},\\[0.3em]
        0, & \text{otherwise.}
    \end{cases}
    \label{eq:protection_marker}
\end{equation}
The marker $\mathcal P_G$ is not an energy term by itself. It is a sector label. Energetically, it enters as a barrier rule:
\begin{equation}
    \Delta E_{\rm rec}(\mathcal K_G\rightarrow \mathcal K'_G)=
    \begin{cases}
        <\infty, & Q_G(\mathcal K_G)=Q_G(\mathcal K'_G),\\[0.3em]
        \infty \text{ or macroscopically large}, & Q_G(\mathcal K_G)\neq Q_G(\mathcal K'_G),
    \end{cases}
    \label{eq:nonabelian_barrier_rule}
\end{equation}
where $Q_G$ is a conserved coloring invariant under the permitted move set.

A useful research-track extension of Eq.~\eqref{eq:sst_eff_energy_old} is therefore the scalarized screening form
\begin{equation}
    \frac{E_{\rm eff}^{(G)}[\mathcal K_G]}{E_0}
    =\alpha_C C(K)+\alpha_L L(K)+\alpha_H\widehat{\mathcal H}(K)
    +\alpha_G I_G(\mathcal K_G),
    \label{eq:nonabelian_extended_energy}
\end{equation}
where
\begin{equation}
    I_G:Q_G(\mathcal K_G)\longrightarrow \mathbb R_{\ge 0}
\end{equation}
extracts a real-valued index from the conserved coloring sector.  The barrier scale $\Delta_G$ belongs to reconnection dynamics through Eq.~\eqref{eq:nonabelian_barrier_rule}; it is not promoted to a fitted mass term.

\subsubsection{Dimensional check}
The quantities $C(K)$, $L(K)$ after ropelength normalization, $\widehat{\mathcal H}(K)$, and $I_G$ are dimensionless.  Therefore the coefficients in Eq.~\eqref{eq:nonabelian_extended_energy} are dimensionless if the left-hand side is $E_{\rm eff}/E_0$.  A dimensional energy is recovered only after selecting an SST energy scale $E_0$, for example a calibrated tube energy, a core line-tension scale, or a controlled numerical relaxation scale.  The group object $Q_G$ itself is not added to the energy; only a scalar index extracted from it may enter a screening functional.

\subsubsection{Minimal falsifier}
If two SST candidate particles have the same $(C,L,\mathcal H)$ but different observed stability, the model predicts that a missing internal sector label exists:
\begin{equation}
    (C_1,L_1,\mathcal H_1)=(C_2,L_2,\mathcal H_2),
    \qquad
    \tau_1\neq \tau_2
    \quad\Rightarrow\quad
    Q_G(\mathcal K_1)\neq Q_G(\mathcal K_2).
\end{equation}
Failure to find such a sector in simulation or experiment would demote this appendix from Research Track to rejected model branch.

\subsection{Research Appendix: Reconnection Decay as a Stability Falsifier}
\label{app:reconnection_decay_falsifier}

\paragraph{Canon status.}
\textbf{[ORTHODOX] Constraint + [SPECULATIVE] SST interpretation.}
This appendix records a negative constraint on SST. A single-component Gross--Pitaevskii-like medium with permitted reconnections does not automatically protect knotted filaments \cite{Kleckner2016SuperfluidKnotsUntie}. Therefore, SST particle stability requires additional structure beyond generic phase-defect topology.

\subsubsection{Baseline statement}
In a dissipative or reconnection-permitting medium, knotted filament topology can decay. The relevant orthodox lesson is not that SST is false, but that topology alone is insufficient unless the move set is restricted by the dynamics.

Let $K(t)$ denote the topology of a swirl-string centerline. A reconnection path is a sequence
\begin{equation}
    K_0 \longrightarrow K_1 \longrightarrow \cdots \longrightarrow K_N,
    \label{eq:reconnection_path}
\end{equation}
where each arrow is a local reconnection or strand-crossing event. If $K_N$ is an unlink or collection of unknotted loops, then $K_0$ is dynamically unprotected under that move set.

\subsubsection{SST stability criterion}
SST should not claim that every closed filament is stable. A canon-compatible stability criterion is
\begin{equation}
    \tau_{\rm life}(K)
    \sim \tau_0\exp\!\left(\frac{\Delta E_{\rm rec}(K)}{E_{\rm drive}}\right),
    \label{eq:sst_lifetime_barrier}
\end{equation}
where $\Delta E_{\rm rec}(K)$ is the lowest reconnection barrier out of the topological sector, and $E_{\rm drive}$ is the effective environmental or internal driving energy available to trigger reconnection.

The stable-particle condition is therefore not merely
\begin{equation}
    K\neq \bigcirc,
\end{equation}
but rather
\begin{equation}
    \Delta E_{\rm rec}(K)\gg E_{\rm drive}
    \quad\text{and/or}\quad
    Q_G(K)\neq Q_G(\bigcirc).
    \label{eq:sst_stability_condition}
\end{equation}

\subsubsection{Relation to Chronos--Kelvin transport}
For ideal SST evolution without reconnection or external source, the Chronos--Kelvin invariant remains
\begin{equation}
    \frac{D}{Dt}\left(R^2\omega\right)=0.
\end{equation}
Reconnection is precisely a non-ideal event where this transport law is locally invalid or discontinuous. Thus a $\kappa$-type transition in the canon may be interpreted as the event layer at which topology can change.

\textbf{[SPECULATIVE] Identification.}
A reconnection event is an SST Kairos event when it changes the topological sector:
\begin{equation}
    \kappa:\quad K_i\rightarrow K_{i+1},
    \qquad
    Q_G(K_i)\neq Q_G(K_{i+1})
    \quad\text{or}\quad
    \mathcal H(K_i)\neq \mathcal H(K_{i+1}).
\end{equation}

\subsubsection{Dimensionless diagnostic}\subsubsection{Topological Hitting Conditions and Chronos--Kelvin Phase Slips}
\label{subsubsec:chronos_kelvin_hitting_conditions}

\paragraph{Canon status ledger.}
\begin{center}
\small
\begin{tabular}{p{0.18\linewidth}p{0.24\linewidth}p{0.48\linewidth}}
\textbf{[ORTHODOX]} & Geometric thickness & Curvature radius and doubly-critical self-distance control finite tube contact.\\
\textbf{[ORTHODOX]} & Phase-slip/reconnection layer & He-II/GP-like media permit non-ideal reconnection events.\\
\textbf{[CANON-COMPATIBLE]} & Ideal SST transport & Without reconnection/source, Chronos--Kelvin transport remains conserved.\\
\textbf{[RESEARCH TRACK]} & First-hitting boundary & A Kairos event is modeled as a distributional boundary condition, not a smooth correction.\\
\textbf{[SPECULATIVE]} & SST core scale & The identification $a_{\rm rec}=\rc$ requires a resolved-core SST move-set proof.\\
\textbf{[FALSIFIER]} & Two-scale test & He-II coherence length and SST core radius define distinct reconnection thresholds.
\end{tabular}
\end{center}

This block is Research Track. It does not modify the ideal SST evolution law. It defines the non-ideal boundary surface at which the ideal Chronos--Kelvin branch ceases to be a valid smooth continuation and the state must be represented in the Kairos event layer. Orthodox superfluid reconnections and phase-slip phenomenology provide the analogue event class, not the SST core proof \cite{KoplikLevine1993,BewleyPaolettiSreenivasanLathrop2008,Varoquaux2015AndersonPhaseSlippage,Kleckner2016SuperfluidKnotsUntie}.

\paragraph{Orthodox geometric premise.}
In the orthodox ropelength setting, finite tube contact is controlled by curvature and doubly-critical self-distance \cite{CantarellaKusnerSullivan2002,AshtonCantarellaPiatekRawdon2011}. Let $\gamma_K:S^1\times\mathbb R\to\mathbb R^3$ denote the centerline of a finite-thickness vortex tube in knot class $K$, with curvature $\kappa(s,t)$. Its geometric thickness is
\begin{equation}
    \tau_K(t)
    =
    \min\!\left[
        \frac{1}{\kappa_{\max}(t)},
        \frac{1}{2}d_{\rm dcsd}(t)
    \right],
    \label{eq:sst_thickness_dcsd_split}
\end{equation}
where
\begin{equation}
    d_{\rm dcsd}(t)
    =
    \min_{(s_1,s_2)\in {\rm dcsd}}
    \|\gamma_K(s_1,t)-\gamma_K(s_2,t)\|.
    \label{eq:sst_dcsd_definition}
\end{equation}
The min-structure in Eq.~\eqref{eq:sst_thickness_dcsd_split} is the key point: a loss of thickness may occur either by excessive curvature or by self-contact. These are physically distinct events.

\paragraph{Derived split: kink versus reconnection.}
Define $\inf\emptyset:=+\infty$. The first curvature-contact event is
\begin{equation}
    t_{\rm kink}
    =
    \inf\left\{
    t:
    \frac{1}{\kappa_{\max}(t)}\le a_{\rm rec},
    \quad
    D_t^+\!\left(\frac{1}{\kappa_{\max}}\right)<0
    \right\}.
    \label{eq:sst_t_kink}
\end{equation}
The first topological reconnection event is instead
\begin{equation}
    t_{\rm rec}
    =
    \inf\left\{
    t:
    d_{\rm dcsd}(t)\le 2a_{\rm rec},
    \quad
    D_t^+d_{\rm dcsd}(t)<0,
    \quad
    E_{\rm avail}(t)\ge E_{\rm b},
    \quad
    \Delta w\in\mathbb Z\setminus\{0\}
    \right\}.
    \label{eq:sst_t_rec}
\end{equation}
Here $D_t^+$ denotes the right-sided time derivative, and $\Delta w$ denotes the integer winding or phase-slip jump across the local core-contact surface. The first breakdown of ideal smooth evolution is
\begin{equation}
    \boxed{
    t_\ast=\min(t_{\rm kink},t_{\rm rec}).
    }
    \label{eq:sst_first_hitting_time}
\end{equation}
Only $t_{\rm rec}$ is a topology-changing reconnection event. A curvature event $t_{\rm kink}$ is a local core-compression or kink-barrier event and must not be automatically labeled as a topological defect.

\paragraph{Reconnection scale.}
The reconnection radius is medium-dependent:
\begin{equation}
    a_{\rm rec}
    =
    \begin{cases}
        \rc, & \text{SST resolved-core hypothesis},\\
        \xi, & \text{orthodox He-II coherence-length falsifier layer}.
    \end{cases}
    \label{eq:sst_reconnection_scale}
\end{equation}
With $\rc=1.40897017\times10^{-15}\,{\rm m}$ and a representative He-II coherence length $\xi\sim10^{-10}\,{\rm m}$,
\begin{equation}
    \frac{\xi}{\rc}
    \simeq
    7.09738\times10^4,
    \qquad
    d_{\rm rec}^{\rm SST}=2\rc=2.81794034\times10^{-15}\,{\rm m},
    \qquad
    d_{\rm rec}^{\rm He-II}\sim2\times10^{-10}\,{\rm m}.
    \label{eq:sst_heii_core_scale_mismatch}
\end{equation}
Thus He-II reconnection experiments primarily test the allowed topological move-set of a superfluid medium; they do not directly probe the SST core radius unless an additional scale-bridge is derived.

\paragraph{Chronos--Kelvin branch and distributional failure.}
On the ideal branch $t<t_\ast$, define the scalar Chronos--Kelvin quantity
\begin{equation}
    \Lambda_{\rm CK}
    =
    \frac{\vnorm}{r_s}R^2
    \sqrt{1-\left(S_t^{\boldsymbol{\circlearrowleft}}\right)^2}.
    \label{eq:sst_lambda_ck_hitting}
\end{equation}
Then
\begin{equation}
    \frac{D\Lambda_{\rm CK}}{Dt_{ae}}=0,
    \qquad t<t_\ast.
    \label{eq:sst_lambda_ck_conserved_before_hit}
\end{equation}
At a Kairos event the same quantity is not smoothly corrected. It acquires a distributional jump:
\begin{equation}
    \frac{D\Lambda_{\rm CK}}{Dt_{ae}}
    =
    \sum_j \Delta\Lambda_j\,\delta(t-t_j),
    \qquad t_j=t_{{\rm rec},j}.
    \label{eq:sst_lambda_ck_distributional_jump}
\end{equation}
Consequently, the Swirl--Clock branch
\begin{equation}
    S_t^{\boldsymbol{\circlearrowleft}}(R)
    =
    \sqrt{1-
    \left(
        \frac{\Lambda_{\rm CK}r_s}{\vnorm R^2}
    \right)^2}
    \label{eq:sst_swirl_clock_branch_hitting}
\end{equation}
exists only while
\begin{equation}
    R\ge R_{\rm CK}
    =
    \left(\frac{\Lambda_{\rm CK}r_s}{\vnorm}\right)^{1/2}.
    \label{eq:sst_ck_radius_bound}
\end{equation}
The operative transition radius is therefore
\begin{equation}
    \boxed{
    R_{\rm trans}=\max(a_{\rm rec},R_{\rm CK}).
    }
    \label{eq:sst_transition_radius_max}
\end{equation}
Inside $R_{\rm trans}$, the ideal smooth Swirl--Clock continuation is no longer defined; the state must be represented in the non-ideal radiative defect sector.

\paragraph{Energetic jump condition.}
The available reconnection energy must be non-negative. For a proposed move $K_i\to K_{i+1}$ define
\begin{equation}
    E_{\rm avail}^{ij}
    =
    W_{\rm ext}
    +
    \left[ E_{\rm tube}(K_i)-E_{\rm tube}(K_{i+1})\right]_+,
    \qquad
    [x]_+:=\max(x,0).
    \label{eq:sst_positive_part_available_energy}
\end{equation}
The reconnection gate opens only if
\begin{equation}
    \boxed{
    E_{\rm avail}^{ij}\ge E_{\rm b}^{ij},
    }
    \label{eq:sst_reconnection_energy_gate}
\end{equation}
where $E_{\rm b}^{ij}$ is the local core-contact or phase-slip barrier. A useful leading slender-tube scale is
\begin{equation}
    E_{\rm tube}(K)
    \simeq
    \frac{\rhoF\Gamma_0^2}{4\pi}\,
    \ell_K
    \ln\!\left(\frac{R_{\rm out}}{a_{\rm rec}}\right).
    \label{eq:sst_slender_tube_energy_hitting}
\end{equation}
The units are
\begin{equation}
    [\rhoF\Gamma_0^2\ell_K]
    =
    \left(\frac{\rm kg}{{\rm m}^3}\right)
    \left(\frac{{\rm m}^4}{{\rm s}^2}\right)
    ({\rm m})
    =
    {\rm J}.
    \label{eq:sst_hitting_energy_units}
\end{equation}

\paragraph{Canonical transition matrix.}
Let
\begin{equation}
    N_{\rm slip}(t)
    =
    \sum_j n_jH(t-t_j),
    \qquad
    \dot N_{\rm slip}(t)
    =
    \sum_j n_j\delta(t-t_j).
    \label{eq:sst_phase_slip_counter}
\end{equation}
Then the state split is
\begin{equation}
\boxed{
\begin{aligned}
\text{Stable T-phase / no-slip state:}\quad
& \tau_K>a_{\rm rec},
\quad N_{\rm slip}={\rm const.},
\quad E_{\rm avail}<E_{\rm b},
\quad \frac{D\Lambda_{\rm CK}}{Dt_{ae}}=0;\\[0.4em]
\text{Dissipative Kairos event:}\quad
& d_{\rm dcsd}\le2a_{\rm rec},
\quad D_t^+d_{\rm dcsd}<0,
\quad \Delta w\in\mathbb Z\setminus\{0\},
\quad E_{\rm avail}\ge E_{\rm b},\\
& \frac{D\Lambda_{\rm CK}}{Dt_{ae}}
=\sum_j\Delta\Lambda_j\delta(t-t_j)\ne0.
\end{aligned}}
\label{eq:sst_canonical_transition_matrix_hitting}
\end{equation}

\paragraph{Falsification rule.}
A simulation or experiment which labels a curvature-only event $t_{\rm kink}$ as a topological reconnection violates the status separation in Eq.~\eqref{eq:sst_first_hitting_time}. Conversely, a claimed reconnection without $d_{\rm dcsd}\le2a_{\rm rec}$, nonzero phase-slip index, and a positive energy gate violates Eq.~\eqref{eq:sst_t_rec}. This makes the conclusion operational:
\begin{equation}
    \boxed{
    \text{ideal SST conserves }\Lambda_{\rm CK};
    \quad
    \text{Kairos events inject distributional jumps.}
    }
    \label{eq:sst_ck_kairos_conclusion}
\end{equation}

\subsubsection{Dimensionless diagnostic}
Define the stability ratio
\begin{equation}
    \Xi_{\rm rec}(K)
    = \frac{\Delta E_{\rm rec}(K)}{E_{\rm tube}(K)}.
    \label{eq:xi_reconnection}
\end{equation}
For a slender swirl tube, the standard leading envelope energy scale is
\begin{equation}
    E_{\rm tube}(K)
    \simeq
    \frac{\rhoF\Gamma^2}{4\pi}\,\ell_K\ln\!\left(\frac{R}{\rc}\right),
    \label{eq:slender_tube_energy_scale}
\end{equation}
where $\Gamma$ is circulation, $\ell_K$ is centerline length, and $R$ is the outer cutoff. The units check:
\begin{equation}
    [\rhoF\Gamma^2\ell_K]
    = \frac{\mathrm{kg}}{\mathrm m^3}\frac{\mathrm m^4}{\mathrm s^2}\mathrm m
    = \mathrm{kg}\,\mathrm m^2\mathrm s^{-2}=\mathrm J.
\end{equation}

A practical SST classification rule is
\begin{equation}
    \Xi_{\rm rec}(K)\lesssim 1 \Rightarrow \text{hydrodynamically fragile},
    \qquad
    \Xi_{\rm rec}(K)\gg 1 \Rightarrow \text{candidate stable T-phase state}.
\end{equation}

\subsubsection{Minimal simulation protocol}
A candidate SST knot assignment should be tested under three levels of dynamics:
\begin{enumerate}
    \item single-component phase-defect evolution;
    \item constrained incompressible evolution without reconnection;
    \item multi-director or non-Abelian colored evolution with reconnection rules.
\end{enumerate}
Only states stable under the third protocol should be promoted as particle candidates.

\subsection{Research Appendix: Fractional Torus-Phase Winding and Swirl-Clock Phase}
\label{app:fractional_torus_phase}

\paragraph{Canon status.}
\textbf{[SPECULATIVE] Research Track.}
This appendix records a bridge between coordinated rotational phase mechanics and SST R-phase torsional packets, using fractional-order optical torus knots as the mathematical template \cite{Pisanty2018FractionalKnotsLight}. It does not assert that optical polarization knots are literal SST particles. It supplies a useful mathematical template for phase winding, internal rotation, and torus-knot closure.

\subsubsection{Coordinated phase generator}
Let $\theta$ be an external azimuthal coordinate and let $\phi$ be an internal director angle. A coordinated rotation acts as
\begin{equation}
    \theta\mapsto \theta+\alpha,
    \qquad
    \phi\mapsto \phi+\gamma\alpha.
\end{equation}
The generator has the schematic form
\begin{equation}
    J_\gamma = L_{\rm orb}+\gamma S_{\rm int},
    \label{eq:jgamma_sst}
\end{equation}
where $L_{\rm orb}$ shifts external azimuthal phase and $S_{\rm int}$ shifts internal swirl-director phase.

For two commensurate internal frequencies $p\Omega$ and $q\Omega$ carrying azimuthal mode numbers $m_1$ and $m_2$, the coordinated-winding parameter is
\begin{equation}
    \gamma
    =\frac{m_2p-m_1q}{p+q}.
    \label{eq:gamma_fractional_winding}
\end{equation}
The corresponding torus-winding pair is
\begin{equation}
    (m,n)=\left(m_2p-m_1q,\;p+q\right).
    \label{eq:torus_pair_fractional_phase}
\end{equation}
If $\gcd(m,n)=1$, the internal lobe trajectory is a single $(m,n)$ torus knot. If $d=\gcd(m,n)>1$, the trajectory splits into $d$ linked components.

\subsubsection{SST phase-action translation}
To avoid notation conflict with the Swirl-Clock factor $\SwirlClock$, define a phase-action field $\mathscr S(\vec x,t)$. In an R-phase torsional packet, write the reduced phase on a toroidal carrier as
\begin{equation}
    \Theta(\theta,\phi,t)
    = n\phi-m\theta-\Omega t.
    \label{eq:sst_torus_phase}
\end{equation}
The closure condition is
\begin{equation}
    \oint d\Theta = 2\pi N,
    \qquad N\in\mathbb Z.
    \label{eq:phase_closure_condition}
\end{equation}
The phase-action field may be represented locally as
\begin{equation}
    \mathscr S(\vec x,t)=\hbar_{\rm SST}\,\Theta(\vec x,t),
    \label{eq:phase_action_relation}
\end{equation}
so that the circulation quantization condition becomes
\begin{equation}
    \oint \nabla\mathscr S\cdot d\vec\ell
    =2\pi N\hbar_{\rm SST}.
    \label{eq:phase_action_quantization}
\end{equation}

\subsubsection{Connection to Swirl-Clock modulation}
The local Swirl-Clock factor remains the kinematic clock factor
\begin{equation}
    \SwirlClock
    =\sqrt{1-\frac{\lVert\uswirl(\mathbf{x},t)\rVert^2}{c^2}}.
    \label{eq:swirl_clock_factor_appendix}
\end{equation}
The phase-action field controls wave-like periodicity, while $\SwirlClock$ controls local clock rate. The research-track bridge is that a stable R-phase packet may require simultaneous closure of both:
\begin{align}
    \oint \nabla\mathscr S\cdot d\vec\ell &=2\pi N\hbar_{\rm SST},\label{eq:phase_closure}\\
    \frac{D}{Dt}\left(R^2\omega\right)&=0.\label{eq:kelvin_closure}
\end{align}

\textbf{[SPECULATIVE] Interpretation.}
Equation~\eqref{eq:gamma_fractional_winding} suggests a route to fractional internal winding without fractional circulation. The circulation quantum remains integer-valued, but the internal phase orientation can return only after multiple external turns. This may be useful for modeling R-phase torsional photons, spin-orbit-like packet structure, or generation-family phase layering.

\subsubsection{Dimensional check}
The variables $m,n,p,q,N$ and $\gamma$ are dimensionless. The phase $\Theta$ is dimensionless. Therefore $\mathscr S=\hbar_{\rm SST}\Theta$ has units of action, and Eq.~\eqref{eq:phase_action_quantization} has units of action because $\nabla\mathscr S\cdot d\vec\ell$ has units $({\rm J\,s\,m^{-1}}){\rm m}={\rm J\,s}$.

\subsection{Research Appendix: Topological Fluid Mechanics and the SST Energy Functional}
\label{app:tfm_energy_functional}

\paragraph{Canon status.}
\textbf{[ORTHODOX] Mathematical background + [SPECULATIVE] SST synthesis.}
This appendix records the topological-fluid mechanics background behind the SST knot-energy program \cite{Moffatt1969,Ricca1998KnotFluidMechanics}. The orthodox part is that ideal incompressible flow preserves topology and helicity under the usual assumptions. The speculative part is the identification of selected relaxed knot classes with particle sectors.

\subsubsection{Frozen-field topology}
In an ideal incompressible medium, topology is transported by the flow map. For SST this supports the use of closed swirl strings as dynamically persistent structures under non-reconnecting evolution.

The primary invariants are
\begin{align}
    \Gamma &= \oint_C \uswirl\cdot d\vec\ell,\label{eq:gamma_appendix}\\
    \mathcal H &= \int_V \uswirl\cdot\vec\omega\,d^3x,\label{eq:helicity_appendix}\\
    \vec\omega &= \nabla\times\uswirl.\label{eq:vorticity_appendix}
\end{align}
For linked tubes with circulations $\Gamma_i$, the leading topological helicity decomposition is
\begin{equation}
    \mathcal H
    \simeq
    \sum_i \Gamma_i^2\,SL_i
    +2\sum_{i<j}Lk_{ij}\Gamma_i\Gamma_j,
    \label{eq:helicity_link_decomposition}
\end{equation}
where $SL_i$ is self-linking and $Lk_{ij}$ is Gauss linking number.

\subsubsection{Cohomological helicity and mapping-class sectors}
\label{subsubsec:rt_helicity_mapping_class}
\paragraph{Status.}
\textbf{[ORTHODOX DIFFERENTIAL-TOPOLOGICAL INPUT / RESEARCH-TRACK SST STATE LABEL].}
Let $\alpha=\iota_{\mathbf V}d^3x$ be the two-form dual to a divergence-free
field $\mathbf V$ tangent to the boundary of a compact domain
$\Omega\subset\mathbb R^3$.  Cantarella and Parsley express the same helicity
both as a configuration-space cohomology pairing and as a wedge product with a
Biot--Savart-selected primitive \cite{CantarellaParsley2010Helicity}:
\begin{align}
    \mathcal H(\alpha)
    &=
    \int_{C_2[\Omega]}
    \alpha_x\wedge\alpha_y\wedge g^*(d\Omega_{S^2}),
    \label{eq:rt_helicity_configuration_space}\\
    &=
    \int_{\Omega}\alpha\wedge\operatorname{BS}(\alpha).
    \label{eq:rt_helicity_primitive_form}
\end{align}
The compactified configuration space $C_2[\Omega]$ resolves the pair diagonal
$x=y$ by retaining the limiting direction of approach.  This is a geometric
regularization of the Gauss kernel; it is not an SST core cutoff and does not
set $a_{\rm core}$ or $\rc$.

\paragraph{Primitive and gauge discipline.}
On a multiply connected domain, the value of
$\int_{\Omega}\alpha\wedge\beta$ need not be independent of an arbitrary
primitive $d\beta=\alpha$.  The Biot--Savart/Hodge choice in
Eq.~\eqref{eq:rt_helicity_primitive_form}, together with the independent flux
and boundary data, is part of the observable definition.  A numerical export
of $\mathcal H$ is therefore incomplete unless it declares
\begin{align}
    \left(
        \boldsymbol\omega,
        \{\Gamma_i\},
        [\beta|_{\partial\Omega}],
        \text{Hodge/Biot--Savart gauge}
    \right).
    \label{eq:rt_helicity_gauge_metadata}
\end{align}
For the kinetic-helicity mapping one sets $\mathbf V=\boldsymbol\omega$ and
$\operatorname{BS}(\boldsymbol\omega)=\uswirl$ up to the separately fixed
harmonic sector.

\paragraph{Mapping-class jump law.}
For an orientation-preserving diffeomorphism $f:\Omega\to\Omega'$, the exact
change is
\begin{align}
    \Delta\mathcal H
    &=
    \mathcal H\!\left((f^{-1})^*\alpha\right)-\mathcal H(\alpha)
    =
    \sum_{i,j}c_{ij}\Phi_i\Phi_j
    =
    \boldsymbol\Phi^{\!T}C\boldsymbol\Phi,
    \label{eq:rt_helicity_mapping_class_quadratic}
\end{align}
where the coefficients $c_{ij}$ are fixed by the induced action on boundary
homology.  For a solid torus and $j\in\mathbb Z$ Dehn twists,
\begin{align}
    \Delta\mathcal H=j\Phi^2 .
    \label{eq:rt_helicity_dehn_twist_jump}
\end{align}
With $\mathbf V=\boldsymbol\omega$, the spanning-disk flux is the circulation,
$\Phi=\Gamma$.  One canonical circulation unit therefore gives
\begin{align}
    \Gamma_0
    &=2\pi\rc\vchar
      =9.68361920\times10^{-9}\,{\rm m^2\,s^{-1}},\\
    \Delta\mathcal H_{j=1}
    &=\Gamma_0^2
      =9.37724809\times10^{-17}\,{\rm m^4\,s^{-2}}.
    \label{eq:rt_helicity_dehn_twist_numeric}
\end{align}
The dimensions are those of kinetic helicity.  Equation
\eqref{eq:rt_helicity_dehn_twist_numeric} does not quantize circulation by
itself; integer circulation requires an independent phase or Onsager closure.

\paragraph{Relation to writhe and twist.}
In the thin single-tube limit,
\begin{align}
    \mathcal H\simeq\Gamma^2(Wr+Tw),
    \label{eq:rt_helicity_wr_tw_sector}
\end{align}
and a Dehn twist changes the framing by $\Delta Tw=j$ while leaving the
centerline embedding potentially unchanged.  Thus two states with the same
$K$, $\operatorname{Rop}(K)$, $Wr$, and $\Gamma$ may still differ by
$j\Gamma^2$ in helicity.  The SST state label must therefore be enlarged to
\begin{align}
    \mathfrak T_K
    =
    \left(K,\{\Gamma_i\},Wr,Tw,[f]_{\rm map}\right).
    \label{eq:rt_helicity_state_label}
\end{align}
The twist variable is measured state data, not a residual parameter chosen to
close a target helicity.

\paragraph{Chronos--Kairos separation.}
A smooth ideal material flow $f_t$ with $f_0=\operatorname{id}$ is homotopic to
the identity and conserves helicity.  A transition between distinct mapping
classes is therefore entered in the Kairos ledger and must name its physical
trigger.  Fluxless sectors, $\Phi_i=0$, are invariant under the quadratic jump
law even when the mapping class is nontrivial.

\paragraph{Falsifier.}
If a numerical ideal-Euler evolution without boundary work, phase slip, or
reconnection changes the tuple
$\left(\mathcal H,\{\Gamma_i\},[f]_{\rm map}\right)$ beyond discretization
error, the implementation violates the frozen-field topology gate.

\subsubsection{Energy relaxation under topological constraints}
For a slender swirl tube, the leading hydrodynamic energy scale is
\begin{equation}
    E_{\rm env}
    \simeq
    \frac{\rhoF\Gamma^2}{4\pi}\,\ell_K\ln\!\left(\frac{R}{\rc}\right),
    \label{eq:env_energy_tfm}
\end{equation}
where $\ell_K$ is the tube centerline length. The core contribution can be represented by a line-tension term
\begin{equation}
    E_{\rm core}\simeq T_{\rm core}\ell_K.
    \label{eq:core_line_tension}
\end{equation}
Near contacts and curvature corrections may be encoded by
\begin{equation}
    E_{\rm geom}[K]
    =T_{\rm core}\ell_K
    +\frac{\rhoF\Gamma^2}{4\pi}\ell_K\ln\!\left(\frac{R}{\rc}\right)
    +\beta_{\kappa}\int_K \kappa_s^2\,ds
    +\beta_{\rm nc}\,C_{\rm near}(K).
    \label{eq:geometric_energy_functional}
\end{equation}
Here $\kappa_s$ is centerline curvature and $C_{\rm near}$ is a near-contact or crossing-complexity proxy.

\subsubsection{Isoperimetric domain-shape and curl-spectrum gate}
\label{subsubsec:rt_isoperimetric_domain_shape}
\paragraph{Status.}
\textbf{[ORTHODOX OPERATOR/SHAPE CALCULUS / RESEARCH-TRACK SST DOMAIN OPTIMIZATION].}
For a smooth compact domain $\Omega\subset\mathbb R^3$, define
\begin{align}
    K(\Omega)
    &=
    \left\{
        \mathbf V:\nabla\cdot\mathbf V=0,
        \ \mathbf V\cdot\mathbf n=0
    \right\},\\
    \operatorname{BS}_{K}
    &=
    P_{K(\Omega)}\operatorname{BS} .
    \label{eq:rt_isoperimetric_projected_bs}
\end{align}
The modified operator is compact and self-adjoint on the $L^2$ completion of
$K(\Omega)$.  Its largest positive eigenvalue, the maximal helicity-to-$L^2$
ratio, and the smallest positive curl eigenvalue on
$\operatorname{im}\operatorname{BS}_{K}$ satisfy
\begin{align}
    \lambda_{\max}(\Omega)
    &=
    \max_{\mathbf V\in K(\Omega)\setminus\{0\}}
    \frac{\langle\mathbf V,
    \operatorname{BS}_{K}\mathbf V\rangle}
    {\langle\mathbf V,\mathbf V\rangle}
    =
    \frac{1}{\kappa_{1}^{\rm curl}(\Omega)},
    \label{eq:rt_isoperimetric_equivalent_problems}\\
    \nabla\times\mathbf V_1
    &=
    \kappa_{1}^{\rm curl}\mathbf V_1 .
    \label{eq:rt_isoperimetric_beltrami_field}
\end{align}
Thus centerline ropelength and the spectrum of the filled tube domain are
independent geometric data \cite{CantarellaDeTurckGluckTeytel2000Isoperimetric}.

\paragraph{Energy-symbol guard.}
The operator paper uses
\begin{align}
    \mathscr E_V
    =
    \int_{\Omega}|\mathbf V|^2d^3x .
    \label{eq:rt_isoperimetric_l2_energy}
\end{align}
For the hydrodynamic choice $\mathbf V=\boldsymbol\omega$,
$\mathscr E_V$ is enstrophy while
\begin{align}
    E_{\rm kin}
    =
    \frac{\rhoF}{2}\int_{\Omega}|\uswirl|^2d^3x
    \label{eq:rt_isoperimetric_kinetic_energy_guard}
\end{align}
is kinetic energy.  Maximizing helicity per enstrophy is not automatically the
same problem as minimizing SST rest energy at fixed helicity.  Any use of
$\kappa_1^{\rm curl}$ in a mass functional must state which field occupies the
operator slot and which physical energy is being extremized.

\paragraph{First variation under a volume-preserving shape deformation.}
Let $\mathbf W$ be the initial velocity of a volume-preserving deformation of
$\Omega$ and carry $\mathbf V$ with the domain.  The orthodox shape derivative
is
\begin{align}
    \delta\mathscr E_V
    =
    -2\int_{\Omega}
    \left[
        \mathbf V\times(\nabla\times\mathbf V)
    \right]\cdot\mathbf W\,d^3x
    -
    \int_{\partial\Omega}
    |\mathbf V|^2\,\mathbf W\cdot\mathbf n\,dA .
    \label{eq:rt_isoperimetric_energy_shape_derivative}
\end{align}
For a curl eigenfield the volume term vanishes, leaving a boundary-only shape
sensitivity.  For a simple eigenvalue the corresponding exact Rayleigh-quotient
derivative is
\begin{align}
    \delta\lambda_1
    =
    \lambda_1
    \frac{
        \int_{\partial\Omega}
        |\mathbf V_1|^2\,\mathbf W\cdot\mathbf n\,dA
    }{
        \int_{\Omega}|\mathbf V_1|^2\,d^3x
    } .
    \label{eq:rt_isoperimetric_eigenvalue_shape_derivative}
\end{align}

\paragraph{Necessary boundary conditions for a smooth global optimum.}
At fixed domain volume, any smooth global maximizer must satisfy
\begin{align}
    \left.|\mathbf V_1|\right|_{\partial\Omega_a}
    &=C_a>0
    \quad\text{on each boundary component},
    \label{eq:rt_isoperimetric_boundary_magnitude}\\
    \chi(\partial\Omega_a)&=0,
    \label{eq:rt_isoperimetric_boundary_euler_characteristic}\\
    k_g(\text{boundary orbit})&=0 .
    \label{eq:rt_isoperimetric_boundary_geodesic}
\end{align}
Hence each smooth closed boundary component is a torus and the boundary field
orbits are intrinsic geodesics.  No smooth simply connected domain can be a
global optimizer in this unrestricted fixed-volume problem.

\paragraph{Horn-envelope non-claim.}
The proposed singular ``extreme apple'' or pinched-torus optimizer in the
source is conjectural.  It does not identify a trefoil, does not determine
$R_{\rm horn}$, and does not resolve a finite SST vortex core.  A canon-compatible
research problem is instead the regularized, embedding-restricted program
\begin{align}
    \min_{\gamma_K,\Omega_K}
    \quad
    \kappa_{1}^{\rm curl}(\Omega_K)
    \quad\text{subject to}\quad
    &[\gamma_K]=K,
    \qquad
    \Omega_K=\operatorname{Tube}_{a_{\rm core}}(\gamma_K),\\
    &\operatorname{Vol}(\Omega_K)=V_0,
    \qquad
    \operatorname{reach}(\gamma_K)\ge a_{\rm core}.
    \label{eq:rt_isoperimetric_regularized_sst_target}
\end{align}
The knot/link embedding, flux data, and harmonic circulations are held fixed
within each comparison.  This target is a Research-Track shape optimization,
not an existing SST theorem.

\paragraph{Operational diagnostics.}
A domain-shape solver should report
\begin{align}
    \mathcal D_{\rm iso}
    =
    \left(
        \kappa_{1}^{\rm curl},
        \lambda_{\max},
        \frac{\operatorname{std}_{\partial\Omega}|\mathbf V_1|}
             {\operatorname{mean}_{\partial\Omega}|\mathbf V_1|},
        \|k_g\|_{L^2(\partial\Omega)},
        \operatorname{reach}(\partial\Omega)
    \right).
    \label{eq:rt_isoperimetric_diagnostics}
\end{align}
The extremal eigenfield is a variational benchmark or stationary candidate.
In an ideal inviscid SST evolution, conservation of energy and helicity does
not make it a dynamical attractor without a separately declared relaxation
mechanism.

\subsubsection{Reduction to the SST normal form}
After nondimensionalization by a reference energy $E_0$, Eq.~\eqref{eq:geometric_energy_functional} reduces to the canon-style scoring form
\begin{equation}
    \frac{E_{\rm geom}[K]}{E_0}
    \leadsto
    \alpha_C C(K)+\alpha_L L(K)+\alpha_H\widehat{\mathcal H}(K)+\alpha_G I_G(K)+\cdots.
    \label{eq:tfm_to_sst_eff}
\end{equation}
Here $I_G(K)$ is included only when a non-Abelian coloring sector $Q_G(K)$ has been specified and scalarized by a map $I_G:Q_G\to\mathbb R_{\ge0}$. This justifies the SST normal form as a coarse-grained energy-ordering tool, not as a finished mass theorem.

\subsubsection{Canon discipline}
The functional in Eq.~\eqref{eq:tfm_to_sst_eff} should be used as follows:
\begin{enumerate}
    \item \textbf{[ORTHODOX]} Use $\Gamma$, $\mathcal H$, $Lk$, writhe, twist, and ropelength as established geometric/topological descriptors.
    \item \textbf{[DERIVED]} Within a chosen tube model, derive the units and leading energy scaling from $\rhoF$, $\Gamma$, $\ell_K$, and $\rc$.
    \item \textbf{[SPECULATIVE]} Only after this, map a relaxed topological class to an SST particle sector.
\end{enumerate}

\subsection{Research Appendix: Resolved-Tube Criticality and Contact-Stress Geometry}
\label{app:resolved_tube_contact_stress_geometry}

\paragraph{Canon status.}
\textbf{[ORTHODOX GEOMETRIC BACKGROUND / RESEARCH-TRACK SST INTERPRETATION].}
The orthodox input is the ropelength theory of finite-thickness embedded
curves and the constrained-gradient formulation of polygonal knot tightening.
The SST interpretation is that the active finite-thickness constraints of an
ideal swirl-string may be read as a discrete contact-stress skeleton for the
resolved tube.  This appendix does not canonize a new mass formula; it
canonizes a validation discipline for any future profile-resolved tube model.

\subsubsection{Resolved tube radius as reach}
For a resolved centerline $\gamma_K$, the physical tube radius is not the horn
radius $\rc$.  The resolved geometric radius is the reach/thickness of the
centerline,
\begin{align}
    a_{\rm core}(K)
    \equiv
    \operatorname{Thi}_{\rm rad}(\gamma_K)
    =
    \operatorname{reach}(\gamma_K)
    =
    \operatorname{NIR}(\gamma_K),
    \label{eq:rt_a_core_reach}
\end{align}
with radius-convention ropelength
\begin{align}
    \operatorname{Rop}_{\rm rad}(K)
    =
    \frac{\operatorname{Len}(\gamma_K)}{a_{\rm core}(K)}.
    \label{eq:rt_rop_radius_convention}
\end{align}
For smooth centerlines the controlling geometric obstructions are curvature
and doubly-critical self-distance,
\begin{align}
    a_{\rm core}(K)
    =
    \min\!\left(
        \inf_s\frac{1}{\kappa(s)},\,
        \frac12 d_{\rm dcsd}(\gamma_K)
    \right).
    \label{eq:rt_thickness_curvature_dcsd}
\end{align}
Equations~\eqref{eq:rt_a_core_reach}--\eqref{eq:rt_thickness_curvature_dcsd}
are \textbf{[ORTHODOX]} geometric definitions.  They preserve the canon rule
$a_{\rm core}\neq\rc$ unless a separate profile calculation proves otherwise.

\subsubsection{Struts, kinks, and contact-stress balance}
In a polygonal approximation $V$ of the centerline, finite thickness is
controlled by two active constraint classes:
\begin{align}
    \operatorname{MinRad}(v_i)&\ge a_{\rm core},
    \label{eq:rt_kink_constraint}\\
    \frac12 d(p,q)&\ge a_{\rm core}.
    \label{eq:rt_strut_constraint}
\end{align}
Vertices saturating Eq.~\eqref{eq:rt_kink_constraint} form the kink set
$\operatorname{Kink}(V)$; pairs saturating Eq.~\eqref{eq:rt_strut_constraint}
form the strut set $\operatorname{Strut}(V)$.  The constrained-gradient
criterion for a thickness-critical state has the schematic KKT form
\begin{align}
    -\nabla E[V]
    +
    \sum_{(p,q)\in\operatorname{Strut}(V)}
        \lambda_{pq}\,
        \nabla\!\left(\frac{d(p,q)}{2}\right)
    +
    \sum_{i\in\operatorname{Kink}(V)}
        \mu_i\,\nabla\operatorname{MinRad}(v_i)
    =0,
    \qquad
    \lambda_{pq},\mu_i\ge0 .
    \label{eq:rt_contact_kkt_balance}
\end{align}
For the mathematical ropelength problem, $E[V]=\operatorname{Len}(V)$.  For a
profile-resolved SST tube one should instead test a functional of the form
\begin{align}
    E_{\rm SST}[\gamma]
    =
    T_{\rm eff}\operatorname{Len}(\gamma)
    +E_{\rm bend}[\gamma]
    +E_{\rm twist}[\gamma]
    +E_{\rm swirl}[\gamma]
    +E_{\rm core}[\gamma] .
    \label{eq:rt_resolved_tube_energy}
\end{align}
The multipliers $\lambda_{pq}$ and $\mu_i$ are then research-track candidates
for a discrete swirl-stress/contact-force map, not yet canonical physical
forces.

\subsubsection{Contact-pressure saturation and Swirl-Clock shielding}
\label{subsec:rt_contact_pressure_saturation_swirl_clock}

\paragraph{Status.}
\textbf{[ORTHODOX INPUT / CANON SEMANTIC GUARD / CONDITIONAL DERIVED / RESEARCH-TRACK CONTACT PHYSICS].}
The orthodox input is the finite-thickness constraint structure of ropelength
criticality: active self-distance constraints are struts, active curvature
constraints are kinks, and their critical balance is expressed by non-negative
KKT multipliers.  The canon semantic guard is that the Swirl-Clock is a
kinematic factor determined by the resolved local swirl speed, not by an
unrestricted contact multiplier.  The algebraic saturation result below is
\textbf{[CONDITIONAL DERIVED]} once the declared Bernoulli/Biot--Savart pressure
closure and speed cap are imposed.  The interpretation of the remaining hard
multiplier as a physical contact-stress pressure is \textbf{[RESEARCH TRACK]}.

\paragraph{Pressure split.}
At an active self-contact patch, decompose the normal contact pressure into a
kinematic Bernoulli/Biot--Savart part and a hard geometric constraint part:
\begin{align}
    \Pi_c
    &=
    \Pi_{\rm BS}
    +
    \Pi_{\rm hard},
    \label{eq:rt_contact_pressure_split}
    \\
    \Pi_{\rm BS}
    &=
    \frac12\rho_\star u_\theta^2 .
    \label{eq:rt_contact_pi_bs}
\end{align}
Here \(\rho_\star\) is the declared load-bearing density of the modeled layer.
The term \(\Pi_{\rm hard}\) is the continuum counterpart of the active
strut/kink multiplier sector enforcing the finite-thickness exclusion
\begin{equation}
    d(s,t)\ge 2a_{\rm core}.
    \label{eq:rt_contact_distance_exclusion}
\end{equation}
Thus \(\Pi_{\rm BS}\) carries the kinematic swirl load, while
\(\Pi_{\rm hard}\) carries any additional geometric incompressibility required
by contact.

\paragraph{Swirl-Clock coupling.}
The local Swirl-Clock remains the Lorentz-compatible kinematic factor
\begin{equation}
    \SwirlClock(\mathbf x,t)
    =
    \sqrt{
        1-
        \frac{\lVert\uswirl(\mathbf x,t)\rVert^2}{c^2}
    } .
    \label{eq:rt_contact_swirl_clock_kinematic}
\end{equation}
For a local azimuthal contact model with speed \(u_\theta\),
Eq.~\eqref{eq:rt_contact_pi_bs} gives
\begin{equation}
    \SwirlClock(\Pi_{\rm BS})
    =
    \sqrt{
        1-
        \frac{2\Pi_{\rm BS}}{\rho_\star c^2}
    } .
    \label{eq:rt_contact_swirl_clock_pressure}
\end{equation}
Consequently,
\begin{equation}
    \nabla\SwirlClock
    =
    -\frac{1}{\rho_\star c^2\SwirlClock}\,
    \nabla\Pi_{\rm BS}
    \label{eq:rt_contact_swirl_clock_gradient_pressure}
\end{equation}
inside the declared Bernoulli/Biot--Savart layer.  No corresponding rule is
assigned to \(\Pi_{\rm hard}\); a hard-contact multiplier may diverge without
forcing \(\SwirlClock\to0\).

\paragraph{Canonical speed cap.}
For one canonical circulation unit,
\begin{equation}
    \Gamma_0=2\pi\rc\vchar .
    \label{eq:rt_contact_gamma0_cap}
\end{equation}
The canonical contact-saturation discipline imposes
\begin{align}
    u_\theta &\le \vchar,
    \label{eq:rt_contact_speed_cap}
    \\
    \Pi_{\rm BS}
    &\le
    \Pi_{\rm cap}
    =
    \frac12\rho_\star\vchar{}^2 .
    \label{eq:rt_contact_pressure_cap}
\end{align}
The Swirl-Clock therefore has the finite lower bound
\begin{equation}
    \boxed{
    \SwirlClock
    \ge
    \SwirlClock_{\min}^{(c)}
    =
    \sqrt{1-\frac{\vchar{}^2}{c^2}}
    }
    \label{eq:rt_contact_swirl_clock_min}
\end{equation}
rather than collapsing to zero at contact.  With
\(\vchar=1.09384563\times10^6\,{\rm m\,s^{-1}}\),
\begin{align}
    \SwirlClock_{\min}^{(c)}
    &=
    0.999993343558553,
    \label{eq:rt_contact_swirl_clock_min_numeric}
    \\
    1-\SwirlClock_{\min}^{(c)}
    &=
    6.65644145\times10^{-6} .
    \label{eq:rt_contact_swirl_clock_deficit_numeric}
\end{align}
If the core-density layer is used, the corresponding kinematic pressure cap is
\begin{equation}
    \Pi_{\rm cap}(\rhocore)
    =
    \frac12\rhocore\vchar{}^2
    =
    2.32924460\times10^{30}\,{\rm Pa} .
    \label{eq:rt_contact_core_pressure_cap_numeric}
\end{equation}

\paragraph{Clock-shear shell.}
An active strut can still generate a steep transition layer in the clock field.
For a characteristic transition thickness \(\ell\),
\begin{equation}
    \lVert\nabla\SwirlClock\rVert
    \sim
    \frac{1-\SwirlClock_{\min}^{(c)}}{\ell} .
    \label{eq:rt_contact_clock_shear_shell}
\end{equation}
At the limiting microscopic scale \(\ell\sim\rc\),
\begin{equation}
    \lVert\nabla\SwirlClock\rVert
    \sim
    4.724\times10^9\,{\rm m^{-1}} .
    \label{eq:rt_contact_clock_gradient_numeric}
\end{equation}
This large gradient is a finite clock-shear shell, not an infinite
time-dilation singularity.  Its integrated clock deficit across the shell is
only of order \(6.66\times10^{-6}\).

\paragraph{Contact-Constraint Theorem.}
In any profile-resolved SST tube model that obeys the canonical speed cap,
\begin{equation}
    \boxed{
    \Pi_c\to\infty
    \quad\Longrightarrow\quad
    \SwirlClock
    \to
    \SwirlClock_{\min}^{(c)}
    =
    \sqrt{1-\vchar{}^2/c^2}
    }
    \label{eq:rt_contact_constraint_theorem}
\end{equation}
provided the divergent part of \(\Pi_c\) is assigned to \(\Pi_{\rm hard}\).
Equivalently, the singularity is carried by the contact-stress sector
\begin{equation}
    \nabla\cdot\mathbf T_{\rm contact},
    \label{eq:rt_contact_stress_divergence_guard}
\end{equation}
not by an unbounded collapse of the Swirl-Clock.  The tube becomes
geometrically hard before the local clock becomes infinitely slow.

\paragraph{Guard against denominator drift.}
A factor of the form \(\sqrt{1-u_\theta^2/\vchar{}^2}\) may be used as a
local topological saturation or transport-efficiency factor inside a declared
model, but it is not the Lorentz-compatible Swirl-Clock.  The operational
clock factor in this canon retains the denominator \(c^2\), as in
Eq.~\eqref{eq:rt_contact_swirl_clock_kinematic}.

\paragraph{Falsification protocol.}
A coupled knot-tightening/Biot--Savart audit must report, for each active strut,
\begin{equation}
    \Pi_{\rm BS},\qquad
    \Pi_{\rm hard},\qquad
    \lambda_{pq},\qquad
    \SwirlClock .
    \label{eq:rt_contact_saturation_audit_tuple}
\end{equation}
The saturation guard fails if \(\Pi_{\rm BS}\) reproducibly exceeds
\(\Pi_{\rm cap}=\tfrac12\rho_\star\vchar{}^2\) without activating a compensating
contact-stress multiplier or without reducing the local speed field back below
\(\vchar\).  A standalone ropelength minimizer does not compute
\(\Pi_{\rm BS}\) by itself; it supplies the contact skeleton that must be coupled
to a Biot--Savart pressure postprocessor.


\subsubsection{Contact-map refinement of the topological kernel}
\label{subsubsec:rt_contact_map_refinement}
A scalar ropelength is not enough to determine all shape-sensitive SST
observables.  A resolved ideal tube carries a contact skeleton
\begin{align}
    \mathcal C_K
    =
    \bigl\{(p,q,\lambda_{pq})\bigr\}_{\operatorname{Strut}}
    \cup
    \bigl\{(i,\mu_i)\bigr\}_{\operatorname{Kink}} .
    \label{eq:rt_contact_skeleton}
\end{align}
A future profile-resolved version of the dimensionless topology kernel may
therefore have the schematic expansion
\begin{align}
    \Xi_K
    =
    \Xi_K^{(0)}
    +\epsilon_s\sum_{(p,q)}\lambda_{pq}
    +\epsilon_k\sum_i\mu_i
    +\epsilon_H\mathcal H_K
    +\cdots .
    \label{eq:rt_kernel_contact_refinement}
\end{align}
Equation~\eqref{eq:rt_kernel_contact_refinement} is a
\textbf{[MATCHING ANSATZ]}.  Its value is diagnostic: if the large calibrated
normalization hidden in $\Xi_K$ can be decomposed into reproducible contact,
curvature, and helicity contributions, the mass kernel becomes mechanically
inspectable rather than a single opaque coefficient.

\subsubsection{Golden-layer factor as hyperbolic transfer parametrization}
\paragraph{Status.}
\textbf{[RESEARCH-TRACK / PARAMETRIC REWRITE].}
This subsection records an exact algebraic reparametrization of the
existing golden-layer factor in the SST kernel.  It does not derive the
factor from Euler dynamics, ropelength criticality, or knot stability.

The identity
\begin{align}
    \operatorname{asinh}\!\left(\frac12\right)
    &=
    \ln\!\left(\frac12+\sqrt{1+\frac14}\right)
     =
    \ln\!\left(\frac{1+\sqrt5}{2}\right)
     =
    \ln\varphi
\end{align}
allows the layer term to be written as
\begin{align}
    \eta_\varphi
    &\equiv
    \operatorname{asinh}\!\left(\frac12\right)
    =
    \ln\varphi
    \simeq
    0.4812118251,
    \label{eq:rt_golden_layer_eta_def}\\
    \varphi^{-2k(K)}
    &=
    \exp\!\left[-2k(K)\eta_\varphi\right].
    \label{eq:rt_golden_layer_exp_rewrite}
\end{align}
Equivalently,
\begin{align}
    \sinh\eta_\varphi=\frac12,
    \qquad
    \cosh\eta_\varphi=\frac{\sqrt5}{2},
    \qquad
    \tanh\eta_\varphi=\frac{1}{\sqrt5}.
    \label{eq:rt_golden_hyperbolic_components}
\end{align}
The associated symmetric transfer matrix
\begin{align}
    B_\varphi
    =
    \begin{pmatrix}
        \sqrt5/2 & 1/2\\
        1/2 & \sqrt5/2
    \end{pmatrix},
    \qquad
    \operatorname{spec}(B_\varphi)=\{\varphi,\varphi^{-1}\},
    \label{eq:rt_golden_transfer_matrix}
\end{align}
therefore supplies a compact notation for one elementary layer-transfer
step.  Iteration gives attenuation eigenvalue
\begin{align}
    \lambda_-^{2k(K)}=\varphi^{-2k(K)}.
\end{align}

\paragraph{Canon discipline.}
This is not a hyperbolic-knot claim.  In particular, the electron-sector
assignment to the trefoil is not made hyperbolic by
Eq.~\eqref{eq:rt_golden_layer_exp_rewrite}; the trefoil remains a torus
knot.  The useful content is only that the existing golden dressing factor
can be written as a discrete exponential attenuation.  To promote this from
notation to a derived SST mechanism, a future calculation must identify
\(k(K)\) with a definite filament-transfer, impedance, phase-locking, or
relaxation step and derive the elementary rapidity
\(\eta_\varphi=\operatorname{asinh}(1/2)\) from that model.

\paragraph{Dimensional check.}
The quantities \(\varphi\), \(k(K)\), \(\eta_\varphi\), and all entries of
\(B_\varphi\) are dimensionless.  Therefore
Eq.~\eqref{eq:rt_golden_layer_exp_rewrite} can modify only a dimensionless
kernel or impedance factor; it cannot by itself introduce a mass, length,
energy, or force scale.

\subsubsection{Geometric lower-bound gates for SSTcore}
For any nontrivial knot in the radius-thickness convention, orthodox
ropelength bounds give
\begin{align}
    \operatorname{Rop}_{\rm rad}(K_{\rm nontrivial})
    \ge
    4\pi+2\pi\sqrt2
    \simeq 21.45213649 .
    \label{eq:rt_nontrivial_rop_lower_bound}
\end{align}
Equivalently,
\begin{align}
    \operatorname{Rop}_{\rm diam}(K_{\rm nontrivial})
    \ge
    2\pi+\pi\sqrt2
    \simeq 10.72606825 .
    \label{eq:rt_nontrivial_rop_diam_bound}
\end{align}
Using the canonical radius $\rc=1.40897017\times10^{-15}\,\mathrm{m}$, the
radius-convention bound corresponds to
\begin{align}
    L_{\min}
    =
    \rc\left(4\pi+2\pi\sqrt2\right)
    \simeq
    3.02254\times10^{-14}\,\mathrm{m} .
    \label{eq:rt_nontrivial_physical_length_bound}
\end{align}
For the trefoil benchmark used elsewhere in the canon,
\begin{align}
    \operatorname{Rop}_{\rm rad}(3_1)
    &\approx 2\times16.371637=32.743274,\\
    L_{3_1}
    &\approx
    32.743274\,\rc
    \simeq
    4.61343\times10^{-14}\,\mathrm{m} .
    \label{eq:rt_trefoil_physical_length_benchmark}
\end{align}
An SSTcore ideal-knot routine that reports a nontrivial knot below
Eq.~\eqref{eq:rt_nontrivial_rop_lower_bound}, or mixes the radius and diameter
conventions without an explicit factor of two, fails the geometry gate.

\subsubsection{Composite-link non-additivity warning}
For composite links and connect-sum-like configurations, ideal ropelength is
not generally additive.  The ropelength literature reports a typical deficit
of order
\begin{align}
    \Delta\mathcal L_{\rm comp}
    \sim
    4\pi-4
    \simeq 8.56637 .
    \label{eq:rt_composite_link_deficit}
\end{align}
Therefore any SST baryon or multi-component particle branch using an additive
length proxy must expose a correction term,
\begin{align}
    \mathcal L_{\rm comp}
    =
    \sum_i\mathcal L_{K_i}
    -\Delta_{\rm link}
    +\delta_{\rm stress},
    \label{eq:rt_composite_nonadditive_length}
\end{align}
rather than assuming $\mathcal L_{\rm comp}=\sum_i\mathcal L_{K_i}$.  This is a
research-track guard, not a universal law for all SST composites.


\subsubsection{Polygonal-to-smooth certification ladder}
\label{subsubsec:rt_polygonal_smooth_certification}
\paragraph{Status.}
\textbf{[ORTHODOX NUMERICAL GEOMETRY / RESEARCH-TRACK SST INPUT GATE].}
This subsection isolates the numerical implications of the Ridgerunner
constrained-gradient method \cite{AshtonCantarellaPiatekRawdon2011}.  It
certifies an approximately thickness-critical polygon and a smooth ropelength
upper bound; it does not certify a unique global minimizer or a physical SST
energy ground state.

\paragraph{Equilateral constraint gate.}
For an equilateral polygon of edge length $\ell$, Ridgerunner's constraint
thickness and Rawdon polygonal thickness define the same feasible set,
\begin{align}
    \operatorname{CThi}(\ell,V)\ge1
    \quad\Longleftrightarrow\quad
    \operatorname{Thi}_p(V)\ge1 .
    \label{eq:rt_cthi_thip_equivalence}
\end{align}
This equivalence must not be transferred silently to a strongly
non-equilateral simulation mesh.  Define
\begin{align}
    \bar\ell
    &=\frac1N\sum_{i=1}^{N}\ell_i,\\
    \epsilon_{\rm eq}
    &=
    \max_i\left|\frac{\ell_i}{\bar\ell}-1\right| .
    \label{eq:rt_equilateral_error}
\end{align}
A pipeline must control $\epsilon_{\rm eq}$ by arc-length remeshing or use a
thickness formulation proved for its non-equilateral discretization.

\paragraph{Curvature-stiffness scaling.}
At normalized $\operatorname{MinRad}=1$, an equilateral polygon obeys
\begin{align}
    \left\|\nabla\operatorname{MinRad}^{\pm}\right\|_2
    \ge
    \frac{2}{\ell^2} .
    \label{eq:rt_minrad_stiffness_scaling}
\end{align}
Thus halving the edge length can increase the curvature-constraint stiffness
by a factor of four.  Timestep, conditioning, and constraint-correction
tolerances must scale with resolution; increasing the vertex count is not by
itself a stability improvement.

\paragraph{Polygonal versus smooth ropelength.}
The polygonal value
\begin{align}
    \operatorname{Rop}_p(V)
    =
    \frac{\operatorname{Len}(V)}{\operatorname{Thi}_p(V)}
    \label{eq:rt_polygonal_ropelength}
\end{align}
is not the final smooth certification.  After replacing vertices by controlled
circular or spline arcs, all newly introduced self-distances must be checked
and one must report
\begin{align}
    \operatorname{Rop}_{\rm smooth}^{\rm ub}
    =
    \frac{\operatorname{Len}(\widetilde V)}
         {\operatorname{Thi}(\widetilde V)} .
    \label{eq:rt_smooth_ropelength_upper_bound}
\end{align}
The label ``ideal'' attached to a Fourier, KnotPlot, or atlas file is not a
geometric certificate.  The Ridgerunner authors could not reproduce every
published Fourier/Gilbert ropelength claim under their polygonal thickness
measurement, which makes independent reconstruction and smooth-thickness
verification mandatory.

\paragraph{Constrained-gradient certificate.}
With rigidity matrix $A$ assembled from active strut and kink gradients, the
nonnegative least-squares problem is
\begin{align}
    \min_{\Lambda\ge0}
    \left\|A\Lambda-(-\nabla E_{\rm SST})\right\|_2,
    \label{eq:rt_ridgerunner_nnls}
\end{align}
and the relative first-order residual is
\begin{align}
    \chi_{\rm KKT}
    =
    \frac{\|-\nabla E_{\rm SST}+A\Lambda\|_2}
         {\|\nabla E_{\rm SST}\|_2},
    \qquad \Lambda\ge0 .
    \label{eq:rt_contact_residual}
\end{align}
The published dataset had mean residual approximately $2.99\times10^{-3}$,
with specially refined examples reaching approximately
$2.54\times10^{-5}$.  These values are historical benchmarks, not universal
acceptance thresholds.  A small $\chi_{\rm KKT}$ establishes only first-order
criticality and must be interpreted together with the stability hierarchy in
Eq.~\eqref{eq:rt_gehring_stability_hierarchy}.

\paragraph{Contact-map and refinement convergence.}
Ropelength convergence is insufficient.  The positive strut measure of
Eq.~\eqref{eq:rt_discrete_contact_measure}, the analogous kink measure, their
supports, total masses, and representative test integrals must converge under
refinement.  The arclength contact map should be compared across resolutions;
new or disappearing contact branches are a geometry change even when the
scalar ropelength changes little.

\paragraph{Multistart protocol.}
For each knot/link class use at least five independent initial embeddings, as
in the published computation, and a resolution ladder.  Suitable starts
include analytic/Fourier, diagram/KnotPlot, previously tightened atlas, and
independent isotopic perturbations.  Retain the best certified smooth upper
bound, not merely the shortest polygon encountered, and preserve seed,
remeshing, active-set, correction-step, and stopping-residual metadata.

\paragraph{Precision ceiling.}
Define a conservative geometry-certification envelope
\begin{align}
    \epsilon_{\rm geom}
    =
    \max\!\left(
        \epsilon_{\rm eq},
        \epsilon_{\rm thi},
        \chi_{\rm KKT},
        \epsilon_{\rm smooth},
        \epsilon_{\rm contact}
    \right),
    \label{eq:rt_geometry_precision_ceiling}
\end{align}
where, for the normalized unit-thickness problem,
\begin{align}
    \epsilon_{\rm thi}
    :=\max\!\left(0,1-\operatorname{CThi}(\bar\ell,V)\right),\\
    \epsilon_{\rm smooth}
    :=
    \frac{\left|\operatorname{Rop}_{\rm smooth}^{\rm ub}
    -\operatorname{Rop}_p\right|}
    {\operatorname{Rop}_{\rm smooth}^{\rm ub}} .
    \label{eq:rt_geometry_error_components}
\end{align}
Here $\epsilon_{\rm contact}$ is a declared resolution-to-resolution norm on
the contact measures.  A downstream SST mass, coupling, or constant extraction
must not quote relative precision parametrically smaller than
$\epsilon_{\rm geom}$ unless an independent error-propagation analysis proves
a cancellation.

\subsubsection{Operational SSTcore diagnostics}
A profile-resolved SSTcore tightening or imported ideal-knot pipeline should
report at least
\begin{align}
    \mathcal D_K^{\rm num}
    =
    \Bigg(
        N,
        \bar\ell,
        \epsilon_{\rm eq},
        \operatorname{Thi}_p,
        \operatorname{CThi},
        \operatorname{Rop}_p,
        \operatorname{Rop}_{\rm smooth}^{\rm ub},
        \min_i\operatorname{MinRad}(v_i),
        \frac12d_{\rm dcsd},
        |\operatorname{Strut}|,
        |\operatorname{Kink}|,
        \Lambda,
        \chi_{\rm KKT},
        \operatorname{cond}_2(A_F),
        N_{\rm seed},
        \epsilon_{\rm geom}
    \Bigg),
    \label{eq:rt_sstcore_contact_diagnostics}
\end{align}
where the rigidity-matrix column-sign convention is declared and $A_F$ is the
free-column submatrix in the final NNLS active set.  This tuple supplements,
rather than replaces, the continuum contact-measure audit of
Sec.~\ref{subsec:rt_gehring_contact_measure}.

\subsubsection{Colored contact-channel rank test for the factor nine}
\label{sec:rt_colored_contact_rank_nine}

\paragraph{Status.}
\textbf{[ORTHODOX OPTIMIZATION BACKGROUND / RESEARCH-TRACK SST--QCD ROSETTA / NOT CANONICAL QCD].}
The orthodox input is the finite-thickness ropelength/KKT framework summarized in
Eqs.~\eqref{eq:rt_thickness_curvature_dcsd}--\eqref{eq:rt_contact_kkt_balance}: active
struts and kinks define a contact/rigidity matrix with nonnegative multipliers.
The SST interpretation below is a confinement-sheet diagnostic for a
\(\mathbb Z_3\)-locked colored trefoil triad.  It is not a derivation of the QCD
beta function, the full \(SU(3)\) Yang--Mills action, running couplings, anomaly
cancellation, or the observed hadron spectrum.

\paragraph{Claim discipline.}
A QCD-like numerical factor of nine must not be introduced as a fitted prefactor
or as the statement ``eight gluons plus one''.  In this research-track bridge it
is instead the target positive rank of a resolved colored contact-stress matrix,
\begin{align}
    9
    &=
    3_{\rm color}\times 3_{\rm mech},
    \label{eq:rt_rank9_factorization}\\
    3_{\rm color}
    &:=
    \#\{RG,GB,BR\},\\
    3_{\rm mech}
    &:=
    \#\{A,S,T\}.
\end{align}
Here \(RG,GB,BR\) denote the three \(\mathbb Z_3\) color root-planes of the
ideal triad.  The mechanical labels denote, respectively,
\begin{align}
    A&: \text{antisymmetric-vector / torsional contact mode},\\
    S&: \text{symmetric-traceless in-sheet shear mode},\\
    T&: \text{transverse director-compression mode}.
\end{align}
Only the color triad is a discrete sector label.  The three mechanical modes are
finite-thickness contact modes and must be verified by the solver.

\paragraph{Resolved contact decomposition.}
Let \(\Sigma_Y\) be the putative confinement sheet joining a separated trefoil
sector to the Steiner node of a \(\mathbb Z_3\)-locked triad.  The active strut set
is decomposed as
\begin{align}
    \mathcal S_{\Sigma}
    =
    \bigcup_{c\in\{RG,GB,BR\}}
    \left(
        \mathcal S_{c}^{(A)}
        \cup
        \mathcal S_{c}^{(S)}
        \cup
        \mathcal S_{c}^{(T)}
    \right),
    \label{eq:rt_colored_strut_decomposition}
\end{align}
where a strut is active only when the resolved finite-thickness condition is
saturated,
\begin{align}
    d_s=2a_{\rm core},
    \qquad
    \lambda_s\ge0.
    \label{eq:rt_rank9_active_strut_condition}
\end{align}
The radius entering this condition is the resolved tube radius
\(a_{\rm core}\), not automatically the canonical horn radius \(\rc\).

For each channel \(I=(c,m)\), with
\(c\in\{RG,GB,BR\}\) and \(m\in\{A,S,T\}\), define the corresponding block of the
rigidity/contact matrix by the gradients of its active constraints,
\begin{align}
    \mathcal R_I
    :=
    \left[
        \nabla\!\left(\frac{d_s}{2}\right),
        \nabla\operatorname{MinRad}(v_k)
    \right]_{s,k\in I}.
    \label{eq:rt_rank9_channel_block}
\end{align}
The full reduced matrix is
\begin{align}
    \mathcal R_{\Sigma}^{\rm reduced}
    :=
    \left[\mathcal R_{RG}^{(A)},\mathcal R_{RG}^{(S)},\mathcal R_{RG}^{(T)},
    \mathcal R_{GB}^{(A)},\ldots,
    \mathcal R_{BR}^{(T)}\right]_{\rm reduced}.
    \label{eq:rt_rank9_reduced_matrix}
\end{align}
Reduction removes rigid translations, rotations, reparametrization modes, and
numerically duplicate constraints before the rank is computed.

\paragraph{Positive-rank criterion.}
Define \(\operatorname{rank}_{+}\) as the number of independent contact channels
that survive the KKT balance with nonnegative strut/kink multipliers.  The
factor nine is promoted from a matching diagnostic to a derived contact-channel
count only if
\begin{align}
    \boxed{
    N_{+}
    :=
    \operatorname{rank}_{+}
    \mathcal R_{\Sigma}^{\rm reduced}
    =9 .
    }
    \label{eq:rt_rank9_positive_rank_condition}
\end{align}
This is a rank-stability condition, not exact numerical orthogonality.

The ideal symmetry limit is block diagonal,
\begin{align}
    \mathcal R_{\Sigma}
    =
    \bigoplus_{c\in\{RG,GB,BR\}}
    \left(
        \mathcal R_c^{(A)}
        \oplus
        \mathcal R_c^{(S)}
        \oplus
        \mathcal R_c^{(T)}
    \right).
    \label{eq:rt_rank9_block_diagonal_limit}
\end{align}
A physical finite-thickness relaxation is expected to have controlled mixing,
\begin{align}
    \mathcal R_{\Sigma}
    =
    \mathcal R_0+\delta\mathcal R,
    \label{eq:rt_rank9_perturbed_rigidity_matrix}
\end{align}
where the leading cross-talk is expected between the \(S\) and \(T\) modes within
the same color root-plane.  The integer plateau survives if
\begin{align}
    \sigma_{\min}^{+}(\mathcal R_0)
    >
    \|\delta\mathcal R\|,
    \label{eq:rt_rank9_rank_stability_bound}
\end{align}
where \(\sigma_{\min}^{+}\) denotes the smallest nonzero positive-channel
singular scale after reduction.  Thus the claim is not that cross-talk vanishes;
the claim is that cross-talk does not collapse the positive rank.

A useful orthogonality diagnostic is the normalized channel Gram matrix
\begin{align}
    G_{IJ}
    =
    \frac{\langle\mathcal R_I,\mathcal R_J\rangle}
         {\|\mathcal R_I\|\,\|\mathcal R_J\|}
    =
    \delta_{IJ}+\varepsilon_{IJ},
    \label{eq:rt_rank9_channel_gram_matrix}
\end{align}
with reported mixing score
\begin{align}
    \eta_{\rm orth}
    :=
    \frac{\sum_{I\ne J}|G_{IJ}|}{\sum_I |G_{II}|}.
    \label{eq:rt_rank9_orthogonality_score}
\end{align}
A small \(\eta_{\rm orth}\) is supporting evidence, but Eq.~\eqref{eq:rt_rank9_positive_rank_condition}
is the decisive test.

\paragraph{Energetic decoupling diagnostic.}
For each channel \(I\), define the channel tension proxy
\begin{align}
    \sigma_I
    =
    \frac{1}{\ell_Y}
    \sum_{s\in\mathcal S_I}\lambda_s d_s,
    \label{eq:rt_rank9_channel_tension}
\end{align}
and the full numerical sheet tension
\begin{align}
    \sigma_{\rm num}
    =
    \frac{1}{\ell_Y}
    \sum_{s\in\mathcal S_{\Sigma}}\lambda_s d_s .
    \label{eq:rt_rank9_numerical_sheet_tension}
\end{align}
In the additive limit,
\begin{align}
    \sigma_{\rm num}
    =
    \sum_{I=1}^{9}\sigma_I .
    \label{eq:rt_rank9_additive_tension_limit}
\end{align}
A finite-thickness sheet may instead produce
\begin{align}
    \sigma_{\rm num}
    =
    \sum_{I=1}^{9}\sigma_I
    +
    \Delta\sigma_{\rm mix},
    \qquad
    \left|
    \frac{\Delta\sigma_{\rm mix}}{\sum_I\sigma_I}
    \right|
    \ll1 .
    \label{eq:rt_rank9_mixed_tension}
\end{align}
The factor-nine interpretation survives only if \(N_+=9\) and the mixing fraction
is stable or decreases under mesh refinement.

\paragraph{Connection to the geometric shielding gate.}
If the rank test succeeds, the candidate confinement-sheet scale may be compared
with the canon-level geometric shielding gate,
\begin{align}
    \frac{\lambda_c}{\pi\rc}
    =
    \frac{4}{\alpha},
    \label{eq:rt_rank9_shielding_gate_reference}
\end{align}
through the diagnostic estimate
\begin{align}
    \sigma_{\rm rank9}
    :=
    9
    \left(\frac{\lambda_c}{\pi\rc}\right)
    F_{\rm swirl}^{\max}
    =
    9\left(\frac{4}{\alpha}\right)F_{\rm swirl}^{\max}.
    \label{eq:rt_rank9_tension_diagnostic}
\end{align}
Equation~\eqref{eq:rt_rank9_shielding_gate_reference} is not derived by this
contact-rank lemma; it is imported from the calibrated geometric impedance chain.
The lemma tests whether there are nine independent positive contact channels to
which that gate could be applied.

\paragraph{Falsifiers.}
This research-track bridge fails, or remains a matching diagnostic only, if any
of the following occurs:
\begin{enumerate}
    \item the reduced positive rank satisfies \(N_+\ne9\);
    \item the KKT balance requires any negative strut or kink multiplier;
    \item the \(S\)--\(T\) cross-talk collapses two channels into one under refinement;
    \item the reported \(\sigma_{\rm num}\) is not stable under mesh refinement,
          symmetry-preserving perturbation, and independent initialization;
    \item the color root-plane assignment cannot be maintained as a
          \(\mathbb Z_3\)-locked triad around the central unknot.
\end{enumerate}

\paragraph{Canon-safe summary.}
\begin{align}
    \boxed{
    \text{The integer }9\text{ is not postulated; it is the target positive rank of the resolved colored contact-stress matrix.}
    }
    \label{eq:rt_rank9_canon_safe_summary}
\end{align}
Until Eq.~\eqref{eq:rt_rank9_positive_rank_condition} is numerically and
geometrically verified, Eq.~\eqref{eq:rt_rank9_tension_diagnostic} remains
\textbf{[MATCHING DIAGNOSTIC, NOT DERIVED]}.


\paragraph{Minimal falsifier for this appendix.}
If two independent high-resolution implementations, using the same thickness
convention and topology, converge to incompatible contact skeletons
$\mathcal C_K$ or incompatible $\chi_{\rm contact}$ under the same declared
$E_{\rm SST}$, the proposed contact-kernel refinement
Eq.~\eqref{eq:rt_kernel_contact_refinement} must be demoted from
Research Track to numerical diagnostic only.


\subsection{Minimal falsifier}
If two distinct knots have the same SST particle assignment but relax to incompatible energy minima or different protected sectors under the same admissible dynamics, the assignment must be split or rejected:
\begin{equation}
    K_1\sim_{\rm SST}K_2
    \quad\text{but}\quad
    \min E[K_1]\neq \min E[K_2]
    \quad\Rightarrow\quad
    \text{taxonomy refinement required.}
\end{equation}

\subsection{Research Appendix: Spherical Curl Spectrum and Helicity Capacity}
    \label{subsec:rt_spherical_curl_spectrum_capacity}

    \paragraph{Status.}
    \textbf{[ORTHODOX SOURCE RESULT / SST RESEARCH-TRACK APPLICATION].}
    This appendix records the bounded-domain curl-spectrum consequences of
    Cantarella, DeTurck, Gluck, and Teytel without identifying a spheromak with an
    SST particle state
    \cite{CantarellaDeTurckGluckTeytel2000CurlSpherical}.

    \paragraph{Helicity capacity.}
    For any domain on which the projected Biot--Savart/curl inverse is defined, set
    \begin{align}
        C_H(\Omega)
        &:=
        \frac{1}{\kappa_1(\Omega)},
        &
        [C_H]&=\mathrm{m}.
        \label{eq:rt_helicity_capacity}
    \end{align}
    This is the maximum helicity-to-\(L^2\)-energy ratio in the source convention.
    A scale-invariant resolved-tube diagnostic is
    \begin{align}
        \chi_{\rm curl}(K)
        :=
        a_{\rm core}\,\kappa_1(\Omega_K),
        \qquad
        [\chi_{\rm curl}]=1,
        \label{eq:rt_chi_curl}
    \end{align}
    where \(\Omega_K\) is the actual finite-thickness knot domain.  A pure
    ropelength mass branch requires \(\chi_{\rm curl}(K)\) and the normalized field
    profile to be sufficiently universal across the compared knot classes.

    \paragraph{Thin active-shell cost.}
    Let \(\delta=b-a\ll b\).  Expanding the spherical-shell root condition in the
    thin-shell regime gives the leading radial scaling
    \begin{align}
        \kappa_1(B^3(a,b))
        \sim
        \frac{\pi}{\delta}.
        \label{eq:rt_thin_shell_curl_scaling}
    \end{align}
    Thus a progressively thinner \emph{vorticity- or flux-carrying filled shell}
    acquires a large spectral cost.  This result must not be transferred to an
    irrotational horn envelope surrounding a much smaller vorticity-bearing core;
    those are different field sectors.

    \paragraph{Field-line geometry.}
    The minimum-eigenvalue ball and shell fields are axisymmetric and are organized
    by nested toroidal invariant surfaces.  This supports their use as qualitative
    benchmarks for a bent coarse-grained flux bundle.  It does not imply that a
    remote background flow must be toroidal, nor does it establish the geometry of a
    knotted SST carrier.

    \paragraph{Beltrami residual.}
    For a numerical candidate \(\mathbf{V}\), define
    \begin{align}
        \mathcal R_{\rm CK}(\kappa)
        :=
        \frac{
        \left\|
        \nabla\times\mathbf{V}-\kappa\mathbf{V}
        \right\|_{L^2(\Omega)}
        }{
        \left\|\nabla\times\mathbf{V}\right\|_{L^2(\Omega)}
        +
        |\kappa|\left\|\mathbf{V}\right\|_{L^2(\Omega)}
        }.
        \label{eq:rt_ck_residual}
    \end{align}
    The fitted \(\kappa\), \(\mathcal R_{\rm CK}\), and
    \(a_{\rm core}\kappa\) must be reported together.

    \paragraph{No automatic ideal-Euler relaxation.}
    A lowest curl eigenfield is a stationary variational benchmark.  In the ideal,
    incompressible, unforced Euler limit, kinetic energy and helicity are conserved;
    the dynamics therefore do not automatically dissipate toward the lowest mode.
    Any relaxation claim requires a separately declared dissipative, radiative, or
    coarse-graining mechanism.  Hence
    \begin{align}
        \text{small }\mathcal R_{\rm CK}
        \not\Rightarrow
        \text{dynamical attractor or particle stability}.
        \label{eq:rt_ck_no_relaxation_guard}
    \end{align}

    \paragraph{Minimal falsifier.}
    Compute \(\kappa_1\) on resolved unknot, trefoil, figure-eight, and selected
    higher-knot tubes at common \(a_{\rm core}\) and common volume.  Strong
    knot-dependent variation in \(\chi_{\rm curl}\) falsifies any claim that the
    entire field-energy splitting is determined by scalar ropelength alone.

\subsection{Research Appendix: Writhe Bounds, Twist Guards, and Flat-Torus Benchmark}
    \label{subsec:rt_upper_bounds_writhe_helicity}

    \paragraph{Status.}
    \textbf{[ORTHODOX SOURCE RESULTS / SST DIAGNOSTIC USE].}
    This appendix imports only the bounds and explicit model-domain results of
    Cantarella, DeTurck, and Gluck
    \cite{CantarellaDeTurckGluck2001UpperBounds}.

    \paragraph{Writhe sanity gate.}
    A computed centerline with radius-convention ropelength
    \(\operatorname{Rop}_{\rm rad}\) must satisfy
    \begin{align}
        |\operatorname{Wr}|
        <
        \frac14\operatorname{Rop}_{\rm rad}^{4/3}.
        \label{eq:rt_writhe_sanity_gate}
    \end{align}
    Violation indicates a convention, discretization, or integration error.  Passing
    the bound does not certify ideality because the inequality is intentionally broad.

    \paragraph{Special zero-internal-twist tube.}
    For the particular tube field used in the source---normal to each cross-sectional
    disk and with magnitude depending only on distance from the centerline---the
    internal twist contribution is absent and
    \begin{align}
        \mathscr H_V
        =
        \Phi^2\operatorname{Wr}(K).
        \label{eq:rt_special_zero_twist_helicity}
    \end{align}
    This equality is not a general SST identity.  It may be used only after the
    selected framing and cross-sectional field have independently established zero
    internal twist.  In particular,
    \begin{align}
        \mathscr H_V/\Phi^2
        =
        \operatorname{Wr}
        \quad\Longrightarrow\quad
        \text{verified zero-twist model},
        \label{eq:rt_zero_twist_implication_guard}
    \end{align}
    not a post-hoc choice of twist to match a target helicity.

    \paragraph{Flat solid-torus benchmark.}
    On the product domain \(D^2(a)\times S^1\), the first positive curl eigenvalue is
    \begin{align}
        \kappa_1
        =
        \frac{j_{0,1}}{a},
        \qquad
        j_{0,1}=2.404825558\ldots,
        \label{eq:rt_flat_solid_torus_eigenvalue}
    \end{align}
    with an eigenfield containing both longitudinal and azimuthal components.  This
    supplies a local slender-tube benchmark for \(a\kappa_1\).  The product torus is
    not an embedded curved horn torus in \(\mathbb R^3\); major-radius curvature,
    knotting, and self-contact are absent and must be recomputed in the physical
    domain.

    \paragraph{Research requirement.}
    For every candidate particle tube, report the tuple
    \begin{align}
        \left(
        \operatorname{Rop}_{\rm rad},
        \operatorname{Wr},
        \text{framing/twist prescription},
        \mathscr H_V/\Phi^2,
        a_{\rm core}\kappa_1
        \right).
        \label{eq:rt_writhe_helicity_audit_tuple}
    \end{align}
    A scalar ropelength or writhe value alone is insufficient to specify the
    volumetric field state.

\subsection{Research Appendix: Gehring-Link Criticality and Contact Measures}
    \label{subsec:rt_gehring_contact_measure}

    \paragraph{Status.}
    \textbf{[ORTHODOX VARIATIONAL GEOMETRY / SST CONTACT APPLICATION].}
    This appendix records the continuum balance theorem and its limitations from
    the Gehring link problem
    \cite{CantarellaFuKusnerSullivanWrinkle2006Gehring}.

    \paragraph{Continuum strut measure.}
    Let \(\operatorname{Strut}(L)\) be the set of intercomponent pairs realizing
    \(\operatorname{LThi}(L)\).  For a variation field \(\boldsymbol\xi\), define
    \begin{align}
        (A_S\boldsymbol\xi)(x,y)
        :=
        \frac{
        \left\langle
        \boldsymbol\xi(y)-\boldsymbol\xi(x),y-x
        \right\rangle
        }{|y-x|}.
        \label{eq:rt_gehring_rigidity_operator}
    \end{align}
    Strong criticality is equivalent to the existence of a positive Radon measure
    \(\mu\) on \(\operatorname{Strut}(L)\) satisfying
    \begin{align}
        \delta_{\boldsymbol\xi}\operatorname{Len}(L)
        &=
        \int_{\operatorname{Strut}(L)}
        A_S\boldsymbol\xi\,d\mu,
        &
        A_S^*\mu&=-\mathbf K,
        \label{eq:rt_gehring_continuum_balance}
    \end{align}
    where \(\mathbf K\) is the vector-valued curvature-force measure.

    For a resolved SST energy with both self-contact and curvature constraints, the
    corresponding first-order structure is
    \begin{align}
        -\frac{\delta E_{\rm SST}}{\delta\gamma}
        +
        A_S^*\mu
        +
        A_K^*\nu
        =0,
        \qquad
        \mu,\nu\geq0,
        \label{eq:rt_sst_measure_kkt_balance}
    \end{align}
    where \(\nu\) is the kink/curvature multiplier measure.

    \paragraph{Geometric contact versus load-bearing contact.}
    The multiplier measure need not have support on every active geometric strut:
    \begin{align}
        \operatorname{supp}\mu
        \subseteq
        \operatorname{Strut}(L),
        \label{eq:rt_load_bearing_strut_subset}
    \end{align}
    with possible strict inclusion.  Hence
    \(d(x,y)=a_i+a_j\) does not by itself imply nonzero contact pressure.  A solver
    must distinguish the active geometric strut set from the positive-multiplier
    load-bearing set.

    \paragraph{Discrete-to-continuum convergence.}
    A polygonal contact solution should be represented as
    \begin{align}
        \mu_N
        =
        \sum_{q\in\operatorname{Strut}(V_N)}
        \lambda_q^{(N)}\delta_q.
        \label{eq:rt_discrete_contact_measure}
    \end{align}
    Refinement requires weak-* convergence \(\mu_N\overset{*}{\rightharpoonup}\mu\),
    not merely convergence of the number of struts or total ropelength.  The total
    measure, support geometry, and selected test integrals \(\int f\,d\mu_N\) must be
    reported.

    \paragraph{Criticality is not stability.}
    Equation~\eqref{eq:rt_gehring_continuum_balance} is a first-order KKT balance.
    It does not establish a nonnegative constrained Hessian and does not establish
    stability under Euler or Biot--Savart time evolution.  The required hierarchy is
    \begin{align}
        \text{balanced}
        \quad\not\Rightarrow\quad
        \text{variationally stable}
        \quad\not\Rightarrow\quad
        \text{dynamically stable}.
        \label{eq:rt_gehring_stability_hierarchy}
    \end{align}

    \paragraph{Clasp singularity and finite-core regularization.}
    In the critical unit clasp, the tip curvature behaves as
    \begin{align}
        \kappa(s)
        \sim
        |s|^{-1/3},
        \label{eq:rt_gehring_clasp_curvature}
    \end{align}
    so the curve is \(C^{1,2/3}\), not \(C^{1,1}\), while the central inter-tip
    distance is approximately \(1.05639\) for unit link-thickness.  This demonstrates
    that intercomponent exclusion alone permits unbounded curvature and a non-contact
    gap at the visually closest tips.  A physical SST tube therefore requires the
    ordinary curvature/reach gate or a resolved bending/core regularization in
    addition to Gehring separation.

    \paragraph{Unequal component tensions.}
    For component-dependent effective tensions \(T_i\), the curvature-force side is
    weighted:
    \begin{align}
        A_S^*\mu
        =
        -\sum_i T_i\mathbf K_i.
        \label{eq:rt_gehring_unequal_tension_balance}
    \end{align}
    A geometrically symmetric multi-component state is not automatically critical
    when its components have unequal circulation, profile, or tension data.

    \paragraph{Higher link-homotopy data.}
    Pairwise linking numbers do not classify three-component Brunnian links.  The
    Borromean sector therefore requires a higher invariant such as Milnor's
    \(\mu_{123}\) in addition to \(Lk_{ij}\).  This is a topology label; it does not
    by itself establish the mechanical balance of Eq.~\eqref{eq:rt_sst_measure_kkt_balance}.

    \paragraph{Minimal falsifier.}
    A claimed multi-component SST equilibrium fails this gate if it preserves only
    intercomponent distance while allowing a component self-crossing, if any required
    multiplier is negative, if the load-bearing support does not converge under
    refinement, or if KKT balance is reported as dynamical stability without a
    separate perturbation spectrum.

\subsection{Research Appendix: Biot--Savart Realizability and Far-Field Gate}
    \label{subsec:rt_biot_savart_realizability}

    \paragraph{Status.}
    \textbf{[ORTHODOX OPERATOR RESULTS / SST IMPLEMENTATION REQUIREMENT].}
    This appendix maps the bounded-domain Biot--Savart theorems to SSTcore and
    VortexLab diagnostics
    \cite{CantarellaDeTurckGluck2001BiotSavart}.

    \paragraph{Exact vorticity-to-velocity inversion.}
    With \(\mathbf V=\boldsymbol\omega\), the raw reconstruction
    \(\mathbf u=\operatorname{BS}[\boldsymbol\omega]\) is a left inverse of curl
    only when
    \begin{align}
        \nabla\cdot\boldsymbol\omega=0,
        \qquad
        \boldsymbol\omega\cdot\mathbf n=0
        \quad\text{on }\partial\Omega.
        \label{eq:rt_biot_savart_vorticity_conditions}
    \end{align}
    Boundary leakage or discrete divergence produces the explicit gradient terms in
    Eq.~\eqref{eq:canon_biot_savart_curl_corrections}; it is not a harmless numerical
    detail.

    \paragraph{Potential-sector non-reconstruction.}
    Since boundary-normal gradients lie in \(\ker\operatorname{BS}\), a vorticity
    solver does not recover the complete pressure or potential-gradient sector.
    Therefore
    \begin{align}
        \boldsymbol\omega
        \text{ and }
        \operatorname{BS}[\boldsymbol\omega]
        \quad\not\Rightarrow\quad
        \text{unique pressure, contact stress, or gravity closure}.
        \label{eq:rt_biot_savart_pressure_nonreconstruction}
    \end{align}
    The pressure Poisson problem and its boundary data remain independent parts of
    the SST state.

    \paragraph{Realizability gate.}
    Because
    \(\operatorname{im}\operatorname{BS}\subsetneq\operatorname{im}\nabla\times\),
    a smooth divergence-free target field need not be generated by sources confined
    to the same domain.  Proposed analytic swirl profiles must therefore be tested
    against the projected Biot--Savart image rather than accepted from
    \(\nabla\cdot\mathbf u=0\) alone.

    \paragraph{Optical-size diagnostic.}
    Define the dimensionless geometry factor
    \begin{align}
        \chi_{\rm OS}(K)
        :=
        \frac{\operatorname{OS}(\Omega_K)}{a_{\rm core}},
        \qquad
        [\chi_{\rm OS}]=1.
        \label{eq:rt_biot_savart_optical_size_factor}
    \end{align}
    It supplies an operator-norm bound independent of scalar ropelength and should be
    reported when comparing domains with different compactness or concavity.

    \paragraph{Far-field regression test.}
    For an isolated compact fluid-knot source, fit
    \begin{align}
        \left\|\mathbf u(r)\right\|
        \propto
        r^{-p}
        \quad(r\gg\operatorname{diam}\Omega).
        \label{eq:rt_biot_savart_far_field_fit}
    \end{align}
    The direct Biot--Savart sector requires \(p\simeq3\).  A measured numerical
    exponent near \(1\) or \(2\) indicates a non-compact source, boundary leakage,
    an added background/potential sector, or an implementation error; it must not be
    silently attributed to the compact knot.

    \paragraph{Required residuals.}
    A bounded-domain reconstruction must report
    \begin{align}
        \epsilon_{\rm div}
        &:={
        \|\nabla\cdot\boldsymbol\omega\|_{L^2}
        \over
        \|\nabla\boldsymbol\omega\|_{L^2}},
        &
        \epsilon_{\rm wall}
        &:={
        \|\boldsymbol\omega\cdot\mathbf n\|_{L^2(\partial\Omega)}
        \over
        \|\boldsymbol\omega\|_{L^2(\Omega)}},
        \notag\\
        \epsilon_{\rm curl}
        &:={
        \|\nabla\times\mathbf u-\boldsymbol\omega\|_{L^2}
        \over
        \|\boldsymbol\omega\|_{L^2}}.
        \label{eq:rt_biot_savart_residuals}
    \end{align}
    The denominators and treatment of exact-zero cases must be declared in code.

\subsection{Research Appendix: Hodge-Complete SST Velocity Reconstruction}
    \label{subsec:rt_hodge_complete_reconstruction}

    \paragraph{Status.}
    \textbf{[ORTHODOX HODGE DECOMPOSITION / DERIVED SST RECONSTRUCTION].}
    The following reconstruction uses the domain topology explicitly
    \cite{CantarellaDeTurckGluck2002Hodge}.

    \paragraph{Harmonic basis and circulation coordinates.}
    Let
    \begin{align}
        g
        :=
        \dim H_1(\Omega;\mathbb R),
    \end{align}
    and choose harmonic-knot fields \(\mathbf h_1,\ldots,\mathbf h_g\) normalized by
    \begin{align}
        \oint_{\gamma_i}\mathbf h_j\cdot d\boldsymbol\ell
        =
        \delta_{ij}.
        \label{eq:rt_harmonic_basis_normalization}
    \end{align}
    A Hodge-complete bounded-domain velocity reconstruction has the form
    \begin{align}
        \boxed{
        \mathbf u
        =
        P_{K(\Omega)}\operatorname{BS}[\boldsymbol\omega]
        +
        \sum_{j=1}^{g}\Gamma_j\mathbf h_j
        }
        \label{eq:rt_hodge_complete_velocity}
    \end{align}
    after the flux and gauge conventions defining the projection have been fixed.
    The coefficients \(\Gamma_j\) are independent topological circulation data; they
    cannot be reconstructed from \(\boldsymbol\omega=\nabla\times\mathbf u\).

    \paragraph{Solid-torus sector.}
    For a solid torus, \(g=1\).  In the standard torus of revolution an illustrative
    harmonic-knot field is locally
    \begin{align}
        \mathbf h
        =
        \frac{1}{r}\,\widehat{\boldsymbol\phi},
        \qquad
        \nabla\cdot\mathbf h=0,
        \qquad
        \nabla\times\mathbf h=0.
        \label{eq:rt_solid_torus_harmonic_field}
    \end{align}
    Thus local curl-free flow may carry nonzero circulation around a non-contractible
    loop.  For the SST circulation quantum, the local slender-core scaling is
    \begin{align}
        u_{\rm HK}(r)
        =
        \frac{\Gamma_0}{2\pi r}
        =
        \vchar\frac{\rc}{r}.
        \label{eq:rt_harmonic_circulation_scaling}
    \end{align}
    This is a local benchmark outside an excluded core, not a determination of
    \(r_{\rm horn}\) or of the full curved-knot velocity profile.

    \paragraph{Irrotational envelope guard.}
    Equation~\eqref{eq:rt_hodge_complete_velocity} supplies a mathematically valid
    separation between a vorticity-bearing core sector and a curl-free global
    circulation sector.  It therefore permits an irrotational horn/envelope flow
    around a smaller core.  It does not prove that every SST horn is harmonic, nor
    does it fix the separatrix, pressure profile, or envelope radius.

    \paragraph{Hodge data are not particle identity.}
    Every tubular neighborhood of a single knot is topologically a solid torus, so
    \begin{align}
        \dim HK(\Omega_{3_1})
        =
        \dim HK(\Omega_{4_1})
        =1.
        \label{eq:rt_hodge_not_knot_classifier}
    \end{align}
    Hodge homology alone cannot distinguish a trefoil from a figure-eight knot.
    Particle classification still requires embedding-sensitive and knot-sensitive
    data such as knot group, chirality, writhe/twist, ropelength, contact map, and
    spectral geometry.

    \paragraph{Numerical circulation residual.}
    A bounded-domain solver must report
    \begin{align}
        \epsilon_{\Gamma}
        :=
        \max_j
        \frac{
        \left|
        \oint_{\gamma_j}\mathbf u\cdot d\boldsymbol\ell-\Gamma_j
        \right|
        }{
        \max(|\Gamma_j|,\Gamma_{\rm floor})
        },
        \label{eq:rt_hodge_circulation_residual}
    \end{align}
    with a declared nonzero \(\Gamma_{\rm floor}\).  A solver that matches
    \(\nabla\times\mathbf u\) but fails this circulation test is not Hodge-complete.

    \paragraph{Minimal research requirement.}
    For each resolved knot or link domain, compute a homology basis, normalize the
    harmonic fields, project the raw Biot--Savart field into \(K(\Omega)\), and verify
    both curl recovery and circulation recovery.  The stored state must include
    \begin{align}
        \left(
        \boldsymbol\omega,
        \{\Gamma_j\}_{j=1}^{g},
        \text{boundary conditions},
        \text{projection/gauge convention}
        \right).
        \label{eq:rt_hodge_state_metadata}
    \end{align}

\section{Research-Track particle and clock ansatz}
\label{app:research-track}

\[
\boxed{
\text{SUBSTRATE INTERPRETATION / NON-ORTHODOX RESEARCH TRACK}
}
\]

\paragraph{Status.}
This section introduces candidate SST closures that are not consequences of
orthodox Euler hydrodynamics, Gross--Pitaevskii theory, or established particle
physics.  The purpose is to formulate a constrained and falsifiable research
programme while preserving the distinction between geometric definitions,
canonically calibrated SST quantities, and unverified particle assignments.

Suppose, as an additional hypothesis, that microscopic matter structures are
finite-core vortex knots of an incompressible underlying substrate.  A
predictive implementation must derive its equilibrium states from a declared
continuum energy and finite-thickness constraint.  Arbitrary pairwise
short-range potentials introduced only to prevent strand crossings are
excluded.

\subsection{Resolved finite-core energy functional}

Let
\begin{equation}
\gamma_K:
S^1\rightarrow\mathbb{R}^3,
\qquad
s\mapsto\bm{X}_K(s)
\end{equation}
be an arclength-parametrized centreline of knot type \(K\).  Its resolved tube
radius is defined geometrically by
\begin{equation}
a_{\rm core}(K)
=
\operatorname{reach}(\gamma_K)
=
\min\left[
\inf_s\frac{1}{\kappa(s)},
\frac{1}{2}d_{\rm dcsd}(\gamma_K)
\right],
\label{eq:RT-reach}
\end{equation}
where \(d_{\rm dcsd}\) is the minimum doubly-critical self-distance
\cite{CantarellaKusnerSullivan2002}.

A candidate resolved-tube energy is
\begin{align}
\mathcal{E}_{\rm eff}[\gamma_K]
={}&
E_{\rm SI}^{(a)}[\gamma_K,\Gamma_K]
+
T_{\rm eff}L[\gamma_K]
+
B_{\rm eff}\oint_{\gamma_K}\kappa^2\,\dd s
+
E_{\rm twist}[\gamma_K],
\label{eq:RT-resolved-functional}
\\
L[\gamma_K]
={}&
\oint_{\gamma_K}\dd s .
\end{align}

Here \(E_{\rm SI}^{(a)}\) is a finite-core regularization of the kinetic
self-induction energy.  It may be represented generally as
\begin{equation}
E_{\rm SI}^{(a)}
=
\frac{\rhoF\Gamma_K^2}{8\pi}
\oint_{\gamma_K}
\oint_{\gamma_K}
\bm{t}(s)\cdot\bm{t}(s')\,
\mathcal{K}_{a_{\rm core}}
\left(
\left|
\bm{X}_K(s)-\bm{X}_K(s')
\right|
\right)
\dd s\,\dd s',
\label{eq:RT-regularized-self-induction}
\end{equation}
where
\begin{equation}
\mathcal{K}_{a}(r)
\longrightarrow
\frac{1}{r}
\qquad
(r\gg a)
\end{equation}
and \(\mathcal{K}_{a}(0)\) is finite.  The regularization kernel must be
derived from, or fixed by, a declared core profile before mass predictions
are evaluated.

The coefficients have dimensions
\begin{equation}
[T_{\rm eff}]
=
\mathrm{J\,m^{-1}},
\qquad
[B_{\rm eff}]
=
\mathrm{J\,m}.
\end{equation}
A finite-core scaling ansatz may take the form
\begin{equation}
B_{\rm eff}
=
c_B T_{\rm eff}a_{\rm core}^2,
\label{eq:RT-bending-scaling}
\end{equation}
where \(c_B\) is dimensionless.  Dimensional analysis does not determine
\(c_B=1\); it must follow from a resolved cross-sectional calculation.

The minimization problem is therefore
\begin{equation}
\min_{\gamma_K\in K}
\mathcal{E}_{\rm eff}[\gamma_K]
\qquad
\text{subject to}
\qquad
\operatorname{reach}(\gamma_K)
\geq
a_{\rm core}.
\label{eq:RT-constrained-minimization}
\end{equation}

For a polygonal discretization \(V\), this constraint becomes
\begin{align}
\operatorname{MinRad}(v_i)
&\geq
a_{\rm core},
\\
\frac{1}{2}d(p,q)
&\geq
a_{\rm core}.
\end{align}
The active curvature constraints are kinks and the active self-distance
constraints are struts.  Their non-negative KKT multipliers supply a
contact-stress skeleton for the resolved tube
\cite{AshtonCantarellaPiatekRawdon2011}.

The hard reach constraint is a geometric admissibility condition, not an
auxiliary Lennard--Jones interaction.  Bending stiffness can oppose
homothetic collapse, but it cannot by itself exclude a finite-curvature
self-crossing.

\subsection{Quantized circulation and characteristic velocity}

The circulation is retained in the canonical SST form
\begin{equation}
\Gamma_K
=
n_K\Gamma_0,
\qquad
n_K\in\mathbb{Z},
\label{eq:RT-circulation-quantization}
\end{equation}
with
\begin{equation}
\Gamma_0
=
2\pi\rc\vchar
=
9.68361920\times10^{-9}\,
\mathrm{m^2\,s^{-1}}.
\label{eq:RT-circulation-quantum}
\end{equation}

Thus the conserved circulation must not vary continuously with the knot
radius \(R_K\).  Instead, the characteristic resolved velocity may depend
on geometry through
\begin{equation}
U_K
=
\frac{\Gamma_K}{2\pi R_K}
\,
\mathcal{G}
\left(
\frac{a_{\rm core}}{R_K},
\mathcal{T}_K,
\mathfrak{q}_K
\right),
\label{eq:RT-characteristic-speed}
\end{equation}
where \(\mathcal{G}\) is dimensionless and \(\mathfrak{q}_K\) denotes any
independently specified internal coloring data.

Equation~\eqref{eq:RT-characteristic-speed} separates the canonical
circulation quantum from the geometry-dependent relation between circulation
and a chosen characteristic velocity.

\subsection{Topology-first particle assignment}

A predictive particle dictionary must be fixed before the observed mass
spectrum is used.  Let
\begin{equation}
\mathsf{A}(P)
=
\left(
K_P,
n_P,
\mathfrak{q}_P,
\text{framing data},
\text{component structure}
\right)
\label{eq:RT-assignment-map}
\end{equation}
denote the preregistered assignment for particle \(P\).

If a non-Abelian internal coloring is introduced, the mathematically
appropriate structure is a group homomorphism
\begin{equation}
\mathfrak{q}_K:
\pi_1
\left(
\mathbb{R}^3\setminus K
\right)
\longrightarrow
G,
\label{eq:RT-coloring-homomorphism}
\end{equation}
whose values on Wirtinger generators satisfy the knot-group crossing
relations.

A tuple such as
\begin{equation}
(L\bmod3,S\bmod2,\chi)
\end{equation}
may serve as a collection of discrete labels, but it is not by itself a
homomorphism into
\(\mathfrak{su}(3)\oplus\mathfrak{su}(2)\oplus\mathfrak{u}(1)\).

The trefoil \(K_{2,3}=3_1\) may be retained as the preregistered electron
benchmark of the current SST taxonomy.  Baryon assignments remain under
audit.  A torus-knot candidate \(K_{3,q}\), a three-component colored
composite, or another explicitly specified topology may be tested, but
threefold winding must not be identified with QCD color solely because
both involve the integer three.

\subsection{Preregistered falsifiability test}

Before fitting any mass, the following objects must be fixed:

\begin{enumerate}
\item the particle-to-topology map \(\mathsf{A}\);
\item the core regularization kernel \(\mathcal{K}_{a}\);
\item the circulation numbers \(n_P\);
\item the dimensionless coefficient \(c_B\);
\item the twist-energy law, if retained;
\item the admissible embedding class and numerical convergence criteria.
\end{enumerate}

Only one overall energy normalization may then be calibrated using the
electron mass.  Define the dimensionless minimized energy
\begin{equation}
\widehat{\mathcal{E}}_P
=
\min_{\gamma\in\mathcal{C}_P}
\widehat{\mathcal{E}}_{\rm eff}
\left[
\gamma;
\mathsf{A}(P)
\right],
\label{eq:RT-dimensionless-minimum}
\end{equation}
where \(\mathcal{C}_P\) is the preregistered admissible class.

The out-of-sample proton prediction is
\begin{equation}
\widehat{R}_{p/e}
=
\frac{
\widehat{\mathcal{E}}_p
}{
\widehat{\mathcal{E}}_e
}.
\label{eq:RT-predicted-proton-ratio}
\end{equation}
The model must reproduce
\begin{equation}
\widehat{R}_{p/e}
\simeq
1836.15267
\label{eq:RT-proton-target}
\end{equation}
within a preregistered numerical and theoretical uncertainty, without
changing the topology assignment, the kernel, or any dimensionless
coefficient after inspecting the result \cite{Mohr2025CODATA}.

The same fixed model must then predict additional ratios, including the
neutron-to-electron ratio and at least one further particle sector.  Failure
of any preregistered primary target falsifies the corresponding assignment
or the energy closure.

\subsection{Clock-rate discipline}

Define
\begin{equation}
\beta_K
=
\frac{U_K}{c},
\qquad
\Pi_a
=
\frac{a_{\rm core}}{R_K},
\qquad
\Pi_\kappa
=
a_{\rm core}^2\langle\kappa^2\rangle.
\label{eq:RT-clock-parameters}
\end{equation}

The canon-compatible local clock factor is
\begin{equation}
\frac{\dd\tau_K}{\dd T}
=
\sqrt{1-\beta_K^2}.
\label{eq:RT-canonical-clock}
\end{equation}

A Research-Track deviation may only be introduced in the explicitly
factorized form
\begin{equation}
\frac{\dd\tau_K}{\dd T}
=
\sqrt{1-\beta_K^2}
\left[
1+
\varepsilon_{\rm RT}
\Delta_K
\left(
\Pi_a,
\Pi_\kappa,
\mathcal{T}_K,
\mathfrak{q}_K,
\ldots
\right)
\right],
\label{eq:RT-clock-deviation}
\end{equation}
subject to
\begin{equation}
\Delta_K(0,0,\ldots)=0.
\end{equation}

The correction must be local, universal, independently derived, and
preregistered before experimental comparison.  A ratio involving the
galactic averaging radius \(R_{\rm gal}\) is not admitted into the local
clock law unless a specific non-local coupling mechanism is derived.

In the absence of such a derivation,
\begin{equation}
\varepsilon_{\rm RT}=0
\end{equation}
and Eq.~\eqref{eq:RT-canonical-clock} remains the operative clock law.

\subsection{Reconnection, framing, and twist}

A smooth filled Euler vortex tube has frozen topology and does not possess
the Gross--Pitaevskii reconnection mechanism.  The model must therefore
choose among:

\begin{enumerate}
\item exact topological conservation in the ideal limit;
\item a declared finite-scale nonideal reconnection law;
\item a compressible or order-parameter core theory.
\end{enumerate}

A framed finite-core tube may carry twist.  Any twist term
\(E_{\rm twist}\) must specify the material frame, its transport equation,
and the partition of helicity between linking, writhe, and twist
\cite{Moffatt1969}.  Candidate falsifiers include Kelvin-wave dispersion,
reconnection scaling, and helicity-transfer measurements.

\section{Minimal numerical programme}
\label{app:minimal-numerical-programme}

A staged numerical test must separate geometry, hydrodynamics, topology
assignment, and clock interpretation.

\begin{enumerate}

\item
Preregister the candidate particle topologies, circulation integers, core
profile, regularization kernel, and all dimensionless coefficients.

\item
Generate centreline discretizations for the preregistered knot and link
classes.  Measure
\[
L,\qquad
\operatorname{MinRad},\qquad
d_{\rm dcsd},\qquad
\operatorname{reach},
\qquad
\int\kappa^2\,\dd s.
\]

\item
Minimize Eq.~\eqref{eq:RT-resolved-functional} subject to the hard
finite-thickness conditions
\[
\operatorname{MinRad}\geq a_{\rm core},
\qquad
\frac12d_{\rm dcsd}\geq a_{\rm core}.
\]
Record the active struts, kinks, KKT multipliers, and constrained-gradient
residual.

\item
Perform core-resolution and discretization convergence tests.  The predicted
energy ratios must be stable under changes in vertex count, quadrature order,
kernel resolution, and optimizer tolerance.

\item
Calibrate only the overall energy scale on the electron.  Evaluate all
preregistered baryon and additional particle ratios without further fitting.

\item
For a fully contained closed vortex loop, verify the exact global identity
\begin{equation}
\int_{V}\bm{\omega}\,\dd V
=
\Gamma_K\oint_{\gamma_K}\dd\bm{X}
=
\bm{0}.
\label{eq:RT-closed-loop-zero-monopole}
\end{equation}

\item
Test the local coarse-grained relation
\begin{equation}
\overline{\omega}_z(\bm{x})
\simeq
\Gamma_K
\mathcal{L}(\bm{x})
p_z(\bm{x})
\label{eq:RT-local-polarization}
\end{equation}
only as a local segment average.  For fully contained closed loops, verify
that positive axial-vorticity regions are balanced by return-vorticity
regions.

\item
Compare separately:
\begin{enumerate}
\item straight vortex lines crossing the averaging boundary;
\item boundary-anchored or periodic vortex lines;
\item fully closed vortex knots;
\item space-filling vortex-cell models.
\end{enumerate}
Only the first two classes can support a nonzero global vector-vorticity
monopole without an explicit return region.

\item
Measure finite-size and homogenization errors as
\[
\frac{L_{\rm cg}}{R_{\rm gal}}\rightarrow0,
\qquad
\frac{\ell_v}{L_{\rm cg}}\rightarrow0.
\]

\item
Use Gross--Pitaevskii simulations for depleted-core quantum superfluids and
Euler, vortex-blob, or resolved finite-core methods for filled classical
tubes.  Do not transfer reconnection laws between these models without a
declared bridge.

\item
Do not introduce a modified clock law until the circulation, velocity,
energy, and topology closures have passed their independent tests.

\end{enumerate}



% ----------------------------------------------------------
% Bibliography entries to merge into the main thebibliography
% ----------------------------------------------------------

\ifresearchtrackincluded\else
\begin{thebibliography}{99}

\bibitem{Annala2022ProtectedVortexKnots}
T.~Annala, R.~Zamora-Zamora, and M.~M{\"o}tt{\"o}nen,
``Topologically protected vortex knots and links,''
\emph{Communications Physics} \textbf{5}, 309 (2022).
DOI: \href{https://doi.org/10.1038/s42005-022-01071-2}{10.1038/s42005-022-01071-2}.

\bibitem{Kleckner2016SuperfluidKnotsUntie}
D.~Kleckner, L.~H.~Kauffman, and W.~T.~M.~Irvine,
``How superfluid vortex knots untie,''
\emph{Nature Physics} \textbf{12}, 650--655 (2016).
DOI: \href{https://doi.org/10.1038/nphys3679}{10.1038/nphys3679}.

\bibitem{Pisanty2018FractionalKnotsLight}
E.~Pisanty, G.~Jim{\'e}nez, V.~Vicu{\~n}a-Hern{\'a}ndez, A.~Pic{\'o}n, A.~Celi, J.~P.~Torres, and M.~Lewenstein,
``Knotting fractional-order knots with the polarization state of light,''
\emph{arXiv:1808.05193} (2018).
Permalink: \href{https://arxiv.org/abs/1808.05193}{arXiv:1808.05193}.

\bibitem{CantarellaKusnerSullivan2002}
J.~Cantarella, R.~B. Kusner, and J.~M. Sullivan,
``On the Minimum Ropelength of Knots and Links,''
\emph{Inventiones Mathematicae} \textbf{150}, 257--286 (2002).
DOI: \href{https://doi.org/10.1007/s00222-002-0234-y}{10.1007/s00222-002-0234-y};
arXiv: \href{https://arxiv.org/abs/math/0103224}{math/0103224}.

\bibitem{AshtonCantarellaPiatekRawdon2011}
T.~Ashton, J.~Cantarella, M.~Piatek, and E.~J. Rawdon,
``Knot Tightening by Constrained Gradient Descent,''
\emph{Experimental Mathematics} \textbf{20}(1), 57--90 (2011).
DOI: \href{https://doi.org/10.1080/10586458.2011.544581}{10.1080/10586458.2011.544581};
arXiv: \href{https://arxiv.org/abs/1002.1723}{1002.1723}.

\bibitem{CantarellaDeTurckGluckTeytel2000Isoperimetric}
J.~Cantarella, D.~DeTurck, H.~Gluck, and M.~Teytel,
``Isoperimetric problems for the helicity of vector fields and the Biot--Savart and curl operators,''
\emph{Journal of Mathematical Physics} \textbf{41}, 5615--5641 (2000).
DOI: \href{https://doi.org/10.1063/1.533429}{10.1063/1.533429}.


\bibitem{Ricca1998KnotFluidMechanics}
R.~L.~Ricca,
``Applications of knot theory in fluid mechanics,''
\emph{Banach Center Publications} \textbf{42}, 321--346 (1998).
Permalink: \href{https://eudml.org/doc/252224}{EuDML:252224}.

\bibitem{Moffatt1969HelicityKnottedness}
H.~K.~Moffatt,
``The degree of knottedness of tangled vortex lines,''
\emph{Journal of Fluid Mechanics} \textbf{35}, 117--129 (1969).
DOI: \href{https://doi.org/10.1017/S0022112069000991}{10.1017/S0022112069000991}.

\bibitem{CantarellaParsley2010Helicity}
J.~Cantarella and J.~Parsley,
``A new cohomological formula for helicity in $\mathbb R^{2k+1}$ reveals the effect of a diffeomorphism on helicity,''
\emph{Journal of Geometry and Physics} \textbf{60}, 1127--1155 (2010).
DOI: \href{https://doi.org/10.1016/j.geomphys.2010.04.001}{10.1016/j.geomphys.2010.04.001};
arXiv: \href{https://arxiv.org/abs/0903.1465}{0903.1465}.

\bibitem{ArnoldKhesin1998TopologicalMethods}
V.~I.~Arnold and B.~A.~Khesin,
\emph{Topological Methods in Hydrodynamics},
Springer (1998).
DOI: \href{https://doi.org/10.1007/978-1-4612-0645-3}{10.1007/978-1-4612-0645-3}.

\end{thebibliography}
\fi


% ============================================================
% legacy insertion block, retained in v0.8.17 research track:
% SST CANON-v0.8.1 Research-Track Appendix Pack
% Topic: Resonance Selection, Time Kernels, Swirl Impedance,
%        and Layered Relaxation in Swirl--String Theory
% Status: [RESEARCH-TRACK], not main-canon proof.
% ============================================================

% --- Minimal macro safety for standalone insertion ---
\providecommand{\vswirl}{\mathbf{v}_{\!\boldsymbol{\circlearrowleft}}}
\providecommand{\vnorm}{\lVert \mathbf{v}_{\!\boldsymbol{\circlearrowleft}}\rVert}
\providecommand{\rhoF}{\rho_{\!f}}
\providecommand{\rhoE}{\rho_{\!E}}
\providecommand{\rhoM}{\rho_{\!m}}
\providecommand{\rhocore}{\rho_{\text{core}}}
\providecommand{\rc}{r_c}
\providecommand{\omC}{\omega_C}
\providecommand{\Fswirlmax}{F_{\text{swirl}}^{\max}}

% ============================================================

\section{Swirl-String Response and Resonance Framework}

\subsection{Resonance Selection Principle}
    \label{sec:resonance-selection-principle}
    
    \noindent
    \textbf{Status:} [RESEARCH-TRACK; candidate Canon principle].
    
    A swirl-string state \(K\) is assigned a causal response function
    \(\chi_K(\omega)\). An external perturbation \(F(t)\), with Fourier
    spectrum \(\widetilde F(\omega)\), excites the state \(K\) with amplitude
    \begin{equation}
        \mathcal A_K[F]
        =
        \int_0^\infty
        \widetilde F(\omega)\,
        \chi_K(\omega)\,
        d\omega .
    \end{equation}
    The excitation probability is
    \begin{equation}
        P_K[F]
        =
        |\mathcal A_K[F]|^2 .
    \end{equation}
    
    For isolated narrow resonances, the response function may be written as
    \begin{equation}
        \chi_K(\omega)
        =
        \sum_n
        \frac{B_{Kn}\gamma_{Kn}}
        {\omega-\omega_{Kn}+i\gamma_{Kn}},
    \end{equation}
    where \(\omega_{Kn}\) are internal swirl-string eigenfrequencies,
    \(\gamma_{Kn}\) are linewidths, and \(B_{Kn}\) are coupling weights.
    
    The natural SST frequency scale is
    \begin{equation}
        \Omega_{\rm SST}
        =
        \frac{\vchar}{r_c},
    \end{equation}
    so that internal modes may be parametrized by
    \begin{equation}
        \omega_{Kn}
        =
        \eta_{Kn}\Omega_{\rm SST},
    \end{equation}
    with dimensionless mode factors \(\eta_{Kn}\).
    
    This principle unifies delay-loop selection, Swirl-Clock phase-locking,
    mode-locking, and knot-state excitation as instances of one spectral
    overlap law.
    
    
\subsection{Resonance Selection Principle for Swirl--String Excitation}
    \label{app:resonance-selection-principle}
    
    \subsubsection*{Status and purpose}
        
        \noindent
        \textbf{Status:} [RESEARCH-TRACK].
        This appendix promotes the frequency-overlap idea from a specific beam--swirl model to a general SST excitation principle. It is not a proof of nuclear activation or fusion enhancement. It is a controlled response ansatz compatible with the present canon language of swirl strings, mode selection, circulation, helicity, and Kelvin-compatible dynamics.
    
    \subsubsection*{Claim}
        
        A swirl-string state \(K\) is excited most efficiently when the spectrum of an external perturbation overlaps the internal response spectrum of \(K\). Define a normalized external spectral density
        \begin{equation}
            p_{\mathrm{ext}}(\omega)\ge 0,
            \qquad
            \int_0^\infty p_{\mathrm{ext}}(\omega)\,d\omega = 1,
            \label{eq:p_ext_normalized}
        \end{equation}
        where \(p_{\mathrm{ext}}\) has units of seconds, because \(d\omega\) has units \(\mathrm{s}^{-1}\). Let the dimensionless response function of the swirl-string state be
        \begin{equation}
            R_K(\omega)
            =
            \sum_{n}
            B_{Kn}
            \frac{\gamma_{Kn}^{2}}
            {(\omega-\omega_{Kn})^{2}+\gamma_{Kn}^{2}},
            \label{eq:RK_lorentzian}
        \end{equation}
        where:
        \begin{itemize}
            \item \(\omega_{Kn}\) is the \(n\)-th internal mode frequency,
            \item \(\gamma_{Kn}\) is the corresponding linewidth,
            \item \(B_{Kn}\) is a dimensionless coupling weight.
        \end{itemize}
        
        The dimensionless excitation yield is then
        \begin{equation}
            \boxed{
                \mathcal{Y}_K
                =
                \int_0^\infty
                p_{\mathrm{ext}}(\omega)\,
                R_K(\omega)\,
                d\omega
            }
            \qquad
            [\mathrm{RESEARCH\text{-}TRACK}].
            \label{eq:SST_resonance_yield}
        \end{equation}
    
    \subsubsection*{Dimensional check}
        
        Since
        \[
            [p_{\mathrm{ext}}]=\mathrm{s},
            \qquad
            [d\omega]=\mathrm{s}^{-1},
            \qquad
            [R_K]=1,
        \]
        we obtain
        \[
            [\mathcal{Y}_K]
            =
            \mathrm{s}\cdot 1 \cdot \mathrm{s}^{-1}
            =
            1.
        \]
        Thus \(\mathcal{Y}_K\) is a dimensionless excitation efficiency.
    
    \subsubsection*{Canonical scale}
        
        The natural SST angular-frequency scale is
        \begin{equation}
            \Omega_{\mathrm{SST}}
            =
            \frac{\vnorm}{\rc}.
            \label{eq:Omega_SST_scale}
        \end{equation}
        Using the Canon values
        \[
            \vnorm = 1.09384563\times10^{6}\ \mathrm{m\,s^{-1}},
            \qquad
            \rc = 1.40897017\times10^{-15}\ \mathrm{m},
        \]
        gives
        \begin{equation}
            \Omega_{\mathrm{SST}}
            =
            7.76344066\times10^{20}\ \mathrm{s^{-1}}.
            \label{eq:Omega_SST_numeric}
        \end{equation}
        The corresponding cycle period is
        \begin{equation}
            T_{\mathrm{SST}}
            =
            \frac{2\pi}{\Omega_{\mathrm{SST}}}
            =
            8.09329985\times10^{-21}\ \mathrm{s}.
            \label{eq:T_SST_numeric}
        \end{equation}
    
    \subsubsection*{Interpretation}
        
        Equation~\eqref{eq:SST_resonance_yield} does not assert that all nuclear, atomic, or particle transitions are already derived in SST. It states a weaker and safer rule:
        
        \begin{quote}
            External excitation is strongest when the perturbing spectrum overlaps an allowed internal swirl-string response spectrum.
        \end{quote}
        
        This principle can later be specialized to:
        \begin{itemize}
            \item radiative transitions,
            \item qubit-like R/T phase transitions,
            \item knot-mode excitation,
            \item delayed mode selection,
            \item controlled response experiments in optical, plasma, or condensed-matter analogues.
        \end{itemize}
    
    \subsubsection*{Falsifiability}
        
        For a fixed target state \(K\), SST predicts that changing the spectral center \(\omega_0\) of a narrow-band perturbation should change the activation yield according to \(R_K(\omega_0)\), up to normalization and detector response. If no resonance-like dependence appears after controlling for beam power, fluence, and thermal effects, then Eq.~\eqref{eq:SST_resonance_yield} is falsified as a useful excitation law.

% ============================================================
\subsection{Time-Domain Kernel Form of Swirl--String Response}
    \label{app:time-domain-kernel-response}
    
    \subsubsection*{Status and purpose}
        
        \noindent
        \textbf{Status:} [RESEARCH-TRACK].
        This appendix converts the frequency-domain selection rule into a time-domain response law. It is intended to connect SST mode selection to pulsed experiments, delay-loop systems, and Swirl-Clock phase locking.
    
    \subsubsection*{Pulse-response formulation}
        
        Let \(f(t)\) be a real external driving pulse. A Gaussian-windowed monochromatic pulse may be written as
        \begin{equation}
            f(t)
            =
            f_0
            \exp\!\left(-\frac{t^2}{2\tau_p^2}\right)
            \cos(\omega_0 t+\phi_0),
            \label{eq:gaussian_drive_pulse}
        \end{equation}
        where \(\tau_p\) is the pulse duration and \(\omega_0\) is the carrier angular frequency.
        
        Let \(k_K(t)\) be the causal impulse-response kernel of the swirl-string state \(K\), with
        \begin{equation}
            k_K(t<0)=0.
        \end{equation}
        The induced response is
        \begin{equation}
            \boxed{
                s_K(t)
                =
                \int_{-\infty}^{t}
                f(t')\,k_K(t-t')\,dt'
            }
            \qquad
            [\mathrm{ORTHODOX\ form;\ RESEARCH\text{-}TRACK\ SST\ interpretation}].
            \label{eq:SST_convolution_response}
        \end{equation}
    
    \subsubsection*{Frequency-domain equivalence}
        
        With the Fourier convention
        \[
            \hat f(\omega)
            =
            \int_{-\infty}^{\infty}
            f(t)e^{i\omega t}\,dt,
        \]
        the convolution theorem gives
        \begin{equation}
            \hat s_K(\omega)
            =
            \hat f(\omega)\,\chi_K(\omega),
            \label{eq:frequency_response_chi}
        \end{equation}
        where \(\chi_K(\omega)\) is the complex response function of the swirl-string state. The measurable spectral response is proportional to
        \begin{equation}
            R_K(\omega)
            \propto
            |\chi_K(\omega)|^2.
            \label{eq:response_modulus}
        \end{equation}
    
    \subsubsection*{Bandwidth--duration relation}
        
        A Gaussian pulse obeys the scaling
        \begin{equation}
            \Delta\omega\,\tau_p \sim 1.
            \label{eq:time_bandwidth_relation}
        \end{equation}
        Therefore:
        \begin{itemize}
            \item short pulses have broad spectra and can excite many nearby modes;
            \item long pulses have narrow spectra and selectively excite one mode;
            \item coherent excitation requires phase stability over at least one internal response time.
        \end{itemize}
        
        For the canonical SST scale,
        \[
            T_{\mathrm{SST}} = 8.09329985\times10^{-21}\ \mathrm{s}.
        \]
        Thus, a pulse with duration \(\tau_p\gg T_{\mathrm{SST}}\) is narrow relative to the core-scale carrier, whereas \(\tau_p\lesssim T_{\mathrm{SST}}\) is broadband at the canonical core scale.
    
    \subsubsection*{Connection to Swirl Clock}
        
        If the internal phase of a swirl string is denoted by \(\theta_K(t)\), then a resonant perturbation can be represented as a phase-locking term
        \begin{equation}
            \frac{d\theta_K}{dt}
            =
            \omega_K
            +
            \epsilon\,f(t)\sin(\omega_0 t-\theta_K),
            \label{eq:phase_locking_response}
        \end{equation}
        where \(\epsilon\) is a coupling parameter. Phase locking occurs when
        \begin{equation}
            |\omega_0-\omega_K| \lesssim \epsilon |f_0|.
            \label{eq:locking_condition}
        \end{equation}
        This gives a direct bridge between frequency-selective excitation and Swirl-Clock synchronization.
    
    \subsubsection*{Falsifiability}
        
        For a given \(K\), changing only the pulse duration \(\tau_p\) while keeping total pulse energy fixed should alter the response through the bandwidth \(\Delta\omega\). A null result after controlling for heating and total fluence would falsify this time-domain response model.

% ============================================================
\subsection{Swirl Impedance and Mode Linewidth}
    \label{app:swirl-impedance-linewidth}
    
    \subsubsection*{Status and critical correction}
        
        \noindent
        \textbf{Status:} [RESEARCH-TRACK].
        The earlier expression
        \[
            Z_{\mathrm{eff}} \sim \frac{P_{\mathrm{core}}\rc}{\vnorm\hbar}
        \]
        should \emph{not} be canonised as written. Its dimensions are
        \[
            \left[
                \frac{P_{\mathrm{core}}\rc}{\vnorm\hbar}
            \right]
            =
            \mathrm{m^{-3}},
        \]
        so it is not an impedance and not dimensionless.
        
        A Canon-compatible replacement is to define either:
        
        \begin{enumerate}
            \item a dimensionless rigidity ratio,
            \item or an acoustic-like pressure-to-speed impedance.
        \end{enumerate}
    
    \subsubsection*{Dimensionless swirl-rigidity ratio}
        
        Let \(V_K\) be the effective active volume of a swirl-string configuration and \(P_K\) the pressure scale associated with its deformation. Define
        \begin{equation}
            \boxed{
                \Pi_K(\omega)
                =
                \frac{P_K V_K}{\hbar\omega}
            }
            \qquad
            [\mathrm{RESEARCH\text{-}TRACK}].
            \label{eq:dimensionless_swirl_rigidity}
        \end{equation}
        This compares stored mechanical energy \(P_KV_K\) to the quantum energy scale \(\hbar\omega\). It is dimensionless:
        \[
            [P_KV_K]=\mathrm{J},
            \qquad
            [\hbar\omega]=\mathrm{J}.
        \]
    
    \subsubsection*{Acoustic-like swirl impedance}
        
        An alternative intensive impedance is
        \begin{equation}
            \boxed{
                Z_{\mathrm{swirl}}
                =
                \frac{P_K}{v_K}
            }
            \label{eq:swirl_acoustic_impedance}
        \end{equation}
        with units
        \[
            [Z_{\mathrm{swirl}}]
            =
            \frac{\mathrm{kg\,m^{-1}\,s^{-2}}}{\mathrm{m\,s^{-1}}}
            =
            \mathrm{kg\,m^{-2}\,s^{-1}},
        \]
        the same units as acoustic impedance. Here $v_K$ denotes the coarse-grained velocity amplitude of the same mode family $K$ used to define $P_K$; it may be instantiated as a phase-transport speed, local carrier-speed amplitude, or resonant surface-speed amplitude, but the choice must be declared in any benchmark. It is not the canonical scalar speed $\vchar$ unless the benchmark explicitly specializes to the canonical carrier limit. This object measures resistance to velocity-driven pressure excitation.
    
    \subsubsection*{Connection to linewidth and quality factor}
        
        A resonant mode with angular frequency \(\omega_K\) and quality factor \(Q_K\) has linewidth
        \begin{equation}
            \gamma_K
            =
            \frac{\omega_K}{2Q_K}.
            \label{eq:linewidth_quality_factor}
        \end{equation}
        A minimal phenomenological relation is
        \begin{equation}
            Q_K
            =
            Q_0\,\Pi_K^\eta,
            \qquad
            \eta>0,
            \label{eq:Q_rigidity_ansatz}
        \end{equation}
        so that larger stored-energy rigidity produces narrower response lines.
    
    \subsubsection*{Numerical check using canonical scales}
        
        Using
        \[
            \Fswirlmax = 29.053507\ \mathrm{N},
            \qquad
            \rc = 1.40897017\times10^{-15}\ \mathrm{m},
        \]
        define the core pressure scale
        \begin{equation}
            P_{\max}
            =
            \frac{\Fswirlmax}{\pi \rc^2}
            =
            4.65848920\times10^{30}\ \mathrm{Pa}.
            \label{eq:Pmax_numeric}
        \end{equation}
        \textbf{[CRITICAL NOTE]}: $V_{\mathrm{horn}}=2\pi^2\rc^3$ is the $\rc$-scale envelope (horn-torus) volume, \emph{not} the resolved tube cross-section, for which Sec.~2.2 stipulates $a_{\mathrm{core}}\neq\rc$. The resulting $\Pi_K\approx\pi$ is hence a geometric normalization at the envelope scale, not a stiffness theorem for the physical tube; a profile-resolved $a_{\mathrm{core}}$ model is required for the latter.
        For a horn-torus active volume
        \begin{equation}
            V_{\mathrm{horn}}
            =
            2\pi^2 \rc^3
            =
            5.52122107\times10^{-44}\ \mathrm{m^3},
            \label{eq:horn_volume_numeric}
        \end{equation}
        the stored pressure energy is
        \begin{equation}
            P_{\max}V_{\mathrm{horn}}
            =
            2.57205487\times10^{-13}\ \mathrm{J}.
            \label{eq:pressure_energy_numeric}
        \end{equation}
        At the canonical Compton angular frequency
        \[
            \omega_c
            =
            7.76344066\times10^{20}\ \mathrm{s^{-1}},
        \]
        we have
        \begin{equation}
            \hbar\omega_c
            =
            8.18710572\times10^{-14}\ \mathrm{J}
            =
            510998.946\ \mathrm{eV}.
            \label{eq:hbarOmega_numeric}
        \end{equation}
        Therefore,
        \begin{equation}
            \Pi_{\mathrm{horn}}
            =
            \frac{P_{\max}V_{\mathrm{horn}}}
            {\hbar\omega_c}
            =
            3.14159235
            \approx \pi.
            \label{eq:Pi_horn_numeric}
        \end{equation}
    
    \subsubsection*{Interpretation}
        
        Equation~\eqref{eq:Pi_horn_numeric} is not yet a derivation of a particle mass or coupling constant. It is a useful dimensional coincidence produced by canonical SST scales and the horn-torus active-volume choice. It should be labelled:
        \[
            \boxed{
                \text{[CALIBRATED / RESEARCH-TRACK NUMERICAL STRUCTURE]}
            }
        \]
        until derived from the SST action or from a topological minimization principle.
    
    \subsubsection*{Falsifiability}
        
        If \(\Pi_K\) controls linewidth, then families of analog swirl-string excitations should satisfy
        \[
            \gamma_K \propto \omega_K \Pi_K^{-\eta}.
        \]
        A measured absence of correlation between active-volume rigidity and linewidth would falsify Eq.~\eqref{eq:Q_rigidity_ansatz}.

% ============================================================
\subsection{Layered Relaxation and Multi-Mode Decay}
    \label{app:layered-relaxation}
    
    \subsubsection*{Status and warning}
        
        \noindent
        \textbf{Status:} [RESEARCH-TRACK].
        This appendix preserves the useful part of the delayed-decay idea: a multi-layer swirl-string system may relax through multiple modes. It explicitly rejects the naive identification
        \[
            \tau_i = \frac{r_i}{\vnorm}
        \]
        with observed delayed-neutron group lifetimes. For nuclear-scale radii, this gives \(\tau_i\sim 10^{-21}\ \mathrm{s}\), whereas delayed-neutron groups are macroscopic nuclear decay times. Therefore direct radial-transit mapping is dimensionally possible but physically wrong in scale.
    
    \subsubsection*{General multi-mode relaxation}
        
        If a perturbed swirl-string state has several metastable channels, its observable relaxation can be expanded as
        \begin{equation}
            \boxed{
                A(t)
                =
                \sum_{i=1}^{N}
                A_i
                e^{-t/\tau_i}
                \cos(\omega_i t+\phi_i)
            }
            \qquad
            [\mathrm{ORTHODOX\ form;\ RESEARCH\text{-}TRACK\ SST\ interpretation}].
            \label{eq:multimode_relaxation}
        \end{equation}
        The non-oscillatory limit is
        \begin{equation}
            A(t)
            =
            \sum_{i=1}^{N}
            A_i e^{-t/\tau_i}.
            \label{eq:multigroup_decay}
        \end{equation}
    
    \subsubsection*{Pole interpretation}
        
        Equation~\eqref{eq:multimode_relaxation} corresponds to a response function with poles
        \begin{equation}
            \omega
            =
            \omega_i
            -
            \frac{i}{\tau_i}.
            \label{eq:response_poles}
        \end{equation}
        Thus, each decay time \(\tau_i\) is a lifetime of a response pole, not necessarily a physical radial crossing time.
    
    \subsubsection*{Allowed SST interpretation}
        
        In SST, a layered or nested swirl-string configuration may contain:
        \begin{itemize}
            \item a core-scale response,
            \item a sheath-scale response,
            \item a linking/helicity response,
            \item a global Swirl-Clock phase response.
        \end{itemize}
        The decay constants \(\tau_i\) should therefore be treated as \emph{effective metastable relaxation times}. A safe parametrization is
        \begin{equation}
            \tau_i
            =
            \frac{Q_i}{\omega_i},
            \label{eq:tau_Q_over_omega}
        \end{equation}
        where \(Q_i\) is the quality factor of the \(i\)-th mode.
    
    \subsubsection*{Critical scale check}
        
        For the canonical circulation scale,
        \begin{equation}
            \tau_{\rm cross,c}
            =
            \frac{\rc}{\vnorm}
            =
            1.28808868\times10^{-21}\ \mathrm{s}.
            \label{eq:core_crossing_time_numeric}
        \end{equation}
        This is a single radial crossing time, distinct from the full circulation
        period \(\tau_{\rm circ,c}=2\pi\rc/\vchar\) of Eq.~\eqref{eq:compton_delay}
        in the main canon.
        Therefore, an observed lifetime of \(1\ \mathrm{s}\) would require
        \begin{equation}
            Q
            \sim
            \omega_i\tau_i
            \sim
            7.76\times10^{20}.
            \label{eq:large_Q_warning}
        \end{equation}
        Such a huge \(Q\)-factor is not impossible as a formal metastability parameter, but it cannot be interpreted as a simple geometric shell crossing.
    
    \subsubsection*{Experimental analogy}
        
        Multi-exponential decay models are common in nuclear delayed-neutron kinetics and in linear response theory. SST may use this only as an analogy unless a topological barrier calculation derives the \(Q_i\) or \(\tau_i\) values independently.
    
    \subsubsection*{Falsifiability}
        
        The layered relaxation ansatz predicts that perturbing a structured target should produce reproducible decay constants \(\tau_i\) independent of total pulse amplitude, provided the response remains linear. If the apparent \(\tau_i\) vary strongly with drive amplitude, heating, or detector threshold, the model is likely observing instrumental or thermal relaxation rather than swirl-string metastability.

% ============================================================
\subsection*{Summary of Canon Placement}
    \label{app:canon-placement-summary}
    
    The four additions should be placed as follows:
    
    \begin{center}
        \begin{tabular}{|p{0.25\linewidth}|p{0.22\linewidth}|p{0.42\linewidth}|}
            \hline
            \textbf{Appendix} & \textbf{Canon status} & \textbf{Recommended placement} \\
            \hline
            Resonance Selection Principle &
            [RESEARCH-TRACK] strong candidate &
            CANON-v0.8.17 companion appendix; possible future main-canon principle after validation. \\
            \hline
            Time-Domain Kernel Response &
            [RESEARCH-TRACK] strong candidate &
            Near Swirl Clock, delay-mode selection, or quantum-response sections. \\
            \hline
            Swirl Impedance and Linewidth &
            [RESEARCH-TRACK] with corrected dimensions &
            Use as a response/linewidth ansatz; do not canonise the older dimensionally incorrect formula. \\
            \hline
            Layered Relaxation &
            [RESEARCH-TRACK] analogy only &
            Keep as metastable multi-mode response; reject direct delayed-neutron shell-crossing interpretation. \\
            \hline
        \end{tabular}
    \end{center}

% ============================================================
    \ifresearchtrackincluded\else
\begin{thebibliography}{99}
        
        \bibitem{Kubo1957}
        R. Kubo. (1957).
        \textit{Statistical-Mechanical Theory of Irreversible Processes. I. General Theory and Simple Applications to Magnetic and Conduction Problems}.
        Journal of the Physical Society of Japan \textbf{12}, 570--586.
        doi: \href{https://doi.org/10.1143/JPSJ.12.570}{10.1143/JPSJ.12.570}.
        
        \bibitem{BreitWigner1936}
        G. Breit and E. Wigner. (1936).
        \textit{Capture of Slow Neutrons}.
        Physical Review \textbf{49}, 519--531.
        doi: \href{https://doi.org/10.1103/PhysRev.49.519}{10.1103/PhysRev.49.519}.
        
        \bibitem{Batchelor1967}
        G. K. Batchelor. (1967).
        \textit{An Introduction to Fluid Dynamics}.
        Cambridge University Press.
        
        \bibitem{Saffman1992}
        P. G. Saffman. (1992).
        \textit{Vortex Dynamics}.
        Cambridge University Press.
        doi: \href{https://doi.org/10.1017/CBO9780511624063}{10.1017/CBO9780511624063}.
        
        \bibitem{Helmholtz1858}
        H. von Helmholtz. (1858).
        \textit{\"Uber Integrale der hydrodynamischen Gleichungen, welche den Wirbelbewegungen entsprechen}.
        Journal f\"ur die reine und angewandte Mathematik \textbf{55}, 25--55.
        doi: \href{https://doi.org/10.1515/crll.1858.55.25}{10.1515/crll.1858.55.25}.
        
        \bibitem{Kelvin1867}
        W. Thomson, Lord Kelvin. (1867).
        \textit{On Vortex Atoms}.
        Proceedings of the Royal Society of Edinburgh \textbf{6}, 94--105.
        
        \bibitem{Onsager1949}
        L. Onsager. (1949).
        \textit{Statistical Hydrodynamics}.
        Il Nuovo Cimento \textbf{6}, 279--287.
        doi: \href{https://doi.org/10.1007/BF02780991}{10.1007/BF02780991}.
        
        \bibitem{Feynman1955}
        R. P. Feynman. (1955).
        \textit{Application of Quantum Mechanics to Liquid Helium}.
        In C. J. Gorter (ed.), \textit{Progress in Low Temperature Physics}, Vol. 1, pp. 17--53.
        North-Holland.
        doi: \href{https://doi.org/10.1016/S0079-6417(08)60077-3}{10.1016/S0079-6417(08)60077-3}.
        
        \bibitem{Bracewell2000}
        R. N. Bracewell. (2000).
        \textit{The Fourier Transform and Its Applications}.
        3rd ed., McGraw--Hill.
        
        \bibitem{Keepin1965}
        G. R. Keepin. (1965).
        \textit{Physics of Nuclear Kinetics}.
        Addison-Wesley.
        
        \bibitem{Dhibar2018}
        M. Dhibar, I. Mazumdar, G. Anil Kumar, S. P. Weppner, A. K. Rhine Kumar, S. M. Patel, and P. B. Chavan. (2018).
        \textit{Study of \(^{12}\mathrm C(p,p'\gamma)^{12}\mathrm C\) reaction}.
        arXiv:1806.02126.
        \href{https://arxiv.org/abs/1806.02126}{https://arxiv.org/abs/1806.02126}.
        
        \bibitem{Rose2018}
        P. B. Rose Jr. et al. (2018).
        \textit{High-energy gamma rays resulting from low-energy nuclear reactions}.
        Physical Review C \textbf{97}, 064305.
        doi: \href{https://doi.org/10.1103/PhysRevC.97.064305}{10.1103/PhysRevC.97.064305}.
    
    \end{thebibliography}
\fi




\section{Extended Vortex, Helicity, and Kairos Research Tracks}

\subsection{Research Track: Nested Vortex--Antivortex States as Entropy-Minimizing Swirl Structures}
    
    \subsection{Motivation}
        
        Consider a coherent swirl structure consisting of an outer circulation region
        with vorticity $\omega_{\mathrm{out}}>0$ and an embedded inner core with
        opposite vorticity
        
        \[
            \omega_{\mathrm{in}}<0.
        \]
        
        Such configurations appear naturally in fluid dynamics, magnetic flux ropes,
        and vortex-ring interactions.
        
        Within SST, opposite swirl chirality corresponds to opposite temporal
        orientation of the Swirl Clock:
        
        \[
            S_t^{\circlearrowleft}
            \quad\leftrightarrow\quad
            \text{matter},
        \]
        
        \[
            S_t^{\circlearrowright}
            \quad\leftrightarrow\quad
            \text{antimatter}.
        \]
    
    \subsection{Helicity Cancellation}
        
        The total helicity may be decomposed as
        
        \[
            H_{\mathrm{tot}}
            =
            H_{\mathrm{out}}
            +
            H_{\mathrm{in}}.
        \]
        
        Because the embedded structure carries opposite chirality,
        
        \[
            H_{\mathrm{in}}<0.
        \]
        
        Hence
        
        \[
            |H_{\mathrm{tot}}|
            <
            |H_{\mathrm{out}}|.
        \]
        
        The configuration therefore possesses reduced net topological complexity.
    
    \subsection{Entropy Hypothesis}
        
        We introduce the conjecture
        
        \[
            S_{\mathrm{top}}
            \propto
            |H|,
        \]
        
        where $S_{\mathrm{top}}$ denotes topological entropy.
        
        Consequently,
        
        \[
            S_{\mathrm{top}}^{\mathrm{nested}}
            <
            S_{\mathrm{top}}^{\mathrm{single}},
        \]
        
        suggesting that nested vortex--antivortex states may represent preferred
        low-energy attractors of the swirl medium.
    
    \subsection{Status}
        
        \textbf{Research Track.}
        
        Not yet derived from the canonical SST Lagrangian.
        Proposed as a candidate mechanism for stability enhancement and
        matter--antimatter coexistence.
    
    
\subsection{Research Track: Helicity Conversion and Torsional Radiation}
    
    \subsection{Motivation}
        
        Experiments and simulations of knotted vortex structures indicate that
        reconnection events frequently preserve a substantial fraction of helicity.
        
        Rather than disappearing, helicity is redistributed between knotting,
        linking, writhe, and twist.
    
    \subsection{Helicity Partition}
        
        Write total helicity as
        
        \[
            H
            =
            H_{\mathrm{knot}}
            +
            H_{\mathrm{twist}}
            +
            H_{\mathrm{link}}.
        \]
        
        During reconnection:
        
        \[
            H_{\mathrm{knot}}
            \downarrow
        \]
        
        while
        
        \[
            H_{\mathrm{twist}}
            \uparrow .
        \]
        
        Thus
        
        \[
            H_{\mathrm{tot}}
            \approx
            \mathrm{const}.
        \]
    
    \subsection{SST Interpretation}
        
        Within SST, mass-bearing T-phase structures correspond to localized knot
        helicity.
        
        Radiative R-phase excitations correspond to propagating torsional modes.
        
        We therefore conjecture
        
        \[
            H_{\mathrm{knot}}
            \rightarrow
            H_{\mathrm{twist}}
            \rightarrow
            \gamma_R ,
        \]
        
        where $\gamma_R$ denotes an R-phase torsional pulse.
    
    \subsection{Mass--Radiation Conversion Chain}
        
        The proposed conversion pathway is
        
        \[
            M
            \leftrightarrow
            H_{\mathrm{knot}}
            \leftrightarrow
            H_{\mathrm{twist}}
            \leftrightarrow
            \gamma_R.
        \]
        
        In this picture:
        
        \begin{itemize}
            \item stable particles store helicity as topology,
            \item reconnections release helicity as torsional twist,
            \item torsional twist propagates as photon-like R-phase radiation.
        \end{itemize}
    
    \subsection{Connection to Photon Ontology}
        
        This hypothesis is consistent with the SST identification
        
        \[
            \text{photon}
            =
            \text{torsional R-phase pulse},
        \]
        
        and provides a possible mechanism for
        
        \begin{itemize}
            \item particle decay,
            \item annihilation,
            \item radiation emission,
            \item energy transport.
        \end{itemize}
    
    \subsection{Falsifiable Prediction}
        
        If the hypothesis is correct, vortex reconnection experiments should show
        
        \[
            \Delta H_{\mathrm{knot}}
            \approx
            -
            \Delta H_{\mathrm{twist}}
        \]
        
        with emitted torsional wave packets carrying the missing helicity.
    
    \subsection{Status}
        
        \textbf{Research Track.}
        
        Motivated by helicity-conservation studies in knotted vortex systems.
        Requires derivation from the SST radiation sector.
    
    
\subsection{Research Track: Kairos Surfaces as Chirality Transition Layers}
    
    \subsection{Definition}
        
        The Temporal Ontology introduces the Kairos event
        
        \[
            \kappa,
        \]
        
        representing an irreversible topological transition.
        
        We extend this concept spatially.
        
        Let
        
        \[
            \omega(\mathbf{x})=0
        \]
        
        define a surface separating opposite chiral domains:
        
        \[
            \omega>0
            \quad\leftrightarrow\quad
            S_t^{\circlearrowleft},
        \]
        
        \[
            \omega<0
            \quad\leftrightarrow\quad
            S_t^{\circlearrowright}.
        \]
        
        We define the locus
        
        \[
            \Sigma_\kappa
            =
            \{
            \mathbf{x}
            \;|\;
            \omega(\mathbf{x})=0
            \}
        \]
        
        as a \emph{Kairos Surface}.
    
    \subsection{Physical Interpretation}
        
        Across $\Sigma_\kappa$ the swirl orientation reverses sign.
        
        The surface therefore acts as a boundary between distinct temporal sectors.
        
        \[
            S_t^{\circlearrowleft}
            \longleftrightarrow
            S_t^{\circlearrowright}.
        \]
        
        Such regions may act as:
        
        \begin{itemize}
            \item annihilation interfaces,
            \item vortex reconnection sites,
            \item phase-transition boundaries,
            \item chirality-conversion layers.
        \end{itemize}
    
    \subsection{Kairos Criterion}
        
        A local Kairos transition is conjectured whenever
        
        \[
            \omega
            \rightarrow
            0
        \]
        
        and
        
        \[
            \nabla\omega
            \neq
            0.
        \]
        
        The gradient supplies a preferred transition direction.
    
    \subsection{Status}
        
        \textbf{Research Track.}
        
        Consistent with the Temporal Ontology but not yet derived from the SST field
        equations.

\section{EM, Gravity, Diagnostics, Relativity, and Inertia Research Tracks}

\subsection{Euler-Decomposed Swirl Gravity and Probe Transport}
\label{sec:appendix_euler_probe_transport}

\subsection{Purpose and status}
    This appendix separates the canonical scalar pressure acceleration from
    orientation-dependent transverse transport effects. The pressure law is retained
    as the canon-level bridge, while Magnus-type terms are treated as probe-level
    corrections. The motivation is the orthodox transverse-force/Magnus result for
    moving vortices in acoustic-geometry and superfluid settings
    \cite{Zhang2005TransverseForce,Saffman1992}.

    \textbf{[GUARDRAIL]} The transverse term is not promoted here to the primary
    cause of universal free fall. In the present canon, universal fall and the
    Newtonian limit remain attached to the scalar pressure/clock sector. Magnus-type
    contributions are orientation-, circulation-, and chirality-sensitive probe
    transport corrections until a universal macroscopic coupling law is derived.

\subsection{Canonical pressure acceleration}
    \begin{align}
        \mathbf{a}_{P}
        =
        -\frac{1}{\rho_{\!f}}\nabla p_{\mathrm{swirl}}.
        \label{eq:rt_scalar_pressure_acceleration}
    \end{align}

\subsection{Euler decomposition}
    For an incompressible inviscid background flow,
    \begin{align}
        \frac{\partial \mathbf{u}}{\partial t}
        +
        \boldsymbol{\omega}\times\mathbf{u}
        +
        \nabla\left(\frac{u^2}{2}\right)
        =
        -\frac{1}{\rho_{\!f}}\nabla p_{\mathrm{swirl}}.
        \label{eq:rt_euler_decomposed_swirl_balance}
    \end{align}

\subsection{Filament Magnus transport}
    For a circulation-carrying swirl-string probe $K$,
    \begin{align}
        \mathbf{F}_{\mathrm{M}}^{(K)}
        =
        \int_K
        \rho_{\!f}\Gamma_K
        \mathbf{T}(s)\times
        \left[
            \mathbf{U}_K(s)-\mathbf{u}(\mathbf{X}(s))
        \right]ds.
        \label{eq:rt_filament_magnus_transport}
    \end{align}

\subsection{Curvature self-transport}
    In the local-induction approximation,
    \begin{align}
        \mathbf{f}_{\kappa}
        \simeq
        -\frac{\rho_{\!f}\Gamma_K^2}{4\pi}
        \Lambda_K\kappa\,\mathbf{N}.
        \label{eq:rt_curvature_self_transport}
    \end{align}

\subsection{Dimensional and density check}
    Equation~\eqref{eq:rt_filament_magnus_transport} is a force on a closed
    filament. The corresponding local density
    \begin{align}
        \mathbf{f}_{\mathrm{M}}
        =
        \rho_{\!f}\Gamma_K\,
        \mathbf{T}\times(\mathbf{U}_K-\mathbf{u})
    \end{align}
    has units
    \begin{align}
        [\rho_{\!f}\Gamma_K\Delta U]
        =
        \mathrm{kg\,m^{-3}}\,\mathrm{m^2\,s^{-1}}\,\mathrm{m\,s^{-1}}
        =
        \mathrm{N\,m^{-1}}.
    \end{align}
    With the canonical values
    \begin{align}
        \Gamma_0
        &=
        2\pi \rc\vchar
        =
        9.6836192\times10^{-9}\ \mathrm{m^2\,s^{-1}},\\
        \rho_{\!f}\Gamma_0
        &=
        6.7785334\times10^{-15}
        \frac{\mathrm{N\,m^{-1}}}{\mathrm{m\,s^{-1}}},\\
        \rho_{\!f}\Gamma_0\vchar
        &=
        7.4146692\times10^{-9}\ \mathrm{N\,m^{-1}}.
    \end{align}
    Thus, with \(\rho_{\!f}\), the Magnus layer is a weak background-transport
    term rather than a replacement for the saturated core force scale
    \(F_{\mathrm{swirl}}^{\max}=29.053507\ \mathrm{N}\). Replacing
    \(\rho_{\!f}\) by \(\rho_{\mathrm{core}}\) would change the physical sector and
    is not allowed in this transport estimate.

\subsection{Axis/off-axis swirl interpretation}
    For an axisymmetric galactic swirl field
    \begin{align}
        \mathbf{u}=u_\theta(r)\mathbf{e}_\theta,
        \qquad
        \boldsymbol{\omega}=\omega_z(r)\mathbf{e}_z,
    \end{align}
    an exactly axial probe suppresses the transverse Magnus correction by symmetry.
    Off-axis motion samples the azimuthal flow. For a flat-tail branch
    \(u_\theta(r)\to v_0\),
    \begin{align}
        \omega_z
        =
        \frac{1}{r}\frac{d}{dr}\left(r v_0\right)
        =
        \frac{v_0}{r},
        \qquad
        \boldsymbol{\omega}\times\mathbf{u}
        =
        -\frac{v_0^2}{r}\mathbf{e}_r.
    \end{align}
    The inward sign is therefore already contained in the Euler decomposition of
    the background swirl. This observation should be used as an interpretation of
    probe drift and centripetal balance, not as an independent gravity postulate.

\subsection{Status}
    The scalar pressure term is canon-level. The Magnus and curvature terms are
    research-track transport corrections until a universal coupling law for
    macroscopic matter is derived. Their role is to describe path curvature,
    chiral drift, and frame-dragging-like probe effects on top of the scalar
    pressure/clock bridge.
    
    % ============================================================
% Research Track Appendix: Swirl--Electromagnetic Normalization
% ============================================================
    

% ============================================================
% Research Track Appendix: Hybrid density-source Swirl-Clock benchmark
% ============================================================

\subsection{Research Track: Hybrid Density-Source Swirl-Clock Benchmark}
\label{app:research_hybrid_density_source_swirl_clock}

% ------------------------------------------------------------
% Local macro safety block
% ------------------------------------------------------------
\providecommand{\swirlarrow}{\boldsymbol{\circlearrowleft}}
\providecommand{\vswirl}{\mathbf{v}_{\!\boldsymbol{\circlearrowleft}}}
\providecommand{\rhoF}{\rho_{\!f}}
\providecommand{\rhoE}{\rho_{\!E}}
\providecommand{\rhoM}{\rho_{\!m}}
\providecommand{\SwirlClock}{S_{(t)}^{\boldsymbol{\circlearrowleft}}}

\paragraph{Status.}
\textbf{[RESEARCH TRACK / BENCHMARK CANDIDATE].}
This subsection records a dimensionally corrected benchmark ansatz for coupling
local Swirl-Clock dilation to density-source gravitation across microscopic and
macroscopic regimes.  It is not a canonical master equation.  It may be promoted
only if the transition function, source density, and rotational-length scale are
independently derived from SST medium dynamics and pass the recovery tests
listed below.

\paragraph{Hybrid clock ansatz.}
For a localized swirl-string configuration $K$, define the candidate clock factor
\begin{equation}
\boxed{
\SwirlClock_{\rm hyb}(r;K)
=
\left[
1
-
\frac{2G_{\rm hyb}(r)M_{\rm hyb}(r;K)}{r c^2}
-
\frac{\left\|\mathbf{u}_{\rm rel}+\mathbf{v}_{\theta}(r;K)\right\|^2}{c^2}
-
\frac{\ell_{\Omega}^{2}(r;K)\Omega_{\rm int}^{2}(r;K)}{c^2}
\right]^{1/2}
}
\label{eq:rt_hybrid_density_source_swirl_clock}
\end{equation}
where $\mathbf{u}_{\rm rel}$ is the relative velocity of the probe or localized
configuration with respect to the background foliation, $\mathbf{v}_{\theta}$ is
the resolved circular swirl velocity, and $\ell_{\Omega}\Omega_{\rm int}$ is the
linear speed associated with an unresolved internal rotational mode.

\paragraph{Hybrid interpolation.}
The benchmark interpolation is written as
\begin{align}
\mu(r) &= \exp\!\left[-\left(\frac{r}{R_0}\right)^2\right],
\label{eq:rt_hybrid_mu}\\[2mm]
G_{\rm hyb}(r)
&=
\mu(r)G_{\rm swirl}+\bigl[1-\mu(r)\bigr]G_N,
\label{eq:rt_hybrid_G}\\[2mm]
M_{\rm hyb}(r;K)
&=
\mu(r)M_{\rm eff}^{\rm SST}(r;K)
+\bigl[1-\mu(r)\bigr]M_{\rm bulk}.
\label{eq:rt_hybrid_M}
\end{align}
Here $R_0$ is a research-track transition length, not a canonical constant, and
$M_{\rm bulk}$ denotes the independently measured macroscopic source mass used
in the Newtonian recovery regime.

\paragraph{Density-source mass closure.}
The SST-compatible source is the mass-equivalent swirl-energy density, not a bare
geometric density:
\begin{equation}
\boxed{
M_{\rm eff}^{\rm SST}(r;K)
=
\int_{B_r}\rho_{\!m}^{\rm eff}(\mathbf{x};K)\,d^3x,
\qquad
\rho_{\!m}^{\rm eff}
=
\frac{\rhoE}{c^2}
=
\frac{\rhoF\,\|\vswirl\|^2}{2c^2}.
}
\label{eq:rt_density_source_mass_closure}
\end{equation}
This form keeps the source term aligned with the canonical density hierarchy
$\rhoE=\tfrac12\rhoF\|\vswirl\|^2$ and $\rhoM=\rhoE/c^2$.

\paragraph{Phenomenological spherical test profile.}
For numerical benchmarks one may use the non-canonical profile
\begin{equation}
\rho_{\!m}^{\rm eff}(r)=\rho_0 e^{-r/r_*},
\end{equation}
which gives the enclosed source mass
\begin{equation}
\boxed{
M_{\rm eff}^{\rm exp}(r)
=
8\pi\rho_0 r_*^3
\left[
1-
\left(
1+\frac{r}{r_*}
+\frac{1}{2}\frac{r^2}{r_*^2}
\right)e^{-r/r_*}
\right].
}
\label{eq:rt_exponential_density_source_mass}
\end{equation}
The small-radius limit is
\begin{equation}
M_{\rm eff}^{\rm exp}(r)
=
\frac{4\pi}{3}\rho_0 r^3+O(r^4),
\qquad r\ll r_*,
\end{equation}
so the enclosed source remains regular at the origin.  Omitting the quadratic
term in Eq.~\eqref{eq:rt_exponential_density_source_mass} breaks this recovery
limit and is therefore inadmissible.

\paragraph{Irrotational compression and boundary response.}
The exterior circular-flow branch used for bubble-boundary benchmarks is
\begin{equation}
\mathbf{v}_{\theta}(r;K)
=\frac{\Gamma_K}{2\pi r}\,\hat{\boldsymbol{\theta}},
\qquad r\ge r_{\rm core},
\label{eq:rt_irrotational_boundary_velocity}
\end{equation}
with pressure response inherited from Euler radial balance,
\begin{equation}
\frac{1}{\rhoF}\frac{dp_{\rm swirl}}{dr}
=\frac{v_{\theta}^{2}}{r}.
\label{eq:rt_boundary_pressure_balance}
\end{equation}
Thus compression of the boundary at fixed circulation increases
$v_{\theta}\propto r^{-1}$, raises the local kinetic energy density, and lowers
the Swirl-Clock factor through Eq.~\eqref{eq:rt_hybrid_density_source_swirl_clock}.
Energy addition may be represented by increasing $\Gamma_K$, by increasing the
internal mode speed $\ell_{\Omega}\Omega_{\rm int}$, or by changing the boundary
radius; these alternatives must be reported separately in benchmarks.

\paragraph{Dimensional guard for internal rotation.}
A bare factor $\Omega^2/c^2$ is dimensionally inadmissible.  The only admissible
clock contribution from an angular rate is
\begin{equation}
\boxed{
\frac{v_{\Omega}^{2}}{c^2}
=
\frac{\ell_{\Omega}^{2}(r;K)\Omega_{\rm int}^{2}(r;K)}{c^2},
}
\label{eq:rt_omega_dimensional_guard}
\end{equation}
where $\ell_{\Omega}$ is the length scale on which the internal angular mode is
realized.  If $\ell_{\Omega}=r_c$ and $\Omega_{\rm int}=\|\vswirl\|/r_c$, this term
reduces to $\|\vswirl\|^2/c^2$ and must not be counted again as an independent
correction.

\paragraph{Required benchmark reporting.}
Any numerical use of Eq.~\eqref{eq:rt_hybrid_density_source_swirl_clock} must
report the four dimensionless contributions separately:
\begin{equation}
\mathcal{B}_{G}
=\frac{2G_{\rm hyb}M_{\rm hyb}}{r c^2},
\qquad
\mathcal{B}_{\theta}
=\frac{\|\mathbf{u}_{\rm rel}+\mathbf{v}_{\theta}\|^2}{c^2},
\qquad
\mathcal{B}_{\Omega}
=\frac{\ell_{\Omega}^{2}\Omega_{\rm int}^{2}}{c^2},
\qquad
\mu(r).
\end{equation}
At minimum the benchmark set must include electron-core, proton-core, Earth,
Sun, and neutron-star limits.  The candidate is rejected if it fails the weak
field recovery $\SwirlClock_{\rm hyb}\to\sqrt{1-2G_NM_{\rm bulk}/(r c^2)}$ for
$\mu(r)\to0$, or if any contribution is dimensionally non-scalar.

\subsection{Research Track: Swirl--Electromagnetic Normalization by Vorticity}
\label{app:research_swirl_em_vorticity}

% ------------------------------------------------------------
% Local macro safety block
% ------------------------------------------------------------
\providecommand{\swirlarrow}{\boldsymbol{\circlearrowleft}}
\providecommand{\vswirl}{\mathbf{v}_{\!\boldsymbol{\circlearrowleft}}}
\providecommand{\vnorm}{\lVert\mathbf{v}_{\!\boldsymbol{\circlearrowleft}}\rVert}
\providecommand{\vchar}{v_{\!\boldsymbol{\circlearrowleft}}^{\ast}}
\providecommand{\uswirl}{\mathbf{u}_{\!\boldsymbol{\circlearrowleft}}}
\providecommand{\omegaswirl}{\boldsymbol{\omega}_{\!\boldsymbol{\circlearrowleft}}}
\providecommand{\rhoF}{\rho_{\!f}}
\providecommand{\rhoE}{\rho_{\!E}}
\providecommand{\rhoM}{\rho_{\!m}}
\providecommand{\rc}{r_c}
\providecommand{\Fmax}{F_{\max}^{\rm swirl}}

\subsection{Purpose and Status}
    
    This appendix records a research-track extension of Swirl--String Theory
    (SST) in which electromagnetic field variables are interpreted as
    SI-normalized coarse-grained measures of swirl kinematics.  The result is
    not proposed as a replacement of Maxwell's equations at the level of
    observed electrodynamics.  Instead, it is a constitutive bridge between
    the SST swirl-medium variables and the ordinary electromagnetic variables
    \(\mathbf{E}\), \(\mathbf{B}\), \(\mathbf{A}\), \(\mathbf{D}\), and
    \(\mathbf{H}\).
    
    The central observation is that a dimensionally correct identification of
    the vector potential is obtained by scaling the local swirl velocity by the
    electron mass-to-charge ratio:
    \begin{equation}
        \boxed{
            \mathbf{A}_{\rm SI}
            =
            \frac{m_e}{e}\,
            \uswirl(\mathbf{x},t)
        } .
        \label{eq:swirl_vector_potential_bridge}
    \end{equation}
    Since \([\mathbf{A}]={\rm kg\,m\,C^{-1}\,s^{-1}}\) and
    \([(m_e/e)\uswirl]={\rm kg\,C^{-1}}\cdot{\rm m\,s^{-1}}\), this equation
    has the correct SI units. This use of \(\uswirl(\mathbf{x},t)\) is intentional: the EM bridge is a local-field bridge, not a curl of the calibrated scalar speed \(\vchar\).

    \textbf{[CRITICAL NOTE]} Eq.~\eqref{eq:swirl_vector_potential_bridge} is a constitutive normalization, not a proof that the full nonlinear Euler vorticity dynamics are isomorphic to vacuum Maxwell theory. The bridge is admissible only in a coarse-grained, transverse, weak-response regime where advective and vortex-stretching terms are either projected out, averaged away, or shown to be subleading by an explicit perturbation estimate.
    
    The associated magnetic field is then
    \begin{equation}
        \boxed{
            \mathbf{B}_{\rm SST}
            =
            \nabla\times\mathbf{A}_{\rm SI}
            =
            \frac{m_e}{e}\,
            \boldsymbol{\omega}
        },
        \qquad
        \boldsymbol{\omega}
        =
        \nabla\times\uswirl(\mathbf{x},t) .
        \label{eq:B_from_vorticity}
    \end{equation}
    This is the cleanest form of the statement that magnetic flux arises from
    structured vorticity, not merely from charge flow.  The charge-current
    description remains the effective macroscopic description; the SST
    interpretation places the microscopic origin of the magnetic field in
    quantized circulation and vorticity.

\subsection{Why Direct Vorticity Source Terms Are Rejected}
    
    Earlier trial equations attempted to write vorticity as an additional
    source term directly inside Maxwell's equations, for example terms of the
    schematic form
    \begin{equation}
        \nabla\cdot\mathbf{E}
        \sim
        \frac{\rho}{\varepsilon_0}
        +
        \frac{F_{\max}\omega}{\vnorm \rc^2},
        \qquad
        \nabla\times\mathbf{B}
        \sim
        \mu_0\mathbf{J}
        +
        \frac{\vnorm\hbar}{q\rc^2}\boldsymbol{\Omega}.
        \label{eq:rejected_direct_source_terms}
    \end{equation}
    These forms are not retained as canonical equations because the added
    terms do not generally have the SI dimensions required by Maxwell's
    source equations.  In particular,
    \(\nabla\cdot\mathbf{E}\) has units
    \({\rm V\,m^{-2}}\), while the proposed vorticity-force expression does
    not reduce to this unit without an additional calibrated constitutive
    coefficient.  Likewise, the Ampere--Maxwell equation requires units of
    \({\rm T\,m^{-1}}\), not a force-per-charge expression multiplied by an
    unfixed vorticity vector.
    
    Therefore, the research-track correction is:
    
    \begin{quote}
        Vorticity should not be inserted as an uncalibrated extra Maxwell source.
        It should enter through a constitutive identification of the potential and
        field variables.
    \end{quote}
    
    The retained bridge is thus Eq.~\eqref{eq:swirl_vector_potential_bridge},
    not Eq.~\eqref{eq:rejected_direct_source_terms}.

\subsection{Electric Field from Time-Dependent Swirl}
    
    With the vector potential bridge established, the electric field follows
    from the usual potential relation:
    \begin{equation}
        \boxed{
            \mathbf{E}_{\rm SST}
            =
            -\frac{\partial\mathbf{A}_{\rm SI}}{\partial t}
            -
            \nabla\Phi_{\rm SST}
            =
            -\frac{m_e}{e}
            \frac{\partial\uswirl(\mathbf{x},t)}{\partial t}
            -
            \nabla\Phi_{\rm SST}
        } .
        \label{eq:E_from_time_dependent_swirl}
    \end{equation}
    Here \(\Phi_{\rm SST}\) is the SI-normalized scalar potential associated
    with compressive, pressure-like, or boundary-induced swirl response.  In
    strictly incompressible regions, \(\Phi_{\rm SST}\) should be interpreted
    as an effective potential imposed by boundary constraints, topology, and
    coarse-grained pressure balance rather than by local volume compression.
    
    Equation~\eqref{eq:E_from_time_dependent_swirl} gives the SST analogue of
    Faraday induction: time-varying swirl produces an electric field.  Taking
    the curl gives
    \begin{equation}
        \nabla\times\mathbf{E}_{\rm SST}
        =
        -
        \frac{\partial}{\partial t}
        \left(
            \nabla\times\mathbf{A}_{\rm SI}
        \right)
        =
        -
        \frac{\partial\mathbf{B}_{\rm SST}}{\partial t},
        \label{eq:faraday_from_swirl}
    \end{equation}
    since \(\nabla\times\nabla\Phi_{\rm SST}=0\).  Thus the ordinary Faraday
    law is preserved exactly.

\subsection{Maxwell Equations as Effective Medium Equations}
    
    The correct research-track formulation is to retain Maxwell's equations
    in their macroscopic form while adding swirl-induced polarization and
    magnetization terms:
    \begin{align}
        \nabla\cdot\mathbf{D}_{\rm SST}
        &=
        \rho_{\rm free},
        \label{eq:SST_gauss_D}
        \\
        \nabla\cdot\mathbf{B}_{\rm SST}
        &=
        0,
        \label{eq:SST_gauss_B}
        \\
        \nabla\times\mathbf{E}_{\rm SST}
        &=
        -
        \frac{\partial\mathbf{B}_{\rm SST}}{\partial t},
        \label{eq:SST_faraday}
        \\
        \nabla\times\mathbf{H}_{\rm SST}
        &=
        \mathbf{J}_{\rm free}
        +
        \frac{\partial\mathbf{D}_{\rm SST}}{\partial t}.
        \label{eq:SST_ampere}
    \end{align}
    The constitutive relations are
    \begin{align}
        \mathbf{D}_{\rm SST}
        &=
        \varepsilon_0\mathbf{E}_{\rm SST}
        +
        \mathbf{P}_{\rm swirl},
        \label{eq:D_constitutive_swirl}
        \\
        \mathbf{H}_{\rm SST}
        &=
        \frac{1}{\mu_0}\mathbf{B}_{\rm SST}
        -
        \mathbf{M}_{\rm swirl}.
        \label{eq:H_constitutive_swirl}
    \end{align}
    Here \(\mathbf{P}_{\rm swirl}\) and \(\mathbf{M}_{\rm swirl}\) encode the
    coarse-grained response of bound swirl structures, knot cores, circulating
    R-phase modes, and boundary-induced polarization.
    
    This formulation has an important advantage: it does not violate the
    observed Maxwell structure.  It instead interprets charge and current as
    effective descriptions of deeper circulation and vorticity organization.

\subsection{Canonical Circulation-Scale Calibration and the QED Critical Field}
    
    At the SST canonical circulation scale the characteristic angular frequency scale is
    \begin{equation}
        \Omega_{\rm core}
        =
        \frac{\vnorm}{\rc}.
        \label{eq:omega_core}
    \end{equation}
    Using the canonical SST values
    \begin{equation}
        \vnorm
        =
        1.09384563\times10^{6}\ {\rm m\,s^{-1}},
        \qquad
        \rc
        =
        1.40897017\times10^{-15}\ {\rm m},
    \end{equation}
    one obtains
    \begin{equation}
        \Omega_{\rm core}
        =
        7.76344066\times10^{20}\ {\rm s^{-1}}.
        \label{eq:omega_core_numeric}
    \end{equation}
    The corresponding SST magnetic field scale is
    \begin{equation}
        B_{\rm core}^{\rm SST}
        =
        \frac{m_e}{e}\Omega_{\rm core}.
        \label{eq:Bcore_sst}
    \end{equation}
    Numerically,
    \begin{equation}
        \frac{m_e}{e}
        =
        5.68563010\times10^{-12}\ {\rm kg\,C^{-1}},
    \end{equation}
    and therefore
    \begin{equation}
        \boxed{
            B_{\rm core}^{\rm SST}
            =
            4.41400519\times10^{9}\ {\rm T}
        } .
        \label{eq:Bcore_numeric}
    \end{equation}
    
    This is not an arbitrary scale.  The standard QED critical magnetic field
    is
    \begin{equation}
        B_{\rm QED}
        =
        \frac{m_e^2c^2}{e\hbar}
        =
        4.41400522\times10^{9}\ {\rm T}.
        \label{eq:Bqed}
    \end{equation}
    Thus
    \begin{equation}
        \frac{B_{\rm core}^{\rm SST}}{B_{\rm QED}}
        =
        0.999999993.
        \label{eq:Bcore_Bqed_ratio}
    \end{equation}
    The equality follows analytically from the SST Compton--horn-radius relation
    \begin{equation}
        \lambda_c
        =
        \frac{2\pi c\rc}{\vnorm},
        \label{eq:compton_core_relation}
    \end{equation}
    together with the standard identity
    \begin{equation}
        \lambda_c
        =
        \frac{2\pi\hbar}{m_ec}.
    \end{equation}
    Equating these two expressions gives
    \begin{equation}
        \frac{\vnorm}{\rc}
        =
        \frac{m_ec^2}{\hbar}.
        \label{eq:omega_core_compton}
    \end{equation}
    Substitution into Eq.~\eqref{eq:Bcore_sst} yields
    \begin{equation}
        B_{\rm core}^{\rm SST}
        =
        \frac{m_e}{e}
        \frac{m_ec^2}{\hbar}
        =
        \frac{m_e^2c^2}{e\hbar}
        =
        B_{\rm QED}.
        \label{eq:Bcore_equals_Bqed}
    \end{equation}
    
    This is the strongest numerical result of this appendix: the SI-normalized
    vorticity bridge naturally reproduces the Schwinger/QED magnetic field
    scale at the SST core.
    
    The associated electric field scale is
    \begin{equation}
        E_{\rm core}^{\rm SST}
        =
        cB_{\rm core}^{\rm SST}
        =
        1.32328547\times10^{18}\ {\rm V\,m^{-1}},
        \label{eq:Ecore_numeric}
    \end{equation}
    which matches the usual Schwinger electric field
    \begin{equation}
        E_{\rm Schwinger}
        =
        \frac{m_e^2c^3}{e\hbar}.
        \label{eq:Schwinger_field}
    \end{equation}

\subsection{Corrected Maximum Force Relation}
    
    The same calibration fixes the correct factor in the SST maximum-force
    relation.  The expression
    \begin{equation}
        \frac{\vnorm\hbar}{\rc^2}
    \end{equation}
    has units of force, but it is twice the canonical force scale.  The
    consistent relation is
    \begin{equation}
        \boxed{
            \Fmax
            =
            \frac{\vnorm\hbar}{2\rc^2}
        } .
        \label{eq:Fmax_corrected}
    \end{equation}
    Using the canonical constants gives
    \begin{equation}
        \Fmax
        =
        29.0535098\ {\rm N},
    \end{equation}
    in agreement with the canonical value
    \begin{equation}
        F_{\max}^{\rm swirl}
        \simeq
        29.053507\ {\rm N}.
    \end{equation}
    The electron rest energy relation is then
    \begin{equation}
        \boxed{
            m_ec^2
            =
            2\Fmax \rc
        } .
        \label{eq:electron_energy_force_core}
    \end{equation}
    Equivalently,
    \begin{equation}
        m_e
        =
        \frac{2\Fmax\rc}{c^2}.
    \end{equation}
    Numerically,
    \begin{equation}
        \frac{2F_{\max}^{\rm swirl}\rc}{c^2}
        =
        9.10938\times10^{-31}\ {\rm kg},
    \end{equation}
    which matches the electron mass within the stated calibration precision.

\subsection{Photon and R-Phase Normal Modes}
    
    A second research-track result is that photon frequency should not be
    introduced through an ad hoc energy-density law of the form
    \begin{equation}
        u(r,\omega,T)
        \propto
        \frac{F_{\max}\omega^3}{\vnorm r^2}
        \frac{1}{\exp(\hbar\omega/k_BT)-1}.
        \label{eq:rejected_photon_density}
    \end{equation}
    This expression is not retained as a fundamental energy density because
    its bare prefactor does not have the units of energy per volume without
    additional geometric or constitutive factors.
    
    Instead, the research-track photon sector is formulated as a normal-mode
    condition on a closed or quasi-closed R-phase path:
    \begin{equation}
        \boxed{
            k_n
            =
            \frac{2\pi n}{L_{\rm eff}(K)},
            \qquad
            \omega_n
            =
            c_\gamma k_n
            =
            \frac{2\pi n c_\gamma}{L_{\rm eff}(K)}
        } .
        \label{eq:photon_frequency_topological_path}
    \end{equation}
    Here \(L_{\rm eff}(K)\) is the effective optical/swirl path length of the
    supporting configuration \(K\), and \(c_\gamma\) is the propagation speed
    of the R-phase excitation.  In the vacuum limit one expects
    \(c_\gamma=c\).  For a circular ring of radius \(R\),
    \begin{equation}
        L_{\rm eff}=2\pi R,
    \end{equation}
    so that
    \begin{equation}
        \boxed{
            \omega_n
            =
            \frac{n c_\gamma}{R}
        } .
        \label{eq:ring_photon_frequency}
    \end{equation}
    This gives a clean formulation of the statement:
    
    \begin{quote}
        Photon frequency depends on the supporting swirl-ring or path structure.
    \end{quote}
    
    For knotted, linked, or boundary-deformed paths, \(L_{\rm eff}(K)\) may
    include writhe, torsion, self-linking, and medium-response corrections:
    \begin{equation}
        L_{\rm eff}(K)
        =
        L(K)
        +
        \Delta L_{\rm writhe}
        +
        \Delta L_{\rm torsion}
        +
        \Delta L_{\rm boundary}
        +
        \Delta L_{\rm medium}.
        \label{eq:effective_path_length_decomposition}
    \end{equation}
    This places the photon sector in the same mathematical family as the SST
    delay-loop and mode-selection principles.

\subsection{Kelvin Impulse Bridge}
    
    For a thin circular swirl ring with circulation \(\Gamma\), major radius
    \(R\), and core radius \(a\), the classical ring scalings are
    \begin{align}
        E_{\rm ring}
        &=
        \frac{1}{2}\rhoF\Gamma^2 R
        \left[
            \ln\left(\frac{8R}{a}\right)-\alpha_{\rm ring}
        \right],
        \label{eq:ring_energy}
        \\
        P_{\rm ring}
        &=
        \rhoF\Gamma\pi R^2,
        \label{eq:ring_impulse}
        \\
        V_{\rm ring}
        &=
        \frac{\Gamma}{4\pi R}
        \left[
            \ln\left(\frac{8R}{a}\right)-\beta_{\rm ring}
        \right].
        \label{eq:ring_velocity}
    \end{align}
    These relations provide a bridge between Kelvin's impulse theory and an
    SST momentum assignment.  A research-track identification may be written
    as
    \begin{equation}
        \boxed{
            \hbar k
            \leftrightarrow
            \eta_I P_{\rm swirl}
        },
        \qquad
        P_{\rm swirl}
        =
        \rhoF\Gamma\pi R^2.
        \label{eq:hbar_k_kelvin_impulse_bridge}
    \end{equation}
    The coefficient \(\eta_I\) is not a free fit parameter in the final
    theory.  It must be determined from the same SI-normalization that fixes
    Eqs.~\eqref{eq:swirl_vector_potential_bridge}--\eqref{eq:Bcore_equals_Bqed},
    or from an independently specified R-phase normalization.
    
    This suggests a useful research direction:
    
    \begin{quote}
        Planck momentum may correspond to normalized Kelvin impulse of a circulating
        R-phase excitation.
    \end{quote}

\subsection{Topology, Helicity, and Reconnection}
    
    The ideal SST sector assumes conservation of circulation, linkage, and
    helicity in the absence of reconnection.  This corresponds to the
    Helmholtz--Kelvin regime:
    \begin{equation}
        \frac{D\Gamma}{Dt}=0,
        \qquad
        \frac{D\mathcal{H}}{Dt}=0,
        \qquad
        \mathcal{H}
        =
        \int
        \uswirl\cdot\boldsymbol{\omega}\,d^3x.
        \label{eq:helicity_conservation}
    \end{equation}
    In this ideal sector, knot class is preserved.
    
    However, laboratory and numerical studies of knotted and linked fluid
    structures show that reconnection can transfer helicity between linking,
    twist, and writhe.  Therefore, SST should distinguish two regimes:
    
    \begin{align}
        \text{Ideal sector:}
        \quad
        &K(t)=K(0),
        \qquad
        \Gamma=\text{constant},
        \qquad
        \mathcal{H}=\text{constant},
        \label{eq:ideal_topology_sector}
        \\
        \text{Transition sector:}
        \quad
        &K_i\rightarrow K_j,
        \qquad
        \Delta\mathcal{H}
        =
        \Delta\mathcal{L}
        +
        \Delta\mathcal{T}
        +
        \Delta\mathcal{W},
        \label{eq:transition_topology_sector}
    \end{align}
    where \(\mathcal{L}\), \(\mathcal{T}\), and \(\mathcal{W}\) denote
    linking, twist, and writhe contributions respectively.
    
    This is relevant for the SST interpretation of measurement, decay, and
    mode conversion.  A stable particle-like excitation belongs to the ideal
    sector.  A decay or transition corresponds to a controlled failure of
    ideal topological conservation, usually through reconnection or boundary
    interaction.

\subsection{Pressure Envelopes and Spherical Equilibrium}
    
    The electromagnetic bridge above should be kept distinct from the
    pressure-envelope sector of SST.  The pressure sector is governed by the
    Bernoulli-type swirl balance
    \begin{equation}
        p
        +
        \frac{1}{2}\rhoF\vnorm^2
        =
        p_0,
        \label{eq:bernoulli_swirl_pressure}
    \end{equation}
    or, for a perturbation,
    \begin{equation}
        \Delta p
        =
        -
        \frac{1}{2}\rhoF
        \Delta\left(\vnorm^2\right).
        \label{eq:pressure_defect}
    \end{equation}
    This pressure defect can generate attraction-like behavior in the
    gravitational/swirl-clock sector.  It should not be confused with the
    magnetic vorticity bridge
    \begin{equation}
        \mathbf{B}_{\rm SST}
        =
        \frac{m_e}{e}\boldsymbol{\omega}.
    \end{equation}
    
    Thus the research-track separation is:
    
    \begin{align}
        \text{EM-like sector:}
        \quad
        &\mathbf{A}_{\rm SI}
        =
        \frac{m_e}{e}\uswirl(\mathbf{x},t),
        \qquad
        \mathbf{B}_{\rm SST}
        =
        \frac{m_e}{e}\boldsymbol{\omega},
        \label{eq:em_sector_summary}
        \\
        \text{pressure/gravity sector:}
        \quad
        &\Delta p
        =
        -
        \frac{1}{2}\rhoF
        \Delta(\vnorm^2),
        \qquad
        d\tau_{\rm loc}/dt_\infty
        =
        \sqrt{1-\vnorm^2/c^2}.
        \label{eq:pressure_gravity_sector_summary}
    \end{align}

\subsection{Negative-Temperature and Entropy Analogy}
    
    A further research-track idea is that bounded swirl ensembles may possess
    entropy-ordering behavior analogous to negative-temperature vortex states.
    This is not required for the electromagnetic normalization, but it may be
    useful in the classification of high-energy knot ensembles.
    
    A possible entropy functional is
    \begin{equation}
        S_{\rm swirl}
        =
        k_B\ln\Omega_{\rm topo}(E,\Gamma,\mathcal{H}),
        \label{eq:swirl_entropy}
    \end{equation}
    where \(\Omega_{\rm topo}\) counts accessible configurations at fixed
    energy, circulation, and helicity.  A negative-temperature-like regime
    would satisfy
    \begin{equation}
        \frac{\partial S_{\rm swirl}}{\partial E}
        <
        0.
        \label{eq:negative_temperature_condition}
    \end{equation}
    Such a regime would describe increasing order with increasing energy,
    for example through clustering, alignment, or coherent large-scale swirl
    organization.  This idea should remain in the research track until it is
    connected to a concrete SST Hamiltonian and a well-defined finite phase
    space.

\subsection{Research Claims Retained}
    
    The following claims are retained for research-track development:
    
    \begin{enumerate}
        \item Magnetic field strength can be SI-normalized as vorticity:
        \[
            \mathbf{B}_{\rm SST}
            =
            \frac{m_e}{e}\boldsymbol{\omega}.
        \]
        
        \item The vector potential is the SI-normalized swirl velocity:
        \[
            \mathbf{A}_{\rm SI}
            =
            \frac{m_e}{e}\uswirl(\mathbf{x},t).
        \]
        
        \item The core-scale SST magnetic field equals the QED critical field:
        \[
            B_{\rm core}^{\rm SST}
            =
            B_{\rm QED}
            =
            \frac{m_e^2c^2}{e\hbar}.
        \]
        
        \item The associated electric field equals the Schwinger scale:
        \[
            E_{\rm core}^{\rm SST}
            =
            cB_{\rm core}^{\rm SST}
            =
            \frac{m_e^2c^3}{e\hbar}.
        \]
        
        \item Photon frequency should be modeled by normal-mode closure:
        \[
            \omega_n
            =
            \frac{2\pi n c_\gamma}{L_{\rm eff}(K)}.
        \]
        
        \item Kelvin impulse may provide the bridge to photon momentum:
        \[
            \hbar k
            \leftrightarrow
            \eta_I\rhoF\Gamma\pi R^2.
        \]
        
        \item Reconnection belongs to a separate transition sector and should
        not be mixed with ideal knot conservation.
    \end{enumerate}

\subsection{Claims Not Retained as Canonical}
    
    The following statements are not retained as canonical equations:
    
    \begin{enumerate}
        \item Direct insertion of uncalibrated vorticity terms into Maxwell
        source equations.
        
        \item The statement that charge is unnecessary in all electromagnetic
        phenomena.  In this appendix charge remains the effective macroscopic
        source variable.
        
        \item The claim that Podkletnov-type gravitational shielding is
        established evidence for SST.  Such material may be mentioned only as
        speculative historical motivation, not as support for the canon.
        
        \item Photon energy-density laws containing
        \(F_{\max}\omega^3/(\vnorm r^2)\) without a complete dimensional and
        geometric derivation.
    \end{enumerate}

\subsection{Minimal Experimental Tests}
    
    The vorticity-normalization bridge suggests several falsifiable directions.
    
    \paragraph{1. Structured plasmonic or polaritonic media.}
        If electromagnetic fields are coupled to structured circulation-like
        degrees of freedom, then anisotropic or hyperbolic media should show
        geometry-dependent emission and propagation consistent with an effective
        path length \(L_{\rm eff}(K)\).
    
    \paragraph{2. Superfluid or quantum-fluid circulation.}
        Quantized circulation structures should generate measurable effective
        magnetization only when the system contains charged, spinful, or
        electromagnetically coupled degrees of freedom.  A neutral inviscid fluid
        alone should not emit ordinary electromagnetic fields unless an additional
        coupling channel exists.
    
    \paragraph{3. Ring-resonator delay tests.}
        The photon/R-phase normal-mode law predicts that frequency selection
        should scale as
        \[
            \omega_n
            \propto
            \frac{1}{L_{\rm eff}}.
        \]
        Boundary deformation, writhe-like path changes, or torsional loading
        should shift the selected resonances.
    
    \paragraph{4. Core-scale extrapolation.}
        The equality
        \[
            B_{\rm core}^{\rm SST}=B_{\rm QED}
        \]
        is not directly reachable in ordinary laboratory magnetic fields, but it
        can be tested indirectly by checking whether the same normalization
        reproduces Compton, Schwinger, and electron-rest-energy scales without
        additional tuning.

\subsection{Numerical Verification Table}

Using the canonical SST constants
\[
    \vnorm=1.09384563\times10^6\ {\rm m\,s^{-1}},
    \qquad
    \rc=1.40897017\times10^{-15}\ {\rm m},
\]
one obtains:

\begin{center}
    \begin{tabular}{lll}
        \hline
        Quantity & Expression & Numerical value \\
        \hline
        Core angular scale &
        \(\Omega_{\rm core}=\vnorm/\rc\) &
        \(7.76344066\times10^{20}\ {\rm s^{-1}}\) \\
        
        Mass-charge ratio &
        \(m_e/e\) &
        \(5.68563010\times10^{-12}\ {\rm kg\,C^{-1}}\) \\
        
        SST core magnetic field &
        \((m_e/e)\Omega_{\rm core}\) &
        \(4.41400519\times10^9\ {\rm T}\) \\
        
        QED critical magnetic field &
        \(m_e^2c^2/(e\hbar)\) &
        \(4.41400522\times10^9\ {\rm T}\) \\
        
        Ratio &
        \(B_{\rm core}^{\rm SST}/B_{\rm QED}\) &
        \(0.999999993\) \\
        
        SST electric field scale &
        \(cB_{\rm core}^{\rm SST}\) &
        \(1.32328547\times10^{18}\ {\rm V\,m^{-1}}\) \\
        
        Circulation quantum &
        \(2\pi\rc\vnorm\) &
        \(9.68361920\times10^{-9}\ {\rm m^2\,s^{-1}}\) \\
        
        Corrected force ceiling &
        \(\vnorm\hbar/(2\rc^2)\) &
        \(29.0535098\ {\rm N}\) \\
        
        Electron mass from force &
        \(2F_{\max}^{\rm swirl}\rc/c^2\) &
        \(9.10938\times10^{-31}\ {\rm kg}\) \\
        \hline
    \end{tabular}
\end{center}

\subsection{Conclusion}

The strongest result of this appendix is the SI-normalized vorticity
bridge
\[
    \mathbf{A}_{\rm SI}
    =
    \frac{m_e}{e}\uswirl(\mathbf{x},t),
    \qquad
    \mathbf{B}_{\rm SST}
    =
    \frac{m_e}{e}\boldsymbol{\omega}.
\]
It preserves Maxwell's equations as effective field equations, avoids
dimensionally invalid source terms, and reproduces the QED critical field
at the SST core scale.  For this reason it is suitable for inclusion in
the research-track sections of CANON-v0.8.x.

The proposed canonical status is:

\begin{center}
    \begin{tabular}{ll}
        \hline
        Result & Status \\
        \hline
        \(\mathbf{A}_{\rm SI}=(m_e/e)\uswirl(\mathbf{x},t)\) & Research Track, strong candidate \\
        \(\mathbf{B}_{\rm SST}=(m_e/e)\boldsymbol{\omega}\) & Research Track, strong candidate \\
        \(B_{\rm core}^{\rm SST}=B_{\rm QED}\) & Numerically verified calibration result \\
        Direct vorticity source terms in Maxwell equations & Rejected / not canonical \\
        Photon frequency from \(L_{\rm eff}(K)\) & Research Track, compatible with delay-loop sector \\
        Kelvin impulse to \(\hbar k\) bridge & Research Track, unresolved normalization \\
        Negative-temperature swirl entropy & Speculative Research Track \\
        \hline
    \end{tabular}
\end{center}
% ============================================================
% Research Track: Triadic gravity-response diagnostics
% Derived from the EM--gravity sector-split conversation track.
% ============================================================

\subsection{Research Track: Flame-Lock, Photonless-Caustic, and Elastic-Shell Diagnostics}
\label{sec:rt_flame_photonless_shell}

\paragraph{Status.}
\textbf{[RESEARCH TRACK / DIAGNOSTIC ONLY].}
This section records a three-channel diagnostic motivated by anomalous-field
witness reports. The reports are not used as evidence. They serve only as a
catalogue of separable experimental signatures compatible with the canonical
sector-projection law of Section~\ref{subsec:canonical_em_gravity_closure}.

\subsection{Mode I: flame-lock / plume stabilization}

A flame or thermal plume may be stabilized without significant optical ray
deflection if the dominant SST response is local impedance or damping rather
than photon-index curvature. A minimal diagnostic equation is
\begin{align}
    \frac{\partial\mathbf{u}_{\mathrm{plume}}}{\partial t}
    +
    (\mathbf{u}_{\mathrm{plume}}\cdot\nabla)\mathbf{u}_{\mathrm{plume}}
    =
    -\frac{1}{\rho_{\mathrm{plume}}}\nabla p
    +
    \nu\nabla^2\mathbf{u}_{\mathrm{plume}}
    -
    \Gamma_{\swirlarrow}\mathbf{u}_{\mathrm{plume}}.
    \label{eq:rt_flame_lock_plume_equation}
\end{align}
The diagnostic signature is
\begin{align}
    P_{\mathrm{plume}}(f)\downarrow,
    \qquad
    \Delta\phi_\gamma\approx0,
    \qquad
    \nabla_\perp n_\gamma\approx0.
    \label{eq:rt_flame_lock_signature}
\end{align}
This channel is a boundary/impedance response, not a standalone proof of a
gravitational bulk force.

\textbf{Primary orthodox confounders:}
ion wind, acoustic forcing, thermal convection suppression, electric-field
stabilization of charged flame species, camera exposure artifacts, and ordinary
airflow.

\subsection{Mode II: photonless sphere / optical caustic}

A dark or photon-depleted region is interpreted, if real, as an optical caustic,
shadow, or phase-gradient effect, not as a literal astrophysical black hole. The
minimal ray equation is
\begin{align}
    \frac{d}{ds}
    \left(
    n_\gamma\frac{d\mathbf{x}}{ds}
    \right)
    =
    \nabla n_\gamma,
    \qquad
    n_\gamma=\SwirlClock^{-1}.
    \label{eq:rt_photonless_caustic_ray}
\end{align}
The diagnostic signature is
\begin{align}
    \nabla_\perp n_\gamma\neq0,
    \qquad
    \Delta\phi_\gamma
    =
    \frac{2\pi}{\lambda}
    \int(n_\gamma-1)\,ds
    \neq0,
    \qquad
    \mathbf{f}_{\mathrm{matter}}\approx0.
    \label{eq:rt_photonless_caustic_signature}
\end{align}
Macroscopic photon depletion or black-ball-like appearance over laboratory path
lengths requires a resonant enhancement factor:
\begin{align}
    n_\gamma-1
    =
    Q_\gamma
    \frac{\lVert\mathbf{u}_{\swirlarrow}\rVert^2}{2c^2},
    \qquad
    Q_\gamma\gg1.
    \label{eq:rt_optical_gain_factor}
\end{align}
The gain \(Q_\gamma\) is apparatus-specific and is not a canonical constant.

\textbf{Primary orthodox confounders:}
thermal lensing, plasma absorption, smoke or aerosol scattering, detector
saturation, schlieren artifacts, camera auto-exposure, and ordinary optical
shadowing.

\subsection{Mode III: elastic shell / boundary-stress response}

A localized repulsive shell is represented as a boundary-stress potential,
\begin{align}
    U_{\mathrm{shell}}(r)
    =
    U_0
    \exp\left[
    -\frac{(r-R_s)^2}{2\sigma_s^2}
    \right],
    \label{eq:rt_shell_potential}
\end{align}
with radial force
\begin{align}
    F_r(r)
    =
    -\frac{dU_{\mathrm{shell}}}{dr}
    =
    \frac{U_0(r-R_s)}{\sigma_s^2}
    \exp\left[
    -\frac{(r-R_s)^2}{2\sigma_s^2}
    \right].
    \label{eq:rt_shell_force}
\end{align}
The diagnostic signature is
\begin{align}
    \nabla\cdot\boldsymbol{\sigma}^{\mathrm{swirl}}\neq0,
    \qquad
    \Delta\phi_\gamma\approx0
    \quad
    \text{or independently measurable}.
    \label{eq:rt_shell_signature}
\end{align}
This mode is a stress/boundary response. It must not be rebranded as gravity
shielding unless the bulk acceleration sector is independently measured.

\textbf{Primary orthodox confounders:}
electrostatic repulsion, magnetic force, acoustic radiation pressure, airflow,
vibration, thermal expansion, hidden mechanical coupling, and contact charging.

\subsection{Required measurement vector}

A valid test must record the full sectoral observable set
\begin{align}
    \boxed{
    \left\{
    F(t),
    \Delta\phi_\gamma(t),
    \frac{\Delta f}{f}(t),
    P_{\mathrm{plume}}(f,t),
    \mathbf{E}(t),
    \mathbf{B}(t),
    T(t),
    p_{\mathrm{air}}(t)
    \right\}.
    }
    \label{eq:rt_triadic_measurement_vector}
\end{align}
Any reported force must first be decomposed as
\begin{align}
    \mathbf{F}_{\mathrm{meas}}
    =
    \mathbf{F}_{\mathrm{EM}}
    +
    \mathbf{F}_{\mathrm{EHD}}
    +
    \mathbf{F}_{\mathrm{thermal}}
    +
    \mathbf{F}_{\mathrm{acoustic}}
    +
    \mathbf{F}_{\mathrm{mechanical}}
    +
    \mathbf{F}_{\omega}
    +
    \mathbf{F}_{\mathrm{res}}.
    \label{eq:rt_triadic_force_decomposition}
\end{align}
Only \(\mathbf{F}_{\mathrm{res}}\), after uncertainty propagation and null tests,
is admissible as an SST research-track candidate.

\subsection{Null hierarchy}

The minimum null hierarchy is:
\begin{enumerate}
    \item reproduce the measurement with all high-voltage and high-current channels disabled;
    \item randomize phase order at fixed RMS power;
    \item replace the structured geometry with a symmetry-preserving dummy load of matched impedance;
    \item repeat in still air, shielded enclosure, and reduced pressure where safe;
    \item compare optical, force, and plume observables under the same drive waveform.
\end{enumerate}
A positive SST candidate requires a correlated sectoral signature that survives
these nulls and scales with the predicted mode-selection or stress-closure
parameter.


% ============================================================
% Research-track patch: no-monopole transport-stress audit
% ============================================================

\subsection{Research Track: No-Monopole Audit of the Transport-Stress Gravity Candidate}
\label{sec:rt_no_monopole_transport_stress_audit}

\paragraph{Status.}
\textbf{[RESEARCH TRACK / DERIVED CONSTRAINT].}
This section records a constraint on the SST-73 transport-stress gravity
candidate.  It does not modify the canon-level weak-field closure
\(\nabla^2\Phi_{\swirlarrow}=4\pi G_{\mathrm{swirl}}\rhoM\).  It shows that a
double-divergence stress source belongs to the local pressure/stress channel
unless an additional non-divergence density closure or boundary source is
specified.

\subsection{Pressure equivalence lemma}
For an incompressible inviscid medium with constant \(\rhoF\),
\begin{align}
    \partial_i u_i=0,
\end{align}
and no external body force, the Euler equation is
\begin{align}
    \rhoF\left(\partial_t u_i + u_j\partial_j u_i\right)
    =
    -\partial_i p .
\end{align}
Taking the divergence and using \(\partial_i u_i=0\) gives
\begin{align}
    \partial_i\partial_j\left(\rhoF u_i u_j\right)
    =
    -\nabla^2p .
    \label{eq:rt_pressure_equivalence_euler_identity}
\end{align}
Therefore the SST-73 transport-stress candidate
\begin{align}
    \nabla^2\Phi_{\mathrm{tr}}
    =
    \lambda\,\partial_i\partial_j\left(\rhoF u_i u_j\right)
\end{align}
is equivalent to
\begin{align}
    \nabla^2\left(\Phi_{\mathrm{tr}}+\lambda p\right)=0 .
\end{align}
For decaying exterior boundary conditions,
\begin{align}
    \boxed{
    \Phi_{\mathrm{tr}}
    =
    -\lambda\left(p-p_\infty\right)
    }
    \label{eq:rt_transport_pressure_equivalence}
\end{align}
up to an irrelevant harmonic constant.  Candidate C is therefore not an
independent long-range gravity mechanism in the smooth incompressible bulk; it
is the Euler pressure channel written in Poisson form.

\subsection{No-monopole theorem for smooth compact stress sources}
Let
\begin{align}
    S(\mathbf{x})
    =
    \partial_i\partial_j\Sigma_{ij},
    \qquad
    \Sigma_{ij}=\rhoF u_i u_j,
\end{align}
with \(\Sigma_{ij}\) smooth and compactly supported, or decaying sufficiently
fast at infinity.  The monopole moment is
\begin{align}
    Q_0
    =
    \int_{\mathbb{R}^3}S(\mathbf{x})\,dV
    =
    \oint_{\partial\mathbb{R}^3}\partial_j\Sigma_{ij}\,dA_i
    =0 .
\end{align}
The dipole moment also vanishes after integration by parts:
\begin{align}
    Q_k
    =
    \int_{\mathbb{R}^3}x_k S(\mathbf{x})\,dV
    =0 .
\end{align}
Hence the first possible nonzero exterior multipole is quadrupolar:
\begin{align}
    \Phi_{\mathrm{tr}}(r)=O(r^{-3}),
    \qquad
    \lVert\nabla\Phi_{\mathrm{tr}}\rVert=O(r^{-4}).
\end{align}
This cannot supply a Newtonian \(1/r^2\) acceleration by itself.

\subsection{Boundary and singular-core exception}
The no-monopole result assumes a smooth regularized field on all of
\(\mathbb{R}^3\).  If vortex tubes are excised, singular cores are retained,
or material boundaries/shells are present, the surface term
\begin{align}
    \oint_{\partial V}\partial_j\Sigma_{ij}\,dA_i
\end{align}
need not vanish.  Such terms are boundary-stress or shell-response physics
\(\mathcal{R}_{\Pi}\), not bulk Poisson gravity.  A nonzero Newtonian monopole
therefore requires either a genuine scalar density source, such as
\(\rhoM=\rhoE/c^2\), or an explicitly modeled boundary source.

\subsection{Equivalence-principle consequence of the density source}
If the gravitational source density is the same effective mass density that
defines inertial mass,
\begin{align}
    \rhoM^{\mathrm{eff}}
    =
    \frac{\rhoE}{c^2},
    \qquad
    \rhoE=\frac{1}{2}\rhoF\lVert\uswirl\rVert^2,
\end{align}
then
\begin{align}
    M_{\mathrm{grav}}
    =
    \int\rhoM^{\mathrm{eff}}\,dV
    =
    \frac{1}{c^2}\int\rhoE\,dV
    =
    M_{\mathrm{inert}} .
\end{align}
For full topological SST states this statement must be applied to the same
coarse-grained mass functional used for inertial mass, including topology,
geometric gate, and clock-impedance factors.  The numerical value of
\(G_{\mathrm{swirl}}\) remains \textbf{[CALIBRATED]} until derived
independently of orthodox \(G\) or \(t_p\).

\subsection{Minimal numerical tests}
A verification script for this sector must include:
\begin{enumerate}
    \item a smooth compact divergence-free test field for which
    \(\int S\,dV\approx0\) and \(\int x_kS\,dV\approx0\);
    \item a far-field multipole fit verifying
    \(\Phi_{\mathrm{tr}}\sim r^{-3}\), not \(r^{-1}\);
    \item a pressure-equivalence test verifying
    \(\Phi_{\mathrm{tr}}+\lambda p\) is harmonic/constant under the stated
    boundary conditions;
    \item a separate density-source Poisson test verifying
    \(\Phi\sim -G_{\mathrm{swirl}}M/r\) and
    \(\lVert\mathbf{g}\rVert\sim r^{-2}\) for compact positive
    \(\rhoM^{\mathrm{eff}}\).
\end{enumerate}

\subsection{Research Track: Atomic Gravity Closure and Condensate Coherence}
\label{sec:rt_atomic_gravity_closure_condensate}

\paragraph{Status.}
\textbf{[CANON-CANDIDATE / RESEARCH BRIDGE].}
This section records the electron-envelope and condensate-coherence clue for
future gravity-sector tests. It does not claim gravitational shielding.

\subsection{Electron envelope as atomic closure layer}

In a neutral atom, the electron envelope cancels the far electromagnetic charge
sector but does not cancel positive swirl-energy density. Therefore the electron
is not the dominant mass reservoir, but it is the outer phase and linking
boundary through which the neutral atom closes its electromagnetic projection:
\begin{align}
    Q_{\mathrm{atom}}=0,
    \qquad
    \rhoM
    =
    \frac{\rhoE}{c^2}
    >0.
    \label{eq:rt_atomic_charge_mass_split}
\end{align}
The canon-safe interpretation is
\begin{quote}
\textit{
The nucleus supplies most mass-energy, while the electron envelope supplies
the neutral phase boundary and long-range matching layer.}
\end{quote}

\subsection{Condensate coherence extension}

For coherent electronic or atomic condensates, a macroscopic order parameter
\begin{align}
    \Psi_{\mathrm{cond}}(\mathbf{x})
    =
    \sqrt{n_s(\mathbf{x})}\,e^{i\theta(\mathbf{x})}
    \label{eq:rt_condensate_order_parameter}
\end{align}
may coherently reweight boundary-stress and optical/clock projections:
\begin{align}
    \boldsymbol{\sigma}^{\mathrm{swirl}}
    &\rightarrow
    \boldsymbol{\sigma}^{\mathrm{swirl}}_{\mathrm{coh}}[n_s,\theta,\nabla\theta],
    \\
    n_\gamma
    &\rightarrow
    n_{\gamma,\mathrm{coh}}.
    \label{eq:rt_condensate_projection_reweighting}
\end{align}
The proposed clue from Meissner/BEC physics is therefore not mass cancellation,
but collective phase closure: many microscopic carriers can impose one
macroscopic boundary condition.

\textbf{[CRITICAL NOTE].}
This does not imply gravitational shielding or mass cancellation. The testable
claim is only a possible phase-dependent change in boundary-stress, optical
phase, or clock response across a coherence transition such as \(T_c\).

\subsection{Condensate test signature}

A minimally admissible superconducting or BEC test compares the same apparatus
above and below the coherence transition:
\begin{align}
    \Delta\mathcal{O}_{\mathrm{coh}}
    =
    \mathcal{O}(T<T_c)-\mathcal{O}(T>T_c),
    \qquad
    \mathcal{O}
    \in
    \left\{
    F,
    \Delta\phi_\gamma,
    \Delta f/f,
    \Delta\Phi
    \right\}.
    \label{eq:rt_condensate_test_signature}
\end{align}
A positive result must track coherence and vanish under matched non-coherent
thermal, electromagnetic, and mechanical controls.

\subsection{Relativity Emergence Ladder in SST}
\label{sec:rt_relativity_emergence_ladder}

\subsection{Purpose and status}

This section records the precise epistemic status of special-relativistic and
general-relativistic structures inside Swirl--String Theory (SST). The goal is
not to overclaim that general relativity has been derived from incompressible
Euler dynamics. The goal is to separate:
\begin{enumerate}
    \item kinematic consequences of a common propagation cone;
    \item clock and metric bridge structures;
    \item weak-field gravitational closures;
    \item the still-open problem of deriving the Einstein field equations.
\end{enumerate}

\textbf{Status summary.} SST presently supports a conditional reconstruction of
relativistic kinematics for its excitation sector, a conditional weak
equivalence principle, and a conditional Newtonian limit. The full
Einstein--Hilbert dynamics remain an open research target.

\subsection{Primitive structure}

The minimal SST substrate is taken to be an incompressible, inviscid medium on a
preferred foliation,
\begin{align}
    \mathcal{M}_{\mathrm{SST}}
    =
    \mathbb{R}^3 \times \mathbb{R},
    \qquad
    \nabla\cdot \mathbf{v}=0 ,
\end{align}
with structured vorticity, quantized circulation, a canonical transverse
propagation speed \(c\), and a local Swirl--Clock factor
\begin{align}
    S_{(t)}(x)
    =
    \sqrt{1-\frac{u(x)^2}{c^2}} .
    \label{eq:rt_swirl_clock}
\end{align}

\textbf{[ORTHODOX]} The functional form of
Eq.~\eqref{eq:rt_swirl_clock} is the standard Lorentz time-dilation factor.

\textbf{[SST INTERPRETATION]} The velocity \(u(x)\) is interpreted as a local
medium-kinematic state rather than as an abstract spacetime postulate.

\textbf{[CRITICAL NOTE]} If \(c\) is inserted as a primitive signal speed, then
the existence of a limiting cone is not itself derived. What may be derived is
the operational Lorentz kinematics of clocks and rods whose internal dynamics
are mediated only by that cone.

\subsection{Light-clock derivation of the Lorentz factor}

Consider a clock whose tick is defined by a transverse signal moving at speed
\(c\) in the substrate frame. If the clock moves with speed \(u\) relative to
the substrate, then during one half-cycle
\begin{align}
    c^2 \Delta t^2
    =
    u^2 \Delta t^2 + c^2 \Delta \tau^2 .
\end{align}
Therefore
\begin{align}
    \Delta \tau
    =
    \Delta t
    \sqrt{1-\frac{u^2}{c^2}} .
    \label{eq:rt_light_clock_gamma}
\end{align}

\textbf{[DERIVED, conditional]} Equation~\eqref{eq:rt_light_clock_gamma}
derives Lorentz time dilation for clocks whose internal communication is
entirely mediated by substrate excitations with speed \(c\) \cite{Einstein1905}.

\textbf{[OPEN CANON GAP]} Universal Lorentz kinematics for all matter species
requires the stronger monometricity condition below.

\subsection{Monometricity theorem target}

Let \(A\) label low-energy SST excitation species. The principal symbol of the
linearized equations for species \(A\) defines an effective characteristic
surface
\begin{align}
    P_A(k)=0 .
\end{align}
Relativistic universality requires that, at low energy,
\begin{align}
    P_A(k)=0
    \quad\Longleftrightarrow\quad
    g_{\mathrm{eff}}^{\mu\nu} k_\mu k_\nu=0
    \qquad
    \text{for all } A ,
    \label{eq:rt_monometricity}
\end{align}
up to irrelevant conformal factors:
\begin{align}
    g^{\mu\nu}_{\mathrm{eff},A}
    =
    \Omega_A^2(x)\,
    g^{\mu\nu}_{\mathrm{eff}} .
\end{align}

\textbf{[OPEN CANON GAP: MONOMETRICITY]} SST currently contains at least two
distinguished speeds,
\begin{align}
    c,
    \qquad
    \mathbf{v}_{\!\boldsymbol{\circlearrowleft}}
    =
    \frac{\alpha c}{2},
\end{align}
so exact low-energy Lorentz universality is not automatic. The required theorem
is that matter knots, gauge/torsion pulses, and clock excitations share the
same effective cone in the infrared. In condensed-matter language this is the
Fermi-point universality problem \cite{Volovik2003}.

\textbf{Falsifier.} Failure of Eq.~\eqref{eq:rt_monometricity} predicts
species-dependent Lorentz violation and preferred-frame effects.

\subsection{Analogue-metric bridge}

For a hyperbolic perturbation sector with common characteristic speed \(c\),
the linearized excitation equation can be written schematically as
\begin{align}
    \frac{1}{\sqrt{-g_{\mathrm{eff}}}}
    \partial_\mu
    \left(
        \sqrt{-g_{\mathrm{eff}}}\,
        g_{\mathrm{eff}}^{\mu\nu}
        \partial_\nu \phi
    \right)
    =
    0 .
    \label{eq:rt_analogue_metric}
\end{align}

\textbf{[ORTHODOX]} Equations of the form
\eqref{eq:rt_analogue_metric} are standard in analogue-gravity systems
\cite{Unruh1981,BarceloLiberatiVisser2011}.

\textbf{[SST STATUS: CONDITIONAL BRIDGE]} In SST this bridge applies only to the
explicit hyperbolic excitation sector. It does not follow from incompressibility
alone, because incompressible Euler pressure is a constraint field rather than
an ordinary compressional sound mode. The correct conditional statement is: if
the torsion/transverse excitation sector has a hyperbolic principal operator
with speed \(c\), then an effective Lorentz cone follows.

\textbf{[CRITICAL NOTE]} Analogue gravity reproduces curved-spacetime
kinematics for perturbations. It does not, by itself, derive the Einstein field
equations.

\subsection{Weak equivalence principle}

There are two independent conditional routes to the weak equivalence principle.

\paragraph{Energy-source route.}
If both inertial mass and gravitational source strength are determined by the
same localized energy functional,
\begin{align}
    m_i(T)c^2 = E_T,
    \qquad
    m_g(T)c^2 = E_T,
\end{align}
then
\begin{align}
    m_i(T)=m_g(T).
\end{align}

\textbf{[DERIVED, conditional]} This route is valid only if the gravitational
clock/pressure field couples universally to the total localized SST energy
functional.

\paragraph{Metric-coupling route.}
If matter is minimally coupled to a Lorentzian metric and
\(\nabla_\mu T^{\mu\nu}=0\), then the point-particle limit follows geodesic
motion,
\begin{align}
    u^\mu \nabla_\mu u^\nu =0 .
\end{align}

\textbf{[ORTHODOX CONDITIONAL THEOREM]} This is the standard GR result. In SST
it is not yet a primitive-substrate derivation, because the Lorentzian metric
and minimal coupling are assumed at the effective-field-theory level.

\subsection{Newtonian limit}

A Newtonian limit requires a scalar potential satisfying
\begin{align}
    \nabla^2 \Phi
    =
    4\pi G_{\mathrm{eff}}\rho_{\!m}.
    \label{eq:rt_poisson_relativity}
\end{align}
SST can motivate \(\rho_{\!m}=\rho_{\!E}/c^2\) from the local energy accounting,
but Eq.~\eqref{eq:rt_poisson_relativity} remains a closure law unless derived from the
coarse-grained vortex ensemble.

\textbf{[CONDITIONAL]} If Eq.~\eqref{eq:rt_poisson_relativity} is admitted as the
long-wavelength closure, the inverse-square force and Newtonian limit follow.

\textbf{[OPEN CANON GAP]} Derive Eq.~\eqref{eq:rt_poisson_relativity} directly from the
statistical mechanics or coarse-grained pressure response of many swirl strings.

\subsection{Research Track: Compton--Horn Balance Angle}
\label{sec:rt_compton_horn_balance_angle}

\textbf{[RESEARCH-TRACK GEOMETRIC DIAGNOSTIC].}
This subsection records a dimensionless visualization of the calibrated
Compton--horn closure used in the core canon.  It does not introduce a new
gravity law and it must not be used as an independent derivation of
Newton's constant.

Define \(\phi_G\) by
\begin{align}
    G
    =
    \frac{\vchar c^3t_p^2}{M_e\rc}\,\tan^2\phi_G .
    \label{eq:rt_balance_angle_definition}
\end{align}
Equivalently,
\begin{align}
    \tan^2\phi_G
    =
    \frac{GM_e\rc}{\vchar c^3t_p^2}.
    \label{eq:rt_balance_angle_value}
\end{align}
Using the v0.8.18 calibration chain
\begin{align}
    \rc=\frac{\vchar}{\omega_c},
    \qquad
    \omega_c=\frac{M_ec^2}{\hbar},
    \qquad
    t_p^2=\frac{\hbar G}{c^5},
\end{align}
one obtains
\begin{align}
    \tan^2\phi_G
    =
    \frac{M_ec^2}{\hbar\omega_c}
    =1,
    \qquad
    \phi_G=\frac{\pi}{4}.
    \label{eq:rt_balance_angle_pi_over_four}
\end{align}
Thus the \(45^\circ\) value means that the calibrated electron benchmark
weights radial horn confinement and tangential swirl transport equally in this
particular dimensionless projection.  It does \emph{not} mean that gravity acts
at a spatial angle of \(45^\circ\), and it does \emph{not} add a new coupling
constant.

\textbf{[CRITICAL NOTE]} Because \(t_p^2=\hbar G/c^5\) already contains \(G\),
Eq.~\eqref{eq:rt_balance_angle_pi_over_four} is a closure diagnostic, not a
prediction of \(G\) from non-gravitational inputs.  A future promotion would
require deriving the effective long-range coupling or the Planck-time scale
from SST medium dynamics without importing orthodox \(G\).

\subsubsection{Rejected microscopic Kelvin-angle ansatz}
The dimensional ansatz
\begin{align}
    \tan\phi
    =
    \sqrt{
    \frac{I^{3/2}}{N\mu K^{1/2}\sqrt{\pi}}
    }
    \label{eq:rt_rejected_kelvin_angle_ansatz}
\end{align}
is not canon-ready.  In SI units \([I]=\mathrm{kg\,m^2}\), while the present
canon uses \([K]=\mathrm{kg\,m^{-3}\,s}\) for the density--rotation bridge.
Unless \(\mu\) and \(K\) are replaced by explicitly normalized quantities,
Eq.~\eqref{eq:rt_rejected_kelvin_angle_ansatz} is not dimensionless, whereas
\(\tan\phi\) must be dimensionless.

A permissible research-track replacement must use only dimensionless knot
functionals:
\begin{align}
    \tan^2\phi_K
    =
    \frac{\mathcal{I}_K^{3/2}}
    {N_K\mathcal{M}_K\mathcal{K}_K^{1/2}\sqrt{\pi}},
    \label{eq:rt_dimensionless_kelvin_angle}
\end{align}
with, for example,
\begin{align}
    \mathcal{I}_K=\frac{I_K}{M_K\rc^2},
    \qquad
    \mathcal{K}_K=\frac{K_K}{K_0},
    \qquad
    \mathcal{M}_K=\frac{\mu_K}{\mu_0^\ast}.
\end{align}
For the electron/trefoil benchmark the admissible normalization condition is
\begin{align}
    \tan^2\phi_{3_1}=1,
\end{align}
but this is a benchmark closure condition, not yet a microscopic theorem.

\subsection{Einstein dynamics: three candidate routes}

\paragraph{Route I: thermodynamic gravity.}
Jacobson's route derives the Einstein equation as an equation of state from
\(\delta Q=T\,dS\) applied to local Rindler horizons, assuming an entropy-area
law and the Unruh temperature \cite{Jacobson1995}. In SST this requires:
\begin{align}
    S_{\mathrm{defect}}
    \propto
    A,
    \qquad
    T_{\mathrm{SST}}
    =
    \frac{\hbar a}{2\pi c k_B}.
\end{align}

\textbf{[RESEARCH TRACK]} This is the most promising route because it connects
directly to horizon, boundary, and holographic SST sectors.


\paragraph{Source-paper integration for Route I.}
The thermodynamic-gravity route should be linked explicitly to three source
papers.  SST-63 supplies the constraint-holographic boundary reconstruction
layer.  SST-23 supplies the accelerated torsion / Unruh-response candidate.
SST-56 supplies the superfluid topology diagnostics---Hall angle, writhe
barrier, and lifetime hierarchy---needed to test whether a line-tangle or
fabric sector has stable protected response.  Thus the route is not a single
paper claim but a three-asset programme:
\begin{align}
    \boxed{
    \mathrm{SST\mbox{-}63}
    +
    \mathrm{SST\mbox{-}23}
    +
    \mathrm{SST\mbox{-}56}
    \;\longrightarrow\;
    \left(S=\eta A,\;T_{\mathrm{SST}},\;\delta Q=T\,dS\right)
    } .
    \label{eq:rt_63_23_56_thermo_route}
\end{align}
The next open lemma is the Unruh-response derivation in the hyperbolic
torsion sector.  The area law may be attacked through a line-piercing entropy
model, but its coefficient remains calibrated unless the vacuum line density is
derived independently.


\paragraph{Route II: Sakharov induced gravity.}
With ultraviolet cutoff \(\ell_{\mathrm{UV}}\sim r_c\), a schematic induced
gravity estimate gives \cite{Sakharov1967}
\begin{align}
    G_{\mathrm{ind}}
    \sim
    \frac{r_c^2 c^3}{\hbar N}.
\end{align}
Matching Newton's constant requires
\begin{align}
    N_{\mathrm{req}}
    \sim
    \frac{r_c^2 c^3}{\hbar G}
    =
    \left(\frac{r_c}{L_p}\right)^2
    =
    7.60\times10^{39}.
\end{align}

\textbf{[CRITICAL NOTE]} Unless \(N_{\mathrm{req}}\) is independently derived
from SST mode counting, this route only restates the hierarchy
\(r_c/L_p\).

\paragraph{Route III: infrared EFT universality.}
If coarse-grained SST organizes into a diffeomorphism-invariant metric EFT,
then the leading two-derivative action has the form
\begin{align}
    S_{\mathrm{IR}}
    =
    \frac{c^3}{16\pi G_{\mathrm{eff}}}
    \int d^4x \sqrt{-g}\,
    \left(R-2\Lambda\right)
    + S_{\mathrm{matter}}
    + S_{\mathrm{clock}} .
\end{align}

\textbf{[RESEARCH TRACK]} This route explains why Einstein--Hilbert appears as
the leading IR term, but it does not explain why a preferred-foliation
incompressible medium becomes diffeomorphism invariant in the infrared.

\textbf{[CRITICAL NOTE]} The cosmological-constant problem is severe
\cite{Planck2018Cosmology}:
\begin{align}
    \frac{\rho_{\!f}}{\rho_\Lambda}
    \approx
    1.2\times10^{20},
    \qquad
    \frac{\rho_{\mathrm{horn}}^{\mathrm{eff}}}{\rho_\Lambda}
    \approx
    6.6\times10^{44}.
\end{align}
The first ratio is not arbitrary: it is close to one power of the
core--Planck impedance ratio \(r_c/\ell_P\).  This observation motivates
the research-track density route below, without changing the canonical
status of \(\rho_{\!f}\) in the main text.


\subsection{Cosmological Impedance Route for the Background Density}
\label{sec:rt_rhof_cosmological_impedance_route}

\paragraph{Status.}
\textbf{[RESEARCH TRACK / NOT CANON-DERIVED].}
This subsection records a candidate route for reducing the calibrated
background density \(\rho_{\!f}\) to a cosmological vacuum input plus a
UV--IR impedance lift. It does not promote \(\rho_{\!f}\) out of the
primitive calibration set in the main canon.

The orthodox vacuum mass-equivalent density associated with a cosmological
constant is
\begin{align}
    \rho_\Lambda
    =
    \frac{\Lambda c^2}{8\pi G} .
\end{align}
For a spatially flat \(\Lambda\)CDM background,
\begin{align}
    \Lambda
    =
    \frac{3\Omega_\Lambda H_0^2}{c^2},
    \qquad
    \rho_\Lambda
    =
    \Omega_\Lambda\frac{3H_0^2}{8\pi G} .
\end{align}
Using the Planck-2018 baseline values
\(H_0\simeq 67.4\,\mathrm{km\,s^{-1}\,Mpc^{-1}}\) and
\(\Omega_\Lambda\simeq0.685\), this gives
\begin{align}
    \Lambda &\simeq 1.09\times10^{-52}\,\mathrm{m^{-2}},\\
    \rho_\Lambda &\simeq 5.85\times10^{-27}\,\mathrm{kg\,m^{-3}} .
\end{align}
The ratio required to obtain the canonical SST background density is
\begin{align}
    \frac{\rho_{\!f}}{\rho_\Lambda}
    \simeq
    1.20\times10^{20} .
\end{align}
This is close to the one-power core--Planck ratio
\begin{align}
    \frac{r_c}{\ell_P}
    =
    8.72\times10^{19},
\end{align}
which suggests the line-impedance ansatz
\begin{align}
    \rho_{\!f}
    =
    \chi_\Lambda
    \left(\frac{r_c}{\ell_P}\right)
    \rho_\Lambda,
    \qquad
    \chi_\Lambda
    \simeq
    1.37 .
    \label{eq:rt_rhof_lambda_chi}
\end{align}
Two non-equivalent simple closures are numerically relevant:
\begin{align}
    \rho_{\!f}^{(c_s)}
    &:=
    \sqrt{2}
    \left(\frac{r_c}{\ell_P}\right)
    \rho_\Lambda
    \simeq
    7.21\times10^{-7}\,\mathrm{kg\,m^{-3}},
    \\
    \rho_{\!f}^{(\Omega)}
    &:=
    2\Omega_\Lambda
    \left(\frac{r_c}{\ell_P}\right)
    \rho_\Lambda
    \simeq
    6.98\times10^{-7}\,\mathrm{kg\,m^{-3}} .
\end{align}
The first closure uses the internal core-acoustic factor already present in
\(c_s=\mathbf{v}_{\!\boldsymbol{\circlearrowleft}}/\sqrt{2}\), but does not yet
derive why the cosmological vacuum sector should couple through that
impedance. The second closure is numerically sharper for Planck-2018
parameters, but it introduces an explicitly epoch-dependent projection
\(\Omega_\Lambda\) and is therefore not acceptable as a fundamental
constant derivation unless SST supplies a time-independent interpretation of
that factor.

\paragraph{Dimensional interpretation.}
The power of \(r_c/\ell_P\) must be one, not two or three. A volume lift
\((r_c/\ell_P)^3\) or area lift \((r_c/\ell_P)^2\) overshoots catastrophically.
The only viable route is therefore a one-dimensional UV--IR mode-count or
line-impedance theorem. In physical terms, the cosmological vacuum density
would have to be lifted into a local inertial background density by counting
Planck-normal modes along one core/horn length, not across an area or volume.

\paragraph{Canonical conclusion.}
Equation~\eqref{eq:rt_rhof_lambda_chi} is a promising falsifier route, but
not a completed derivation. The main canon should continue to label
\(\rho_{\!f}\) as a calibrated effective background density. Promotion to
\textbf{[DERIVED]} requires an SST theorem that derives either
\(\chi_\Lambda=\sqrt{2}\), or an epoch-invariant replacement for
\(2\Omega_\Lambda\), from the coarse-graining map between the cosmological
vacuum sector and the local incompressible swirl medium.

\subsection{Tensor-speed naturalness}\subsection{Separatrix Surface-Density Lift via the Spherical Monopole DtN Sector}
\label{subsec:ssdl_monopole_dtn}

\paragraph{Status.}
\textbf{[RESEARCH TRACK / SPHERICAL MONOPOLE DtN NORMALIZATION PROVEN / SOURCE COUPLING OPEN]}.
This subsection records a candidate boundary-response route for the
Separatrix Surface-Density Lift (SSDL).  It proves only the geometric
\(R_\partial\)-normalization of the spherical exterior monopole
Dirichlet-to-Neumann (DtN) sector.  It does not prove the constitutive
coupling of the cosmological vacuum density \(\rho_\Lambda\) to an SST
separatrix source, and it does not replace the calibrated status of
\(\rho_{\!f}\).

\paragraph{Exterior boundary problem.}
Let \(\partial\mathcal B\) be a spherical externally resolved separatrix
of radius \(R_\partial\).  In the exterior region \(r\ge R_\partial\), let
an isotropic response potential satisfy
\begin{align}
    \nabla^2\Phi=0,
    \qquad
    \Phi(R_\partial)=\Phi_0,
    \qquad
    \Phi(\infty)=0 .
\end{align}
The unique spherical solution is
\begin{align}
    \Phi(r)=\Phi_0\frac{R_\partial}{r} .
\end{align}
With the normal convention
\begin{align}
    \Lambda_\partial[\Phi_0]
    :=
    -\partial_r\Phi\big|_{r=R_\partial},
\end{align}
the spherical monopole DtN eigenvalue is
\begin{align}
    \Lambda_{\partial,0}
    =
    \frac{1}{R_\partial},
    \qquad
    \Lambda_{\partial,0}^{-1}
    =
    R_\partial .
\end{align}
For a spherical exterior harmonic mode one has more generally
\begin{align}
    \Lambda_{\partial,\ell}
    =
    \frac{\ell+1}{R_\partial},
    \qquad
    \ell=0,1,2,\ldots .
\end{align}
Thus the \(R_\partial\) lift belongs specifically to the isotropic
monopole sector, not to arbitrary tangential boundary data.

\paragraph{Monopole projection.}
Let \(\Pi_0\) denote projection onto the spherical monopole sector,
\begin{align}
    \Pi_0 f
    =
    \frac{1}{4\pi}
    \int_{S^2} f(\theta,\varphi)\,d\Omega .
\end{align}
Then
\begin{align}
    \Pi_0\Lambda_\partial^{-1}\Pi_0[1]
    =
    R_\partial .
\end{align}
If the cosmological vacuum density enters the SST boundary-response
problem as an isotropic normal source, the proposed surface-density
projection is
\begin{align}
    \sigma_{\Lambda,\partial}
    :=
    \Pi_0\Lambda_\partial^{-1}\Pi_0[\rho_\Lambda]
    =
    \rho_\Lambda R_\partial .
\end{align}
Resolving this surface density across a Planck-normal boundary thickness
\(L_p\) gives
\begin{align}
    \rho_{\rm eff}
    =
    \frac{\sigma_{\Lambda,\partial}}{L_p}
    =
    \left(\frac{R_\partial}{L_p}\right)\rho_\Lambda .
\end{align}
Including a present-epoch vacuum projection factor gives the compact
SSDL candidate
\begin{align}
    \rho_{\!f}^{\rm SSDL}
    =
    \frac{\Omega_{\Lambda,0}}{L_p}
    \Pi_0\Lambda_\partial^{-1}\Pi_0[\rho_\Lambda] .
    \label{eq:rt_ssdl_operator_form}
\end{align}

\paragraph{Electron separatrix candidate.}
For the charged electron separatrix, take the external response radius to
be the classical electron radius
\begin{align}
    R_\partial=R_e=2r_c
    =
    \frac{e^2}{4\pi\varepsilon_0m_ec^2} .
\end{align}
The spherical-monopole reduction of Eq.~\eqref{eq:rt_ssdl_operator_form}
then becomes
\begin{align}
    \rho_{\!f}^{\rm SSDL}
    =
    \Omega_{\Lambda,0}
    \left(\frac{R_e}{L_p}\right)
    \rho_\Lambda .
    \label{eq:rt_ssdl_electron_candidate}
\end{align}
Using the audit constants
\begin{align}
    R_e&=2.8179403262\times10^{-15}\,\mathrm{m},
    &L_p&=1.6162550\times10^{-35}\,\mathrm{m},
    \\
    \rho_\Lambda&=5.8450\times10^{-27}\,\mathrm{kg\,m^{-3}},
    &\Omega_{\Lambda,0}&=0.685,
\end{align}
one obtains
\begin{align}
    \rho_{\!f}^{\rm SSDL}
    =
    6.9806682\times10^{-7}\,\mathrm{kg\,m^{-3}},
\end{align}
which differs from the calibrated canon value
\(\rho_{\!f}=7.0\times10^{-7}\,\mathrm{kg\,m^{-3}}\) by
\(-0.2762\%\).  This numerical agreement is not a derivation of
\(\rho_{\!f}\); it is a falsifiable research-track target.

\paragraph{Planck-normal mode-count equivalent.}
The discrete normal-stack version of the same route is obtained by
restricting the active response space to the isotropic sector,
\begin{align}
    \mathcal H_{\rm active}
    =
    \Pi_0\Big(L^2([0,R_\partial])\otimes L^2(S^2)\Big)
    =
    L^2([0,R_\partial]) .
\end{align}
Tangential modes are not claimed to be absent; they are projected out of
the isotropic density response by \(\Pi_0\).  If \(L_p\) is the minimal
resolved normal thickness, the active normal cell count is
\begin{align}
    N_\perp^{\rm cell}
    =
    \frac{R_\partial}{L_p}+O(1).
\end{align}
Equivalently, a one-dimensional spectral convention with half-wavelength
resolution uses
\begin{align}
    k_{\max}=\frac{\pi}{L_p},
\end{align}
so that Weyl counting gives
\begin{align}
    N_\perp^{\rm spec}
    =
    \operatorname{Tr}_\perp
    \Theta\!\left[
        \left(\frac{\pi}{L_p}\right)^2+\partial_n^2
    \right]
    =
    \frac{R_\partial}{L_p}+O(1).
\end{align}
Using \(k_{\max}=L_p^{-1}\) instead would introduce an unwanted factor
\(1/\pi\); the cutoff convention is therefore part of the research-track
claim.

\paragraph{Audit status and open lemmas.}
A v0.2 audit harness supports the operator and counting consistency of
this route: the BEM cross-check recovers \(R_e\) through
\(\Pi_0\Lambda_\partial^{-1}\Pi_0\) with relative error
\(2.68\times10^{-4}\), and the analytic Planck-normal count gives
\(N_\perp=R_e/L_p=1.743499835236\times10^{20}\) up to integer
boundary convention.  These checks are numerical and algebraic consistency
tests, not constitutive proofs.  Promotion to \textbf{[DERIVED]} requires:
\begin{enumerate}
    \item proof that \(\rho_\Lambda\) couples to the SST separatrix as an
    isotropic normal source;
    \item derivation of \(\Omega_{\Lambda,0}\), or its replacement by an
    epoch-invariant SST projection functional;
    \item derivation of \(L_p\) as the correct normal resolution thickness
    for the separatrix boundary layer.
\end{enumerate}



In the Einstein--\AE ther comparison layer,
\begin{align}
    c_T^2
    =
    \frac{c^2}{1-c_{13}},
    \qquad
    c_{13}=c_1+c_3 .
\end{align}
If SST tensor modes are identified with transverse substrate waves, and if the
canonical transverse wave speed is \(c\), then
\begin{align}
    c_T=c
    \quad\Rightarrow\quad
    c_{13}=0 .
\end{align}

\textbf{[DERIVED, conditional]} The luminal spin-2 speed is then not a tuning
but an internal consistency condition of the mode identification
\cite{FosterJacobson2006,Abbott2017GW170817}. This mirrors the main-canon
proposition of Section~\ref{subsec:tensor_speed_naturalness}.

\textbf{[LIMIT]} This does not prove the Einstein field equations. It only
fixes the spin-2 propagation speed inside the effective comparison theory.

\subsection{Final status statement}

The accurate SST statement is:

\begin{quote}
SST conditionally reconstructs relativistic kinematics for its low-energy
excitation sector, provided monometricity holds. It provides conditional routes
to the weak equivalence principle and to the Newtonian limit. It naturally
explains luminal tensor-mode speed if gravitational tensor modes are transverse
substrate waves. The full Einstein field equations remain an open research
target, with thermodynamic, induced-gravity, and infrared-EFT routes under
active investigation.
\end{quote}

\subsection{Observational falsifier ladder}

The relativity-emergence programme has three immediate observational gates:
\begin{enumerate}
    \item \textbf{Tensor-speed gate.} Multimessenger propagation requires
    \(c_T=c\), equivalently \(c_{13}=0\) in the Einstein--\AE ther comparison
    layer. Quantitatively, GW170817/GRB170817A gives
    \begin{align}
        -3\times10^{-15}
        <
        \frac{c_T}{c}-1
        <
        7\times10^{-16},
        \qquad
        |c_{13}|\lesssim10^{-15}
        \label{eq:rt_tensor_speed_observational_gate}
    \end{align}
    in the Einstein--\AE ther comparison convention \cite{LIGOVirgoFermiINTEGRAL2017,OostMukohyamaWang2018}.
    This is the spin-2 version of the transverse-mode identification.
    \item \textbf{Preferred-frame gate.} Any residual preferred-foliation
    coupling must satisfy the post-Newtonian preferred-frame constraints of
    the comparison theory. In practice this means that the effective
    \(\alpha_1\) and \(\alpha_2\) parameters cannot be freely chosen; they must
    be derived from the same substrate couplings that set the clock sector.
    Strong-field pulsar timing gives representative bounds
    \begin{align}
        \hat\alpha_1
        &=
        (-0.4^{+3.7}_{-3.1})\times10^{-5}
        \quad(95\%\ \mathrm{CL};\ \mathrm{PSR}\ J1738+0333),
        \\
        |\hat\alpha_2|
        &<
        1.6\times10^{-9}
        \quad(95\%\ \mathrm{CL};\ \mathrm{PSR}\ B1937+21,J1744-1134).
        \label{eq:rt_preferred_frame_observational_gate}
    \end{align}
    Binary-pulsar damping and PSR~J0337+1715 further reduce the post-GW170817
    viable Einstein--\AE ther parameter region by about one order of magnitude
    \cite{ShaoWex2012,ShaoCaballeroKramerWexChampionJessner2013,GuptaHerreroValeaBlasBarausseCornishYagiYunes2021}.
    \item \textbf{Species-cone gate.} Matter knots, gauge/torsion pulses, and
    clock excitations must share the same infrared characteristic cone. Failure
    of this gate predicts species-dependent Lorentz violation and directly
    falsifies exact monometricity.
\end{enumerate}

\textbf{[CRITICAL NOTE]} In the orthodox Einstein--\AE ther comparison theory,
matter is usually coupled only to the metric, not directly to the \AE ther field.
SST cannot inherit this protection automatically, because matter is modeled as a
substrate excitation. The species-cone gate is therefore stricter for SST than
for a purely metric matter-coupled Einstein--\AE ther EFT.

\textbf{[STATUS]} This is not an additional derivation. It is the observational
test hierarchy attached to the monometricity and tensor-speed theorem targets.


%======================================================================
\subsection{Research Track: Core--Torsion Impedance Matching for Inertia Closure}
\label{sec:rt_core_torsion_impedance_matching}
%======================================================================

\subsection{Purpose and status}

\textbf{[RESEARCH-TRACK].}
This section records the canon-safe form of the inertia-closure programme. It separates the internal NLSE/Gross--Pitaevskii core layer from the transverse torsion/shear radiation layer. The separation is mandatory because the canonical swirl speed \(\vchar\) is much smaller than the observed causal speed \(c\):
\begin{align}
    \frac{\vchar}{c}=3.64867628\times10^{-3}.
\end{align}
A derivation that produces \(u^2/\vchar{}^2\) or \(u^2/c_s^2\) therefore does not yet derive the Lorentz-compatible factor \(u^2/c^2\).

\subsection{NLSE/Gross--Pitaevskii core diagnostic}

For a defocusing NLSE/Gross--Pitaevskii core with Madelung phase velocity,
\begin{align}
    \mathbf u_{\rm NLS}=\frac{\hbar_*}{m_*}\nabla S,
\end{align}
the circulation quantum and healing length may be written
\begin{align}
    \Gamma_q = \frac{h_*}{m_*}N_\Gamma,
    \qquad
    \xi = \frac{\hbar_*}{\sqrt{2}\,m_*c_s}.
\end{align}
For the single-winding convention \(N_\Gamma=1\), imposing
\begin{align}
    \Gamma_q=\Gamma_0=2\pi\rc\vchar,
    \qquad
    \xi=\rc
\end{align}
gives
\begin{align}
    c_s
    =
    \frac{\Gamma_0}{2\sqrt{2}\pi\rc}
    =
    \frac{\vchar}{\sqrt{2}}
    =
    7.73465663\times10^{5}\ \mathrm{m\,s^{-1}}.
    \label{eq:rt_core_torsion_cs}
\end{align}

\textbf{[DERIVED-CONDITIONAL].}
Equation~\eqref{eq:rt_core_torsion_cs} is a useful internal diagnostic of the core layer. It is not a derivation of the Lorentz signal speed.

\subsection{Transverse torsion/shear causal layer}

A minimal transverse torsion displacement field \(\mathbf A\), with \(\nabla\cdot\mathbf A=0\), may be assigned the quadratic bridge Lagrangian
\begin{align}
    \mathcal L_{\rm torsion}
    =
    \frac12\rhoF\lVert\partial_t\mathbf A\rVert^2
    -
    \frac12K_T\lVert\nabla\times\mathbf A\rVert^2 .
\end{align}
The transverse propagation speed is
\begin{align}
    c_T^2=\frac{K_T}{\rhoF}.
\end{align}
The canon-compatible light-speed identification belongs only to this transverse layer:
\begin{align}
    c_T=c.
\end{align}
Its intensive impedance is
\begin{align}
    Z_T=\rhoF c_T=\sqrt{\rhoF K_T},
\end{align}
which has the units of acoustic impedance, \(\mathrm{kg\,m^{-2}\,s^{-1}}\).

\subsection{Core--Torsion Impedance Matching Lemma}

Let \(K\) be a localized closed swirl-string configuration with rest-swirl energy \(E_0[K]\). Define the torsion-dressed inertial tensor by the quadratic stored-energy response
\begin{align}
    \Delta E_{\rm torsion}[K;\mathbf u]
    =
    \frac12\mathbf u^{\mathsf T}
    \mathsf M_{\rm torsion}[K]
    \mathbf u
    +O(\lVert\mathbf u\rVert^4).
\end{align}
The required Lorentz-compatible closure is
\begin{align}
    \boxed{
    \mathsf M_{\rm torsion}[K]
    \stackrel{?}{=}
    \frac{2E_0[K]}{c_T^2}\,\mathsf I
    }
    \label{eq:rt_core_torsion_matching_lemma}
\end{align}
or equivalently
\begin{align}
    \Delta E_{\rm torsion}[K;u]
    \stackrel{?}{=}
    E_0[K]\frac{u^2}{c_T^2}+O(u^4).
\end{align}
If \(c_T=c\), Eq.~\eqref{eq:rt_core_torsion_matching_lemma} supplies the missing bridge from core inertia to the Lorentz-compatible clock factor.

Define the diagnostic residual
\begin{align}
    \chi_K^{(T)}
    =
    \frac{c_T^2}{2E_0[K]}
    \lambda_{\rm iso}\!\left(\mathsf M_{\rm torsion}[K]\right).
\end{align}
The bridge closes only if
\begin{align}
    \chi_K^{(T)}=1
\end{align}
within numerical and topological tolerance.

\textbf{[OPEN LEMMA].}
The equality \(\chi_K^{(T)}=1\) is the falsifiable theorem target. If numerical torsion-dressing returns a stable non-unity value, the Lorentz clock factor remains a constitutive bridge rather than a derived theorem from the core model.

\subsection{Required numerical tests}

A canon-safe numerical programme must report:
\begin{enumerate}
    \item the NLSE boost coefficient \(M_{\rm core}[K]\) and its comparison to \(2E_0[K]/c_s^2\), not to \(2E_0[K]/c^2\);
    \item the torsion-dressed tensor \(\mathsf M_{\rm torsion}[K]\) and the residual \(\chi_K^{(T)}\);
    \item a stiffness sweep \(K_T\) verifying whether \(\chi_K^{(T)}=1\) occurs naturally at \(c_T=c\) or only after retuning;
    \item the circulation-normalization ratio \(\mathcal R_\Gamma=(h_*/m_*)/\Gamma_0\), which must equal unity under the canonical single-winding convention;
    \item the trefoil-locked source test comparing \(T_{2,3}\) against generic helical sources.
\end{enumerate}

\textbf{[CANON INTERPRETATION].}
The canon may retain the two-speed discipline and the open impedance lemma. It should not promote the Lorentz clock factor to hard canon through the NLSE core alone.


% ============================================================
% Research-track patch: EM-to-swirl correspondence and stress closure
% ============================================================

\subsubsection{Research Track: EM-to-Swirl Force-Density Correspondence}
\label{sec:research_em_to_swirl_force_density_correspondence}

\paragraph{Status labels.}
The correspondence proposed in this section is \textbf{[SPECULATIVE]} and
\textbf{[PENDING DERIVATION]}.  It is not a canonical replacement of
Maxwell--Lorentz electrodynamics.  It is a Rosetta-level research bridge
between electromagnetic force density and SST substrate stress.

\paragraph{Structural correspondence.}
The electromagnetic Lorentz force density is
\[
\mathbf f_{\mathrm{EM}}
=
\rho_e\mathbf E+\mathbf J\times\mathbf B .
\]
The magnetic part has the same force-density dimension as the SST
swirl-force density,
\[
[\mathbf J\times\mathbf B]
=
\mathrm{N\,m^{-3}},
\qquad
[\rhoF\mathbf v\times\boldsymbol{\omega}]
=
\mathrm{N\,m^{-3}}.
\]
Hence the most conservative Rosetta correspondence is
\[
\boxed{
\mathbf J\times\mathbf B
\;\longleftrightarrow\;
\lambda_{\mathrm{EM}\to\circlearrowleft}\,
\rhoF\mathbf v\times\boldsymbol{\omega}.
}
\]
Since both sides already carry units of force density,
\[
\boxed{
[\lambda_{\mathrm{EM}\to\circlearrowleft}]=1 .
}
\]
The coefficient
\(\lambda_{\mathrm{EM}\to\circlearrowleft}\) is therefore dimensionless.
It is not fixed by dimensional analysis and remains
\textbf{[SPECULATIVE / PENDING DERIVATION]} until a field-level dictionary or
normalization condition is specified.

\paragraph{Relation to the canonical flux-impulse channel.}
The canonical SST electromagnetic bridge proceeds through phase, flux, and
flux-piercing variables, including the effective Faraday/flux-impulse
channel and the flux quantum
\[
\Phi_0=\frac{h}{2e}.
\]
The bulk correspondence
\[
\mathbf J\times\mathbf B
\;\longleftrightarrow\;
\lambda_{\mathrm{EM}\to\circlearrowleft}
\rhoF\mathbf v\times\boldsymbol{\omega}
\]
is therefore not an independent canonical EM bridge.  It is a candidate
continuum-limit or bulk-stress projection of the already canonical
phase/flux channel.  A canonical upgrade requires showing how the flux
observable, phase winding, and current density reduce to a definite
\(\lambda_{\mathrm{EM}\to\circlearrowleft}\) in the continuum limit.

\paragraph{Possible field-level dictionary.}
A possible but non-canonical dictionary is the classical hydrodynamic analogy
\[
\mathbf A
\;\longleftrightarrow\;
\alpha_A\mathbf v,
\qquad
\mathbf B=\nabla\times\mathbf A
\;\longleftrightarrow\;
\alpha_A\boldsymbol{\omega},
\]
and
\[
\mu_0\mathbf J
=
\nabla\times\mathbf B
\;\longleftrightarrow\;
\alpha_A\nabla\times\boldsymbol{\omega}.
\]
The constants and normalization implied by
\(\alpha_A\) are not fixed here.  This dictionary is retained only as a
research-track route toward deriving
\(\lambda_{\mathrm{EM}\to\circlearrowleft}\).

\paragraph{Stress closure and SST-44 tensor.}
The canonical swirl-stress candidate is the symmetric Cauchy-type stress
used in the SST-44 appendix,
\[
\sigma^{(44)}_{ij}
=
\rhoF v_i v_j
+
\rhoF \rc^2
\left(
\omega_i\omega_j-\frac12\delta_{ij}|\boldsymbol{\omega}|^2
\right)
-
\delta_{ij}\frac12\rhoF|\mathbf v|^2 .
\]
Because this tensor already contains the isotropic dynamic-pressure term
\[
-\delta_{ij}\frac12\rhoF|\mathbf v|^2,
\]
one must not also add an independent
\(-\nabla p_{\circlearrowleft}\) term with
\[
p_{\circlearrowleft}
=
\frac12\rhoF|\mathbf v|^2 .
\]
Otherwise the Bernoulli contribution is counted twice.

\paragraph{Two admissible decompositions.}
The full-stress convention is
\[
\boxed{
\mathbf f_{\mathrm{SST}}
=
\nabla\cdot\boldsymbol{\sigma}^{(44)}
+
\mathbf f_{\mathrm{link}} .
}
\]
The split-stress convention is
\[
\boxed{
\mathbf f_{\mathrm{SST}}
=
-\nabla p_{\circlearrowleft}
+
\nabla\cdot\boldsymbol{\sigma}^{\mathrm{dev}}_{\circlearrowleft}
+
\mathbf f_{\mathrm{link}},
}
\]
where
\[
\boldsymbol{\sigma}^{\mathrm{dev}}_{\circlearrowleft}
\equiv
\rhoF\mathbf v\mathbf v
+
\rhoF\rc^2
\left(
\boldsymbol{\omega}\boldsymbol{\omega}
-
\frac12\mathbf I|\boldsymbol{\omega}|^2
\right)
\]
contains the advective and torsional anisotropic parts but excludes the
isotropic dynamic-pressure term.  These two conventions must not be mixed.

\paragraph{Topological link force.}
The link-force term is not yet canonical.  The intended structure is a
topological potential functional
\[
U_{\mathrm{link}}
=
U_{\mathrm{link}}\!\left[
Lk(C_a,C_b),\Gamma_a,\Gamma_b,H
\right],
\]
with force density obtained schematically by variation,
\[
\boxed{
\mathbf f_{\mathrm{link}}
=
-\frac{\delta U_{\mathrm{link}}}{\delta \mathbf x}
}
\]
or, for a collective coordinate \(q\),
\[
\boxed{
F_{\mathrm{link},q}
=
-\frac{\partial U_{\mathrm{link}}}{\partial q}.
}
\]
The linking number is
\[
Lk(C_a,C_b)
=
\frac{1}{4\pi}
\oint_{C_a}\oint_{C_b}
\frac{
(d\mathbf x_a\times d\mathbf x_b)\cdot(\mathbf x_a-\mathbf x_b)
}{
|\mathbf x_a-\mathbf x_b|^3
}.
\]
Until \(U_{\mathrm{link}}\) is specified and normalized,
\(\mathbf f_{\mathrm{link}}\) remains
\textbf{[PENDING DERIVATION]}.

\paragraph{Research-track conclusion.}
The candidate bulk EM-to-swirl mapping is therefore
\[
\boxed{
\mathbf f_{\mathrm{EM}}
=
\rho_e\mathbf E+\mathbf J\times\mathbf B
\quad\leadsto\quad
\mathbf f_{\mathrm{SST}}
=
\nabla\cdot\boldsymbol{\sigma}^{(44)}
+
\mathbf f_{\mathrm{link}}
}
\]
under the full-stress convention, or equivalently
\[
\boxed{
\mathbf f_{\mathrm{SST}}
=
-\nabla p_{\circlearrowleft}
+
\nabla\cdot\boldsymbol{\sigma}^{\mathrm{dev}}_{\circlearrowleft}
+
\mathbf f_{\mathrm{link}}
}
\]
under the split-stress convention.  The correspondence becomes canonizable
only after
\(\lambda_{\mathrm{EM}\to\circlearrowleft}\),
the field-level dictionary, and
\(U_{\mathrm{link}}\)
are derived or experimentally fixed.



% ============================================================
% Research-track patch: Taylor-column finite-thickness and parity-delay benchmark
% ============================================================
\subsection{Research Track: Taylor-Column Analogues for Finite-Thickness Swirl Strings}
\label{sec:taylor_column_finite_thickness_swirl_strings}

\providecommand{\rhoF}{\rho_{\!f}}
\providecommand{\vSwirl}{\mathbf{v}_{\!\boldsymbol{\circlearrowleft}}}
\providecommand{\rc}{r_c}
\providecommand{\GammaK}{\Gamma_K}
\providecommand{\GammaZero}{\Gamma_0}
\providecommand{\Om}{\Omega}
\providecommand{\Ht}{\mathcal{H}_{\times}}

\paragraph{Status.}
This subsection is \textbf{Research Track}.  The Taylor--Proudman,
Hill-vortex, relative-helicity, inertial-wave, and Feynman--Onsager statements
are orthodox rotating-fluid or superfluid mechanics
\cite{Proudman1916,Taylor1922,Batchelor1967,Greenspan1968,Hill1894,Onsager1949,Feynman1955}.
Their SST use is restricted to analogue benchmarking, scale-separation
diagnostics, and falsification tests.  No macroscopic rotating-tank parameter is
identified with the SST canonical core scale.

\subsubsection{Background kinematics and resolved scale separation}

Consider a rotating cylindrical background of height
\(H=1\,\mathrm m\), radius \(R=0.25\,\mathrm m\), and constant angular velocity
\[
\boldsymbol{\Omega}=\Omega\,\hat{\mathbf z},
\qquad
\Omega=1\,\mathrm{s^{-1}}.
\]
The unperturbed velocity and vorticity fields are
\[
\mathbf v_b=\Omega r\,\hat{\boldsymbol\theta},
\qquad
\boldsymbol\omega_b=\nabla\times\mathbf v_b
=2\Omega\,\hat{\mathbf z}.
\]
The background has zero local kinetic helicity density:
\[
\mathbf v_b\cdot\boldsymbol\omega_b
=(\Omega r\,\hat{\boldsymbol\theta})\cdot(2\Omega\,\hat{\mathbf z})=0.
\]

Let \(a\) denote the resolved separatrix radius of a macroscopic vortex
structure, such as a vortex ring, Hill-type vortex, or trefoil vortex tube.
Let \(r_c\) denote the SST canonical core radius.  The resolved-scale guardrail is
\[
\boxed{a\neq r_c.}
\]
Using
\[
r_c=1.40897017\times10^{-15}\,\mathrm m,
\qquad
\vSwirl=1.09384563\times10^6\,\mathrm{m\,s^{-1}},
\]
the canonical core-scale circulation is
\[
\Gamma_0=2\pi r_c\vSwirl
=9.68361920\times10^{-9}\,\mathrm{m^2\,s^{-1}}.
\]

For a macroscopic separatrix radius
\[
a=\frac{R}{4}=0.0625\,\mathrm m,
\]
the background circulation around the same radius is
\[
\Gamma_b(a)=\oint_{r=a}\mathbf v_b\cdot d\boldsymbol\ell
=2\pi\Omega a^2.
\]
This defines the dimensionless synchronization parameter
\[
\boxed{
\chi_\Omega=\frac{\Gamma_K}{2\pi\Omega a^2}.
}
\]
For \(\Gamma_K=\Gamma_0\) and \(\Omega=1\,\mathrm{s^{-1}}\),
\[
2\pi\Omega a^2=2.45436926\times10^{-2}\,\mathrm{m^2\,s^{-1}},
\]
so that
\[
\boxed{
\chi_\Omega=3.94546141\times10^{-7}.
}
\]
Thus a macroscopic separatrix equipped only with the canonical core circulation
is kinematically far below synchronization with a \(1\,\mathrm{s^{-1}}\)
tank-scale rotation.

\paragraph{Dimensional check.}
\[
[\Gamma_K]=\mathrm{m^2\,s^{-1}},
\qquad
[\Omega a^2]=\mathrm{s^{-1}m^2}=\mathrm{m^2\,s^{-1}},
\]
so \(\chi_\Omega\) is dimensionless.

\subsubsection{Separatrix response and relative helicity}

For a slowly translating finite-thickness structure with axial speed \(U\), the
Rossby number is
\[
Ro=\frac{U}{2\Omega a}.
\]
In the low-Rossby regime \(Ro\ll1\), the leading-order Taylor--Proudman
constraint gives
\[
\frac{\partial\mathbf u}{\partial z}\approx0.
\]
Hence, to leading order, a no-through-flow separatrix behaves externally as a
columnar obstruction.  This does not imply that the internal topology is
sphere-like; it only states that the exterior flow sees a columnar separatrix
constraint.

For an incompressible closed cross-section, the axial flux balance is
\[
\boxed{
\int_0^R2\pi r\,u_z(r)\,dr=0.
}
\]
A simple two-zone model is
\[
u_z(r)=
\begin{cases}
U, & 0\le r<a,\\[1mm]
-U\dfrac{a^2}{R^2-a^2}, & a<r\le R.
\end{cases}
\]
Indeed,
\[
\int_0^a2\pi rU\,dr
+
\int_a^R2\pi r\left(-U\frac{a^2}{R^2-a^2}\right)dr
=\pi a^2U-\pi a^2U=0.
\]

For a vortex ring with circulation \(\Gamma_K\) and ring surface \(S_K\), a
useful diagnostic is the relative or mutual helicity with the rotating
background \cite{Moffatt1969}:
\[
\boxed{
\mathcal H_{\times}
=2\Gamma_K\int_{S_K}\boldsymbol\omega_b\cdot d\mathbf S.
}
\]
Because
\[
\boldsymbol\omega_b=2\Omega\hat{\mathbf z},
\qquad
\int_{S_K}\hat{\mathbf z}\cdot d\mathbf S=\pi a^2,
\]
one obtains
\[
\boxed{
\mathcal H_{\times}=4\pi\Omega a^2\Gamma_K.
}
\]
For \(\Gamma_K=\Gamma_0\),
\[
\boxed{
\mathcal H_{\times,0}
=4.75343546\times10^{-10}\,\mathrm{m^4\,s^{-2}}.
}
\]

\paragraph{Dimensional check.}
\[
[\mathcal H_{\times}]
=[\Omega a^2\Gamma_K]
=\mathrm{s^{-1}m^2\,m^2s^{-1}}
=\mathrm{m^4\,s^{-2}}
=[\Gamma_K]^2.
\]

\subsubsection{Hill-vortex sphere analogue}

A vortex-ring dipole with a closed separatrix may be approximated, in the
continuous Euler limit, by Hill's spherical vortex \cite{Hill1894}.  In the
frame where the undisturbed fluid at infinity is at rest, the kinetic energy
splits into an external irrotational contribution and an internal rotational
contribution:
\[
E_{\rm out}=\frac{1}{3}\pi\rho_{\!f}a^3U^2,
\qquad
E_{\rm in}=\frac{23}{21}\pi\rho_{\!f}a^3U^2.
\]
Therefore
\[
\boxed{
E_{\rm tot}=E_{\rm out}+E_{\rm in}
=\frac{10}{7}\pi\rho_{\!f}a^3U^2.
}
\]
This result is used only as a continuum benchmark for finite-thickness
separatrix energetics.  It does not identify a trefoil vortex with a
featureless sphere; the equivalence is external and leading-order only.

\paragraph{Numerical reference for \(U=1\,\mathrm{m\,s^{-1}}\).}
With
\[
\rho_{\!f}=7.0\times10^{-7}\,\mathrm{kg\,m^{-3}},
\qquad
a=0.0625\,\mathrm m,
\]
one obtains
\[
E_{\rm out}=1.78964425\times10^{-10}\,\mathrm J,
\]
\[
E_{\rm in}=5.88025969\times10^{-10}\,\mathrm J,
\]
and
\[
E_{\rm tot}=7.66990394\times10^{-10}\,\mathrm J.
\]

\subsubsection{Axial force density from vorticity-pressure gradients}

If the local background rotation varies along the axis, \(\Omega=\Omega(z)\),
ordinary Euler/Bernoulli balance gives
\[
\frac{1}{\rho_{\!f}}\frac{\partial p}{\partial r}
=\Omega(z)^2r.
\]
Integrating from the axis to the wall yields
\[
p(0,z)=p(R,z)-\frac{1}{2}\rho_{\!f}\Omega(z)^2R^2.
\]
If \(p(R,z)\) is fixed or externally controlled, the axial force density on the
central region is
\[
\boxed{
f_z=-\frac{\partial p(0,z)}{\partial z}
=\frac{1}{2}\rho_{\!f}R^2
\frac{\partial}{\partial z}\left[\Omega(z)^2\right].
}
\]
This force is part of the ordinary vorticity-pressure impulse budget:
\[
\boxed{
\mathbf F_z\subset\mathbf F_{\omega}.
}
\]
It must be analytically subtracted before any residual SST-specific force claim
is considered.

\paragraph{Dimensional check.}
\[
[\rho_{\!f}R^2\partial_z\Omega^2]
=\mathrm{kg\,m^{-3}}\,\mathrm{m^2}\,\mathrm{m^{-1}s^{-2}}
=\mathrm{kg\,m^{-2}s^{-2}}
=\mathrm{N\,m^{-3}}.
\]
Thus \(f_z\) has the unit of force density.

\subsubsection{Trefoil dynamics and He-II falsification}

A macroscopic trefoil vortex is not a sphere internally.  Its internal geometry
carries writhe, twist, finite-thickness contact structure, and self-induced
Biot--Savart dynamics.  In water-like classical fluids, hydrofoil-generated
trefoil vortices are expected to deform strongly, redistribute helicity, and
undergo reconnection \cite{KlecknerIrvine2013,KoplikLevine1993}.

In orthodox rotating superfluid helium, the background rotation is carried by a
Feynman--Onsager lattice of quantized vortex lines \cite{Onsager1949,Feynman1955}.
For \({}^4\mathrm{He}\),
\[
\kappa_4=\frac{h}{m_4}=9.9693\times10^{-8}\,\mathrm{m^2\,s^{-1}}.
\]
The vortex-line areal density is
\[
n_v=\frac{2\Omega}{\kappa_4}.
\]
For \(\Omega=1\,\mathrm{s^{-1}}\),
\[
\boxed{
n_v=2.0062\times10^7\,\mathrm{m^{-2}}.
}
\]
The number of quantized vortex lines crossing the resolved separatrix area
\(\pi a^2\) is
\[
N_v=n_v\pi a^2.
\]
For \(a=0.0625\,\mathrm m\),
\[
\boxed{
N_v=2.4619\times10^5.
}
\]

The orthodox He-II expectation is therefore a high reconnection rate,
Kelvin-wave cascade, partial helicity redistribution, and likely unknotting or
decay of the trefoil.  By contrast, the ideal SST continuum limit assumes
inviscid, incompressible flow with no reconnecting cuts, so topology is preserved
unless an explicit reconnection or defect-bridge sector is added.  The
disagreement defines a falsification domain:
\[
\boxed{\text{He-II quantized vortex lattice: reconnection-enabled decay,}}
\]
\[
\boxed{\text{ideal SST continuum: topology-preserving transport.}}
\]

\subsection{Research Track: Rotating-Fluid Delay and Swirl-Clock Parity}
\label{sec:rotating_fluid_swirl_clock_parity}

\paragraph{Status.}
This subsection is \textbf{Research Track}.  It uses orthodox inertial-wave and
rotating-vorticity relations as analogue benchmarks and then records the
SST-compatible parity split between odd tracer observables and even clock-rate
observables.  No dimensionally invalid fine-structure normalization is used.

\subsubsection{Arrival delay of the surface signal: inertial-wave delay}

For a small perturbation in a uniformly rotating inviscid incompressible fluid,
the inertial-wave dispersion relation is
\[
\omega=2\Omega\frac{k_z}{k},
\qquad
k^2=k_r^2+k_z^2.
\]
The vertical group velocity is
\[
c_{g,z}=\frac{\partial\omega}{\partial k_z}
=2\Omega\frac{k_r^2}{k^3}.
\]
Therefore the arrival time of a disturbance over height \(H\) is
\[
\boxed{
t_{\mathrm{arr}}
=\frac{H}{c_{g,z}}
=\frac{Hk^3}{2\Omega k_r^2}.
}
\]
Thus
\[
\boxed{t_{\mathrm{arr}}\propto\Omega^{-1}.}
\]
This is a propagation delay of the rotating-fluid disturbance.  It is not a
swirl-clock time-dilation effect.

\paragraph{Dimensional check.}
\[
\left[\frac{Hk^3}{2\Omega k_r^2}\right]
=\frac{\mathrm{m}\,\mathrm{m^{-3}}}{\mathrm{s^{-1}m^{-2}}}
=\mathrm{s}.
\]

\subsubsection{Clock delay from local swirl speed: even-in-direction effect}

In the SST swirl-clock interpretation, the local proper-time rate associated
with a resolved azimuthal swirl speed \(u_\theta\) is taken in the dimensionally
consistent form
\[
\boxed{
\frac{d\tau}{dt}
=\sqrt{1-\frac{u_\theta^2}{c^2}}.
}
\]
For \(|u_\theta|\ll c\),
\[
\frac{d\tau}{dt}
=1-\frac{1}{2}\frac{u_\theta^2}{c^2}
+\mathcal O\!\left(\frac{u_\theta^4}{c^4}\right).
\]
The accumulated proper-time deficit over a laboratory interval \(T\) is
therefore
\[
\Delta\tau-T
=\int_0^T\left(\sqrt{1-\frac{u_\theta(t)^2}{c^2}}-1\right)dt
\approx
-\frac{1}{2c^2}\int_0^Tu_\theta(t)^2\,dt.
\]
Because the dependence is quadratic,
\[u_\theta\rightarrow-u_\theta
\quad\Longrightarrow\quad
u_\theta^2\rightarrow u_\theta^2,
\]
so
\[
\boxed{
\Delta\tau(+u_\theta)=\Delta\tau(-u_\theta)<T.
}
\]
Both swirl directions therefore produce clock slowing.  Directional reversal
changes the sign of tracer rotation, but not the sign of the clock-rate deficit.

\subsubsection{Vorticity response to vertical displacement}

In the low-Rossby rotating-fluid limit, the linearized vorticity equation gives
\[
\partial_t\boldsymbol\omega
\approx
2(\boldsymbol\Omega\cdot\nabla)\mathbf u.
\]
For an axisymmetric vertical displacement field
\[
\mathbf u=w(r,z,t)\,\hat{\mathbf z},
\qquad
w=\partial_t\xi,
\]
the vertical vorticity response is
\[
\partial_t\omega_z
=2\Omega\,\partial_z w
=2\Omega\,\partial_z\partial_t\xi.
\]
Integrating in time with zero initial perturbation vorticity gives
\[
\boxed{
\omega_z(r,z,t)=2\Omega\,\partial_z\xi(r,z,t).
}
\]
Thus column stretching and compression generate opposite-signed vertical
vorticity.

For a Gaussian displacement profile,
\[
\xi(r,z,t)
=Z(t)\exp\left[-\frac{r^2+z^2}{a^2}\right],
\]
one obtains
\[
\omega_z(r,z,t)
=-\frac{4\Omega Z(t)z}{a^2}
\exp\left[-\frac{r^2+z^2}{a^2}\right].
\]
Hence the induced vorticity is odd in \(z\):
\[
\boxed{
\omega_z(r,-z,t)=-\omega_z(r,z,t).
}
\]

\subsubsection{Resolved swirl velocity from induced vorticity}

For an axisymmetric azimuthal velocity \(u_\theta(r,z,t)\),
\[
\omega_z=\frac{1}{r}\frac{\partial}{\partial r}\left(r u_\theta\right).
\]
Using the Gaussian vorticity source above gives
\[
\boxed{
u_\theta(r,z,t)
=-\frac{2\Omega Z(t)z}{a^2}
e^{-z^2/a^2}
\frac{1-e^{-r^2/a^2}}{r}.
}
\]
Therefore
\[u_\theta(r,-z,t)=-u_\theta(r,z,t),
\qquad
u_\theta(r,-z,t)^2=u_\theta(r,z,t)^2.
\]
This establishes the parity split:
\[
\boxed{
\text{tracer angle is odd in swirl direction,}
\qquad
\text{clock deficit is even in swirl direction.}
}
\]

\subsubsection{C1--C4 conclusions}

\paragraph{C1. Inertial-wave benchmark.}
The delayed push--pull surface response in a rotating tank is explained, to
leading order, as an inertial-wave arrival delay:
\[
\boxed{
t_{\mathrm{arr}}=\frac{Hk^3}{2\Omega k_r^2}.
}
\]
It scales as \(t_{\mathrm{arr}}\propto\Omega^{-1}\), and therefore disappears
as a rotating-fluid effect when the background rotation is removed.

\paragraph{C2. Vorticity generation by vertical forcing.}
A vertical displacement in a rotating incompressible fluid produces vertical
vorticity according to
\[
\boxed{
\omega_z=2\Omega\,\partial_z\xi.
}
\]
Stretching and compression of the rotating column generate opposite signs of
relative vorticity above and below the symmetry plane.

\paragraph{C3. Odd tracer response and cancellation over symmetric cycles.}
The induced azimuthal velocity satisfies
\[u_\theta(r,-z,t)=-u_\theta(r,z,t).
\]
Hence tracer angles reverse under reflection or reversed forcing.  For an ideal
symmetric up--down cycle, the leading-order angular displacement cancels:
\[
\boxed{
\Delta\theta_{\mathrm{cycle}}\simeq0
\qquad
\text{to leading order in }Ro.
}
\]

\paragraph{C4. Even clock-rate response and accumulated delay.}
The swirl-clock deficit depends on \(u_\theta^2\), not on \(u_\theta\).  Therefore
\[
\boxed{
\Delta\tau(+u_\theta)=\Delta\tau(-u_\theta)<T.
}
\]
Equivalently,
\[
\boxed{
\Delta\tau-T
\approx
-\frac{1}{2c^2}\int_0^Tu_\theta(t)^2\,dt.
}
\]
Thus a symmetric cycle can have zero net tracer rotation while still producing a
nonzero accumulated clock deficit.

\paragraph{Core-scale guardrail.}
The laboratory rotation rate \(\Omega\) must not be identified with the SST
canonical core turnover rate.  The latter is
\[
\Omega_0=\frac{\vSwirl}{r_c}.
\]
Using the canonical constants gives
\[
\boxed{
\Omega_0=7.76344066\times10^{20}\,\mathrm{s^{-1}}.
}
\]
Therefore
\[
\boxed{
\Omega_{\mathrm{lab}}\neq\Omega_0.
}
\]
The rotating-tank experiment is a resolved analogue benchmark, not a direct
identification of macroscopic rotation with the canonical SST core scale.

\paragraph{Summary.}
The experiment separates two parities:
\[
\boxed{
\theta(+u_\theta)=-\theta(-u_\theta),
}
\]
but
\[
\boxed{
\Delta\tau(+u_\theta)=\Delta\tau(-u_\theta).
}
\]
This provides a falsifiable benchmark: angular tracer motion is direction
sensitive, whereas the clock-rate or energy-density deficit is direction
independent at leading quadratic order.

\paragraph{Child-level analogy.}
A knotted vortex in a rotating fluid is like a small spinning propeller inside a
slowly rotating elevator shaft.  From outside, it may push fluid like a round
object moving up and down, but inside the object the strings are tied into a
knot.  Turning the knot the other way reverses the visible swirl, but the amount
of busy motion stored in the fluid still slows the clock in both directions.


% ----------------------------------------------------------
% Bibliography entries for the Relativity Emergence Ladder
% ----------------------------------------------------------
\ifresearchtrackincluded\else
\begin{thebibliography}{99}


\bibitem{Taylor1922}
G.~I. Taylor,
`The motion of a sphere in a rotating liquid,''
\emph{Proceedings of the Royal Society of London A} \textbf{102} (1922), 180--189.
DOI: \href{https://doi.org/10.1098/rspa.1922.0079}{10.1098/rspa.1922.0079}.

\bibitem{Proudman1916}
J.~Proudman,
`On the motion of solids in a liquid possessing vorticity,''
\emph{Proceedings of the Royal Society of London A} \textbf{92} (1916), 408--424.
DOI: \href{https://doi.org/10.1098/rspa.1916.0026}{10.1098/rspa.1916.0026}.

\bibitem{Hill1894}
M.~J.~M. Hill,
`On a spherical vortex,''
\emph{Philosophical Transactions of the Royal Society of London A} \textbf{185} (1894), 213--245.
DOI: \href{https://doi.org/10.1098/rsta.1894.0006}{10.1098/rsta.1894.0006}.

\bibitem{Greenspan1968}
H.~P. Greenspan,
\emph{The Theory of Rotating Fluids},
Cambridge University Press (1968).
ISBN: 978-0521051477.

\bibitem{LIGOVirgoFermiINTEGRAL2017}
B.~P. Abbott et al. (LIGO Scientific Collaboration, Virgo Collaboration, Fermi Gamma-ray Burst Monitor, and INTEGRAL),
``Gravitational Waves and Gamma-rays from a Binary Neutron Star Merger: GW170817 and GRB 170817A,''
\emph{Astrophysical Journal Letters} \textbf{848}, L13 (2017).
DOI: \href{https://doi.org/10.3847/2041-8213/aa920c}{10.3847/2041-8213/aa920c}.

\bibitem{OostMukohyamaWang2018}
J.~Oost, S.~Mukohyama, and A.~Wang,
``Constraints on Einstein-aether theory after GW170817,''
\emph{Physical Review D} \textbf{97}, 124023 (2018).
DOI: \href{https://doi.org/10.1103/PhysRevD.97.124023}{10.1103/PhysRevD.97.124023}.

\bibitem{ShaoWex2012}
L.~Shao and N.~Wex,
``New tests of local Lorentz invariance of gravity with small-eccentricity binary pulsars,''
\emph{Classical and Quantum Gravity} \textbf{29}, 215018 (2012).
DOI: \href{https://doi.org/10.1088/0264-9381/29/21/215018}{10.1088/0264-9381/29/21/215018}.

\bibitem{ShaoCaballeroKramerWexChampionJessner2013}
L.~Shao, R.~N. Caballero, M.~Kramer, N.~Wex, D.~J. Champion, and A.~Jessner,
``A new limit on local Lorentz invariance violation of gravity from solitary pulsars,''
\emph{Classical and Quantum Gravity} \textbf{30}, 165019 (2013).
DOI: \href{https://doi.org/10.1088/0264-9381/30/16/165019}{10.1088/0264-9381/30/16/165019}.

\bibitem{GuptaHerreroValeaBlasBarausseCornishYagiYunes2021}
T.~Gupta, M.~Herrero-Valea, D.~Blas, E.~Barausse, N.~Cornish, K.~Yagi, and N.~Yunes,
``New binary pulsar constraints on Einstein-\ae ther theory after GW170817,''
\emph{Classical and Quantum Gravity} \textbf{38}, 195003 (2021).
DOI: \href{https://doi.org/10.1088/1361-6382/ac1a69}{10.1088/1361-6382/ac1a69}.

\end{thebibliography}
\fi
